% MAP WORK — Geology Upper Sixth — Cours signé OnBuch+
% Official MINESEC syllabus. Compiler avec : tectonic GEO03-map-work.tex
\newcommand{\DOCMATIERE}{Geology}
\newcommand{\DOCNIVEAU}{Upper Sixth}
\newcommand{\DOCMODULE}{Upper Sixth — Geology}
\newcommand{\DOCLECON}{3}
\newcommand{\DOCTITRE}{MAP WORK}
% OnBuch+ — préambule « cours signés » ENRICHI (compilé avec tectonic / XeLaTeX)
% Système visuel « Pop » d'OnBuch+ : fond crème (jamais blanc), bordures encre
% épaisses, ombres « dures » décalées, titres Archivo Black en capitales.
% Version MATHÉMATIQUES : graphes (pgfplots/tikz), boîtes pédagogiques
% (définition, propriété, méthode, attention, le savais-tu, exercice résolu,
% exercice, TP, bilan), page de garde et signature OnBuch+ ancrée en bas.
%
% Macros attendues dans main.tex AVANT \input{preamble} :
%   \DOCMATIERE  \DOCNIVEAU  \DOCMODULE  \DOCLECON (numéro)  \DOCTITRE
\documentclass[11pt]{article}

\usepackage{fontspec}
\usepackage[english]{babel}
\usepackage[a4paper,margin=2cm,top=3.4cm,bottom=2.8cm,headheight=1.4cm,footskip=1.3cm]{geometry}
\usepackage{amsmath,amssymb,amsfonts}
\usepackage{mathtools} % \xmapsto, \coloneqq... souvent utilisés par les modèles
\usepackage{cancel} % simplifications barrées (bilans de forces)
% \abs{z} (module, valeur absolue) et \norm{u} (norme), utilisés par les modèles
\providecommand{\abs}{}\let\abs\relax\DeclarePairedDelimiter{\abs}{\lvert}{\rvert}
\providecommand{\norm}{}\let\norm\relax\DeclarePairedDelimiter{\norm}{\lVert}{\rVert}
\usepackage[table]{xcolor} % [table] : fournit aussi \rowcolors (alternance de lignes)
\usepackage{pagecolor}
\usepackage{tikz}
\usetikzlibrary{arrows.meta,positioning,calc,shapes.geometric,shapes.misc,shapes.arrows,decorations.pathmorphing,decorations.pathreplacing,decorations.markings,patterns,shadows,mindmap,trees,fit,backgrounds,matrix,intersections,angles,quotes}
\usepackage{pgfplots}
\usepackage[version=4]{mhchem} % équations nucléaires \ce{^{235}_{92}U}
\usepackage[european,siunitx]{circuitikz} % schémas électriques
\pgfplotsset{compat=1.18}
\usepgfplotslibrary{fillbetween,groupplots} % aires entre courbes (calcul intégral)
\usepackage{enumitem}
\usepackage{fancyhdr}
% [most] charge toutes les bibliothèques tcolorbox (listings, minted, raster,
% documentation...) et gonfle inutilement la mémoire TeX sur un document long
% (~300 boîtes) : on ne charge que ce qui est réellement utilisé.
\usepackage[skins,breakable]{tcolorbox}
\usepackage{graphicx}
\usepackage{colortbl}
\usepackage{booktabs}
\usepackage{tabularx}
\usepackage{diagbox} % case d'en-tête diagonale des tableaux à double entrée
% Colonnes extensibles centrée / gauche / droite (C, L, R), souvent utilisées par les modèles
\newcolumntype{C}{>{\centering\arraybackslash}X}
\newcolumntype{L}{>{\raggedright\arraybackslash}X}
\newcolumntype{R}{>{\raggedleft\arraybackslash}X}
\usepackage{array}
\usepackage{multicol}
\usepackage{siunitx}
% parse-numbers=false : accepte aussi \SI{2,5 \times 10^{-3}}{...}
\sisetup{locale=UK,output-decimal-marker={.},group-minimum-digits=4,parse-numbers=false}
\DeclareSIUnit\ohm{\ensuremath{\symup{\Omega}}} % U+2126 absent de la police
\DeclareSIUnit\division{div} % réglages d'oscilloscope (ms/div, V/div)
\DeclareSIUnit\tr{tr} % tours (vitesse de rotation, ex. \SI{1500}{\tr\per\minute})
\usepackage{qrcode}
% \degree : commande fréquemment utilisée par les modèles pour un angle ou une
% température en dehors de \SI{}{} (non fournie par les paquets chargés).
\providecommand{\degree}{^{\circ}}
% \qed (fin de démonstration), utilisé par les modèles sans amsthm
\providecommand{\qed}{\hfill$\square$}
% \barre : double barre des valeurs interdites dans un tableau de variations (tkz-tab)
\providecommand{\barre}{\ensuremath{\|}}
% environnement proof (amsthm non chargé) utilisé par les modèles
\ifdefined\proof\else\newenvironment{proof}[1][Proof]{\par\noindent\textit{#1.}\ }{\qed\par}\fi
\usepackage{lastpage}
\usepackage{hyperref}
\hypersetup{hidelinks}

% ============================================================
% Palette « Pop » OnBuch+ — mêmes hex que la classe P (plus_ui.dart)
% ============================================================
\definecolor{poporange}{HTML}{F59321}
\definecolor{popdark}{HTML}{E07A0C}
\definecolor{popcream}{HTML}{FFF6E8}
\definecolor{popink}{HTML}{1C1714}
\definecolor{poppurple}{HTML}{7A5AE0}
\definecolor{popgreen}{HTML}{2E9E6A}
\definecolor{poppink}{HTML}{E0457B}
\definecolor{popblue}{HTML}{2F7DE1}
\definecolor{popgold}{HTML}{C98A12}
\definecolor{popmuted}{HTML}{8A7F76}
\definecolor{popline}{HTML}{EADFD2}
\definecolor{popgray}{HTML}{8A7F76}
\colorlet{popgrey}{popgray}
% Alias tolérants (le modèle nomme parfois les couleurs différemment)
\colorlet{popwhite}{white}
\colorlet{popblack}{popink}
\colorlet{popred}{poppink}
\colorlet{popyellow}{popgold}
% Teintes claires pour fonds de tableaux / zones
\colorlet{popblueL}{popblue!10!white}
\colorlet{poporangeL}{poporange!14!white}
\colorlet{popgreenL}{popgreen!12!white}
\colorlet{poppurpleL}{poppurple!12!white}
\colorlet{poppinkL}{poppink!12!white}
\colorlet{popgoldL}{popgold!14!white}

\pagecolor{popcream}

% ============================================================
% Typographie
% ============================================================
\defaultfontfeatures{Ligatures=TeX}
\setmainfont{PlusJakartaSans-Medium.ttf}[Path=../../../fonts/,BoldFont=PlusJakartaSans-Bold.ttf]
% Maths en sans-serif (Fira Math) avec chiffres et lettres droites de Plus
% Jakarta, pour que les formules se fondent dans le texte.
\usepackage[math-style=ISO,bold-style=ISO]{unicode-math}
\setmathfont{FiraMath-Regular.otf}[Path=../../../fonts/]
\setmathfont{PlusJakartaSans-Medium.ttf}[Path=../../../fonts/,range={up/num,up/{Latin,latin},\mathup}]
% (icomma retiré : décimales avec point en anglais)
% Caractères mathématiques Unicode absents des polices de texte (Plus Jakarta,
% Archivo Black) : rendus par leur équivalent mathématique, en texte comme en
% formule — sinon ils s'affichent en carré vide (ex. « Arithmétique dans ℤ »).
\usepackage{newunicodechar}
\newunicodechar{ℝ}{\ensuremath{\mathbb{R}}}
\newunicodechar{ℤ}{\ensuremath{\mathbb{Z}}}
\newunicodechar{ℂ}{\ensuremath{\mathbb{C}}}
\newunicodechar{ℕ}{\ensuremath{\mathbb{N}}}
\newunicodechar{ℚ}{\ensuremath{\mathbb{Q}}}
\newunicodechar{𝔻}{\ensuremath{\mathbb{D}}}
\newunicodechar{α}{\ensuremath{\alpha}}
\newunicodechar{β}{\ensuremath{\beta}}
\newunicodechar{γ}{\ensuremath{\gamma}}
\newunicodechar{δ}{\ensuremath{\delta}}
\newunicodechar{ε}{\ensuremath{\varepsilon}}
\newunicodechar{θ}{\ensuremath{\theta}}
\newunicodechar{λ}{\ensuremath{\lambda}}
\newunicodechar{μ}{\ensuremath{\mu}}
\newunicodechar{σ}{\ensuremath{\sigma}}
\newunicodechar{τ}{\ensuremath{\tau}}
\newunicodechar{φ}{\ensuremath{\varphi}}
\newunicodechar{ψ}{\ensuremath{\psi}}
\newunicodechar{ω}{\ensuremath{\omega}}
\newunicodechar{Σ}{\ensuremath{\Sigma}}
% majuscules produites par \MakeUppercase dans les titres (α-aminés -> Α)
\newunicodechar{Α}{\ensuremath{\alpha}}
\newunicodechar{Β}{\ensuremath{\beta}}
\newunicodechar{Γ}{\ensuremath{\Gamma}}
\newunicodechar{Δ}{\ensuremath{\Delta}}
\newunicodechar{Ε}{\ensuremath{\varepsilon}}
\newunicodechar{Θ}{\ensuremath{\Theta}}
\newunicodechar{Λ}{\ensuremath{\Lambda}}
\newunicodechar{Μ}{\ensuremath{\mu}}
\newunicodechar{Τ}{\ensuremath{\tau}}
\newunicodechar{Φ}{\ensuremath{\Phi}}
\newunicodechar{Ψ}{\ensuremath{\Psi}}
\newunicodechar{Ω}{\ensuremath{\Omega}}
\newunicodechar{↦}{\ensuremath{\mapsto}}
\newunicodechar{∈}{\ensuremath{\in}}
\newunicodechar{∉}{\ensuremath{\notin}}
\newunicodechar{∧}{\ensuremath{\wedge}}
\newunicodechar{⇔}{\ensuremath{\Leftrightarrow}}
\newunicodechar{⟺}{\ensuremath{\Longleftrightarrow}}
\newunicodechar{⇒}{\ensuremath{\Rightarrow}}
\newunicodechar{⟹}{\ensuremath{\Longrightarrow}}
\newunicodechar{∘}{\ensuremath{\circ}}
\newunicodechar{∪}{\ensuremath{\cup}}
\newunicodechar{∩}{\ensuremath{\cap}}
\newunicodechar{≡}{\ensuremath{\equiv}}
\newunicodechar{⊂}{\ensuremath{\subset}}
\newunicodechar{⊥}{\ensuremath{\perp}}
\newunicodechar{∃}{\ensuremath{\exists}}
\newunicodechar{∀}{\ensuremath{\forall}}
\newunicodechar{∛}{\ensuremath{\sqrt[3]{\,}}}
\newunicodechar{∜}{\ensuremath{\sqrt[4]{\,}}}
\newunicodechar{⃗}{\ensuremath{{}^{\to}}}
\newunicodechar{ⁿ}{\ifmmode^{n}\else\textsuperscript{\ensuremath{n}}\fi}
\newunicodechar{ˣ}{\ifmmode^{x}\else\textsuperscript{\ensuremath{x}}\fi}
\newunicodechar{ᵇ}{\ifmmode^{b}\else\textsuperscript{\ensuremath{b}}\fi}
\newunicodechar{ʳ}{\ifmmode^{r}\else\textsuperscript{\ensuremath{r}}\fi}
\newunicodechar{ᵘ}{\ifmmode^{u}\else\textsuperscript{\ensuremath{u}}\fi}
\newunicodechar{ʸ}{\ifmmode^{y}\else\textsuperscript{\ensuremath{y}}\fi}
\newunicodechar{ᵃ}{\ifmmode^{a}\else\textsuperscript{\ensuremath{a}}\fi}
\newunicodechar{ᵉ}{\ifmmode^{e}\else\textsuperscript{\ensuremath{e}}\fi}
\newunicodechar{ᵏ}{\ifmmode^{k}\else\textsuperscript{\ensuremath{k}}\fi}
\newunicodechar{ᵖ}{\ifmmode^{p}\else\textsuperscript{\ensuremath{p}}\fi}
\newunicodechar{ᴺ}{\ifmmode^{N}\else\textsuperscript{\ensuremath{N}}\fi}
\newunicodechar{ᶜ}{\ifmmode^{c}\else\textsuperscript{\ensuremath{c}}\fi}
\newunicodechar{ʰ}{\ifmmode^{h}\else\textsuperscript{\ensuremath{h}}\fi}
\newunicodechar{ᵠ}{\ifmmode^{q}\else\textsuperscript{\ensuremath{q}}\fi}
\newunicodechar{ₙ}{\ifmmode_{n}\else\textsubscript{\ensuremath{n}}\fi}
\newunicodechar{ₐ}{\ifmmode_{a}\else\textsubscript{\ensuremath{a}}\fi}
\newunicodechar{ᵢ}{\ifmmode_{i}\else\textsubscript{\ensuremath{i}}\fi}
\newunicodechar{ᵣ}{\ifmmode_{r}\else\textsubscript{\ensuremath{r}}\fi}
\newunicodechar{ₖ}{\ifmmode_{k}\else\textsubscript{\ensuremath{k}}\fi}
\newunicodechar{ᵦ}{\ifmmode_{\beta}\else\textsubscript{\ensuremath{\beta}}\fi}
\newunicodechar{ₓ}{\ifmmode_{x}\else\textsubscript{\ensuremath{x}}\fi}
\newfontfamily\popdisplay{ArchivoBlack-Regular.ttf}[Path=../../../fonts/]
\newfontfamily\poplogo{SpaceGrotesk-Bold.ttf}[Path=../../../fonts/]
\newfontfamily\popsemi{PlusJakartaSans-SemiBold.ttf}[Path=../../../fonts/]
\newfontfamily\popxbold{PlusJakartaSans-ExtraBold.ttf}[Path=../../../fonts/]
\newfontfamily\popxboldls{PlusJakartaSans-ExtraBold.ttf}[Path=../../../fonts/,LetterSpace=3.0]

\setlist[enumerate]{leftmargin=1.7em,itemsep=5pt,topsep=4pt}
\setlist[itemize]{leftmargin=1.7em,itemsep=4pt,topsep=4pt,label={\color{poporange}\textbullet}}
\setlength{\parindent}{0pt}
\setlength{\parskip}{5pt}
\renewcommand{\arraystretch}{1.25}

% pgfplots : style Pop par défaut
\pgfplotsset{
  every axis/.append style={axis line style={popink,line width=1pt},tick style={popink},
    grid style={popline,line width=0.5pt},label style={font=\small},tick label style={font=\footnotesize}},
  popcurve/.style={poporange,line width=1.8pt,no markers,smooth},
}
% Même style disponible dans un \draw TikZ simple (hors pgfplots)
\tikzset{popcurve/.style={poporange,line width=1.8pt}}
\tikzset{popgrid/.style={popline,very thin}}
\colorlet{popcurve}{poporange} % le nom du style est parfois utilisé comme couleur

% ============================================================
% Pill / squiggle
% ============================================================
\newcommand{\poppill}[3]{%
  \tikz[baseline=-0.55ex]\node[draw=popink,line width=1.2pt,rounded corners=2.6pt,%
    fill=#2,inner xsep=6.5pt,inner ysep=3pt]{%
    \popxboldls\fontsize{7}{7}\selectfont\color{#3}\MakeUppercase{#1}};%
}
\newcommand{\popsquiggle}[1]{%
  \tikz[baseline=-0.4ex]{
    \draw[#1,line width=2.6pt,line cap=round]
      plot [smooth] coordinates {(0,0.09) (0.34,0.01) (0.68,0.21) (1.02,0.01) (1.36,0.10)};
  }%
}
% Mot-clé surligné (marqueur orange sous le texte)
\newcommand{\cle}[1]{{\popxbold\color{popdark}#1}}

% ============================================================
% En-tête / pied
% ============================================================
\pagestyle{fancy}
\fancyhf{}
\renewcommand{\headrulewidth}{0pt}
\renewcommand{\footrulewidth}{0pt}
\lhead{\raisebox{-0.15ex}{\poplogo\fontsize{13}{13}\selectfont\color{popink}OnBuch\color{poporange}+}}
\rhead{\poppill{\DOCMATIERE\ \textbullet\ \DOCNIVEAU\ \textbullet\ Chapter \DOCLECON}{white}{popink}}
\lfoot{\popxbold\fontsize{7.6}{7.6}\selectfont\color{popmuted}\MakeUppercase{Signed course OnBuch+}\ \textbullet\ {\popsemi\DOCTITRE}}
\rfoot{\tikz[baseline=-0.6ex]\node[rounded corners=8.5pt,fill=popink,minimum height=17pt,minimum width=17pt,inner xsep=5pt,inner sep=0]{\popxbold\fontsize{8}{8}\selectfont\color{popcream}\thepage\,/\,\pageref*{LastPage}};}

% ============================================================
% Titres
% ============================================================
\newcommand{\chaptitre}[2]{%
  \begin{tcolorbox}[enhanced,colback=poporange,colframe=popink,boxrule=2.6pt,arc=11pt,
      drop shadow=popink,left=15pt,right=15pt,top=12pt,bottom=11pt,before skip=3pt,after skip=18pt]
    \poppill{#2}{popink}{popcream}\par\vspace{9pt}
    {\raggedright\hyphenpenalty=10000\exhyphenpenalty=10000\popdisplay\fontsize{22}{24}\selectfont\color{popcream}\MakeUppercase{#1}\par}
  \end{tcolorbox}
}
% Section numérotée : pastille orange avec le numéro + titre Archivo
\newcounter{popsec}
\newcommand{\coursec}[1]{%
  \stepcounter{popsec}\setcounter{popsub}{0}%
  \par\vspace{16pt}\noindent
  \begin{minipage}{\linewidth}
  \tikz[baseline=-0.7ex]\node[draw=popink,line width=1.6pt,fill=poporange,rounded corners=4pt,minimum size=22pt,inner sep=0]{\popdisplay\fontsize{12}{12}\selectfont\color{popcream}\thepopsec};%
  \hspace{8pt}{\popdisplay\fontsize{15}{17}\selectfont\color{popink}\MakeUppercase{#1}}\par
  \vspace{4pt}\noindent{\color{poporange}\rule{36pt}{3.6pt}}
  \end{minipage}\par\nopagebreak\vspace{8pt}
}
\newcounter{popsub}
\newcommand{\courssub}[1]{%
  \stepcounter{popsub}%
  \par\vspace{9pt}\noindent{\popxbold\fontsize{11.5}{13}\selectfont\color{popink}{\color{poporange}\thepopsec.\thepopsub}\hspace{6pt}#1}\par\nopagebreak\vspace{4pt}
}

% ============================================================
% Boîtes pédagogiques — même recette : fond blanc, bordure épaisse colorée,
% ombre dure encre, badge plein. Titre optionnel [#1].
% ============================================================
\tcbset{popbase/.style={breakable,enhanced,colback=white,boxrule=2.3pt,arc=9pt,
  shadow={3pt}{-3pt}{0pt}{popink},left=14pt,right=14pt,top=11pt,bottom=11pt,before skip=11pt,after skip=15pt,
  fontupper=\popsemi\fontsize{10.6}{14.5}\selectfont\color{popink},
  fontlower=\popsemi\fontsize{10.6}{14.5}\selectfont\color{popink}}}
% Coche dessinée (le glyphe ✓ n'existe pas dans les polices du document)
\newcommand{\popcheck}[1]{\tikz[baseline=-0.5ex]\draw[#1,line width=1.6pt,line cap=round,line join=round](-0.13,0.0)--(-0.04,-0.09)--(0.13,0.1);}
\providecommand{\coche}{\popcheck{popgreen}}
\providecommand{\resultat}[1]{\par\smallskip\noindent\fcolorbox{popgreen}{popgreenL}{\parbox{\dimexpr\linewidth-2\fboxsep-2\fboxrule\relax}{\textbf{Result.} #1}}\par\smallskip}
\newcommand{\popboxhead}[3]{\poppill{#1}{#2}{white}\ifx\relax#3\relax\else\hspace{6pt}{\popxbold\fontsize{10.6}{12}\selectfont\color{popink}#3}\fi\par\vspace{7pt}}

\newtcolorbox{aretenir}[1][]{popbase,colframe=popgreen,before upper={\popboxhead{Key points}{popgreen}{#1}}}
\newtcolorbox{exemplebox}[1][]{popbase,colframe=poporange,before upper={\popboxhead{Exemple}{poporange}{#1}}}
\newtcolorbox{definition}[1][]{popbase,colframe=popblue,before upper={\popboxhead{Definition}{popblue}{#1}}}
\newtcolorbox{propriete}[1][]{popbase,colframe=poppurple,before upper={\popboxhead{Property}{poppurple}{#1}}}
\newtcolorbox{methode}[1][]{popbase,colframe=popgold,before upper={\popboxhead{Method}{popgold}{#1}}}
\newtcolorbox{attention}[1][]{popbase,colframe=poppink,before upper={\popboxhead{Warning}{poppink}{#1}}}
\newtcolorbox{experience}[1][]{popbase,colframe=popblue,colback=popblueL,before upper={\popboxhead{Experiment}{popblue}{#1}}}
% « Le savais-tu ? » : carte violette pleine (contexte, culture, Cameroun)
\newtcolorbox{savaistu}[1][]{popbase,colback=poppurple,colframe=popink,
  fontupper=\popsemi\fontsize{10.4}{14.2}\selectfont\color{white},
  before upper={\poppill{Did you know?}{white}{poppurple}\ifx\relax#1\relax\else\hspace{6pt}{\popxbold\color{white}#1}\fi\par\vspace{7pt}}}
% Exercice résolu : énoncé en haut, \tcblower puis la solution sur fond crème
\newcounter{popexo}\newcounter{popexr}
\newtcolorbox{exoresolu}[1][]{popbase,colframe=poporange,colbacklower=popcream,
  segmentation style={popink,line width=1.2pt,dashed},
  before upper={\stepcounter{popexr}\popboxhead{Worked example \thepopexr}{poporange}{#1}},
  before lower={\poppill{Solution}{popink}{popcream}\par\vspace{6pt}}}
% Exercice (non résolu en place) — la correction est dans la partie Corrigés
\newtcolorbox{exercice}[1][]{popbase,colframe=popink,
  before upper={\stepcounter{popexo}\popboxhead{Exercise \thepopexo}{popink}{#1}}}
% Corrigé
\newtcolorbox{corrige}[1][]{popbase,colframe=popgreen,colback=popgreenL,
  before upper={\popboxhead{Solution}{popgreen}{#1}}}
% Objectifs / prérequis / bilan
\newtcolorbox{objectifs}{popbase,colframe=popink,colback=poporangeL,
  before upper={\popboxhead{Chapter objectives}{popink}{}}}
\newtcolorbox{prerequis}{popbase,colframe=popink,colback=popblueL,
  before upper={\popboxhead{Prerequisites}{popblue}{}}}
\newtcolorbox{bilan}{popbase,colframe=popink,colback=white,boxrule=2.8pt,
  before upper={\popboxhead{Fiche bilan}{poporange}{L'essentiel à connaître pour l'examen}}}
% Figure encadrée avec légende
\newtcolorbox{popfigure}[1][]{enhanced,breakable=false,colback=white,colframe=popline,boxrule=1.4pt,arc=8pt,
  left=10pt,right=10pt,top=10pt,bottom=8pt,before skip=10pt,after skip=12pt,halign=center,
  lower separated=false,halign lower=center,
  fontlower=\popsemi\fontsize{9}{11.5}\selectfont\color{popmuted},
  #1}
\newcommand{\legende}[1]{\par\vspace{6pt}{\popsemi\fontsize{9}{11.5}\selectfont\color{popmuted}#1}}

% ============================================================
% Signature OnBuch+
% ============================================================
\newcommand{\poplogoinline}[1]{{\poplogo\fontsize{#1}{#1}\selectfont\color{popink}OnBuch\color{poporange}+}}
\newcommand{\signatureonbuch}{%
  \begin{tcolorbox}[enhanced,colback=popink,colframe=popink,boxrule=2.2pt,arc=10pt,
      drop shadow=poporange,left=14pt,right=14pt,top=10pt,bottom=10pt,before skip=0pt,after skip=0pt]
    \tikz[baseline=-0.7ex]\node[circle,fill=popgreen,draw=popcream,line width=1.3pt,minimum size=22pt,inner sep=0]{\popcheck{white}};%
    \hspace{9pt}\begin{minipage}[c]{0.62\linewidth}
      {\popxbold\fontsize{12}{13}\selectfont\color{popcream}Signed\ \textbullet\ {\poplogo OnBuch\color{poporange}+}}\par\vspace{2pt}
      {\raggedright\popsemi\fontsize{8}{10.5}\selectfont\color{popline}\MakeUppercase{OnBuch+ teaching team}\\\MakeUppercase{Follows the official MINESEC syllabus}\par}
    \end{minipage}\hfill
    {\poplogo\fontsize{22}{22}\selectfont\color{popcream}OnBuch\color{poporange}+}
  \end{tcolorbox}%
}

% ============================================================
% Page de garde
% #1 titre · #2 matière · #3 niveau · #4 module · #5 n° leçon · #6 durée ·
% #7 description / objectifs (texte ou itemize)
% ============================================================
\newcommand{\pagegarde}[7]{%
  \thispagestyle{empty}
  \newgeometry{a4paper,margin=1.7cm,top=1.6cm,bottom=1.4cm}%
  \begin{tikzpicture}[remember picture,overlay]
    \fill[poporange!18] ([xshift=-3.2cm,yshift=-3.6cm]current page.north east) circle (4.6cm);
    \fill[poppurple!14] ([xshift=2.4cm,yshift=7.6cm]current page.south west) circle (3.8cm);
  \end{tikzpicture}%
  {\poplogo\fontsize{30}{30}\selectfont\color{popink}OnBuch\color{poporange}+}\hfill
  \raisebox{0.6ex}{\poppill{Signed course \textbullet\ Premium}{popink}{popcream}}\par
  \vspace{1pt}\popsquiggle{poppurple}\par
  \vspace{8pt}
  \begin{tcolorbox}[enhanced,colback=poporange,colframe=popink,boxrule=2.8pt,arc=13pt,
      drop shadow=popink,left=18pt,right=18pt,top=12pt,bottom=13pt,after skip=10pt]
    \poppill{#2 \textbullet\ #3}{popink}{popcream}\hspace{5pt}\poppill{#4}{white}{popink}\par\vspace{8pt}
    {\popdisplay\fontsize{14}{14}\selectfont\color{popink}CHAPTER #5}\par\vspace{4pt}
    {\raggedright\hyphenpenalty=10000\exhyphenpenalty=10000\popdisplay\fontsize{25}{27}\selectfont\color{popcream}\MakeUppercase{#1}\par}
    \vspace{8pt}
    {\popxbold\fontsize{9.5}{9.5}\selectfont\color{popink}\MakeUppercase{Suggested time: #6}}
  \end{tcolorbox}
  \begin{tcolorbox}[enhanced,colback=white,colframe=popink,boxrule=2.2pt,arc=10pt,
      drop shadow=popink,left=16pt,right=16pt,top=10pt,bottom=10pt,after skip=8pt]
    \poppill{In this chapter}{poppurple}{white}\par\vspace{5pt}
    \popsemi\fontsize{9.3}{12.6}\selectfont\color{popink}#7
  \end{tcolorbox}
  \noindent\begin{minipage}[t]{0.487\linewidth}
    \begin{tcolorbox}[enhanced,colback=popgreen,colframe=popink,boxrule=2.2pt,arc=10pt,
        drop shadow=popink,left=11pt,right=11pt,top=10pt,bottom=10pt,height=3.1cm,valign=center]
      \tcbox[colback=white,colframe=popink,boxrule=1.3pt,arc=5pt,boxsep=0pt,nobeforeafter,
        left=3pt,right=3pt,top=3pt,bottom=3pt,valign=center]{\qrcode[height=1.9cm]{https://play.google.com/store/apps/details?id=com.luvvix.onbuch&hl=en}}%
      \hspace{8pt}\begin{minipage}[c]{0.5\linewidth}
        {\popdisplay\fontsize{11}{12.5}\selectfont\color{white}\MakeUppercase{Download OnBuch}}\par\vspace{4pt}
        {\raggedright\popsemi\fontsize{8.4}{11}\selectfont\color{white}Free \textbullet\ Google Play\\AI tutor, past papers, results\par}
      \end{minipage}
    \end{tcolorbox}
  \end{minipage}\hfill
  \begin{minipage}[t]{0.487\linewidth}
    \begin{tcolorbox}[enhanced,colback=poppurple,colframe=popink,boxrule=2.2pt,arc=10pt,
        drop shadow=popink,left=11pt,right=11pt,top=10pt,bottom=10pt,height=3.1cm,valign=center]
      \tikz[baseline=-0.7ex]\node[circle,fill=white,draw=popink,line width=1.3pt,minimum size=1.6cm,inner sep=0]{\popdisplay\fontsize{24}{24}\selectfont\color{poppurple}+};%
      \hspace{8pt}\begin{minipage}[c]{0.6\linewidth}
        {\popdisplay\fontsize{11}{12.5}\selectfont\color{white}\MakeUppercase{Go OnBuch+}}\par\vspace{4pt}
        {\raggedright\popsemi\fontsize{8.4}{11}\selectfont\color{white}All signed courses, unlimited AI tutor, PDF sheets — OnBuch+ tab in the app\par}
      \end{minipage}
    \end{tcolorbox}
  \end{minipage}
  \vspace{8pt}
  \popcontenu
  \vfill
  \signatureonbuch
  \restoregeometry
  \newpage
}

% Rangée « Ce cours contient » (page de garde)
\newcommand{\poptile}[3]{% #1 couleur #2 gros texte #3 légende
  \begin{tcolorbox}[enhanced,colback=#1,colframe=popink,boxrule=2pt,arc=9pt,shadow={2.5pt}{-2.5pt}{0pt}{popink},
      left=6pt,right=6pt,top=9pt,bottom=9pt,halign=center,valign=center,height=1.75cm,width=\linewidth,nobeforeafter]
    {\popdisplay\fontsize{12.5}{13}\selectfont\color{white}\MakeUppercase{#2}}\par\vspace{3pt}
    {\centering\popsemi\fontsize{7.8}{9.5}\selectfont\color{white}#3\par}
  \end{tcolorbox}}
\newcommand{\popcontenu}{%
  \par\noindent{\popxboldls\fontsize{7.5}{7.5}\selectfont\color{popmuted}THIS SIGNED COURSE CONTAINS}\par\vspace{7pt}
  \noindent\begin{minipage}[t]{0.235\linewidth}\poptile{popblue}{Course}{complete, illustrated, step by step}\end{minipage}\hfill
  \begin{minipage}[t]{0.235\linewidth}\poptile{poporange}{Methods}{and worked examples}\end{minipage}\hfill
  \begin{minipage}[t]{0.235\linewidth}\poptile{poppink}{Exercises}{exam style + solutions}\end{minipage}\hfill
  \begin{minipage}[t]{0.235\linewidth}\poptile{popgreen}{Summary sheet}{the essentials to remember}\end{minipage}\par}

% Sommaire
\newtcolorbox{sommaire}{enhanced,colback=white,colframe=popink,boxrule=2.2pt,arc=10pt,
  drop shadow=popink,left=18pt,right=18pt,top=13pt,bottom=13pt,before skip=6pt,after skip=20pt,
  before upper={\poppill{Contents}{poppurple}{white}\par\vspace{9pt}\popsemi\fontsize{10.6}{14.5}\selectfont\color{popink}}}

% Signature de fin de cours — poussée tout en bas de la dernière page
\newcommand{\findecours}{%
  \par\vfill\nopagebreak
  \begin{center}\popsquiggle{poporange}\end{center}\vspace{4pt}
  \signatureonbuch
}
\AtBeginDocument{\def\llbracket{\mathopen{[\mkern-3.5mu[}}\def\rrbracket{\mathclose{]\mkern-3.5mu]}}} % intervalles d'entiers (glyphe ⟦ absent de la police)
% Glyphes absents de Fira Math : redéfinis à partir de glyphes disponibles
\AtBeginDocument{%
  \def\setminus{\mathbin{\backslash}}\let\smallsetminus\setminus
  \def\vdots{\mathord{\vbox{\baselineskip3.2pt\lineskiplimit0pt\kern4pt\hbox{.}\hbox{.}\hbox{.}}}}%
  \def\ddots{\mathinner{\mkern1mu\raise6.5pt\hbox{.}\mkern2mu\raise3.6pt\hbox{.}\mkern2mu\raise0.7pt\hbox{.}\mkern1mu}}%
  \def\triangle{\mathord{\scalebox{0.9}{$\symup{\Delta}$}}}%
  \def\nabla{\mathord{\raisebox{\depth}{\scalebox{1}[-1]{$\symup{\Delta}$}}}}%
  \def\checkmark{\popcheck{popdark}}%
  \def\star{\mathbin{\ast}}\def\bigstar{\mathord{\ast}}%
  \def\diamond{\mathbin{\scalebox{0.6}{$\Diamond$}}}%
  \def\bigcup{\mathop{\mathchoice{\raisebox{-0.3em}{\scalebox{1.7}{$\cup$}}}{\scalebox{1.25}{$\cup$}}{\cup}{\cup}}\displaylimits}%
  \def\bigcap{\mathop{\mathchoice{\raisebox{-0.3em}{\scalebox{1.7}{$\cap$}}}{\scalebox{1.25}{$\cap$}}{\cap}{\cap}}\displaylimits}%
  \def\lesseqqgtr{\mathrel{\lessgtr}}\def\bigoplus{\mathop{\oplus}\displaylimits}%
}
\providecommand{\ssi}{if and only if} % abréviation fréquente
\AtBeginDocument{\providecommand{\wideparen}{\overparen}} % arc de cercle

% Fonctions usuelles que les modèles utilisent sans que amsmath les fournisse
\providecommand{\arsinh}{\operatorname{arsinh}}\providecommand{\arcosh}{\operatorname{arcosh}}
\providecommand{\artanh}{\operatorname{artanh}}\providecommand{\arsech}{\operatorname{arsech}}
\providecommand{\arcsch}{\operatorname{arcsch}}\providecommand{\arcoth}{\operatorname{arcoth}}
\providecommand{\arcsinh}{\operatorname{arsinh}}\providecommand{\arccosh}{\operatorname{arcosh}}
\providecommand{\arctanh}{\operatorname{artanh}}\providecommand{\sech}{\operatorname{sech}}
\providecommand{\csch}{\operatorname{csch}}\providecommand{\arcsec}{\operatorname{arcsec}}
\providecommand{\arccsc}{\operatorname{arccsc}}\providecommand{\arccot}{\operatorname{arccot}}
\providecommand{\sgn}{\operatorname{sgn}}\providecommand{\lcm}{\operatorname{lcm}}
\providecommand{\Var}{\operatorname{Var}}\providecommand{\Cov}{\operatorname{Cov}}

\begin{document}

\pagegarde{MAP WORK}{Geology}{Upper Sixth}{Upper Sixth — Geology}{3}{30 periods}{This chapter equips the learner with the complete toolkit of geological map work: reading directions, scales, symbols and strata geometry, solving three-point problems, constructing cross-sections, and interpreting structures and photographs. It is the most heavily weighted practical chapter of the GCE Advanced Level Geology paper.
\vspace{4pt}
\begin{multicols}{2}\small
\begin{itemize}
\item By the end of this chapter I can distinguish the types and categories of geological maps and state their uses.
\item By the end of this chapter I can describe the essential characteristics of a good geological map (title, key, scale, orientation, grid).
\item By the end of this chapter I can determine directions, bearings and grid references on a topographic or geological map.
\item By the end of this chapter I can use map scales to convert map distances to ground distances and vice versa.
\item By the end of this chapter I can recognise and use conventional geological symbols (dip, strike, faults, folds, minerals, boreholes).
\item By the end of this chapter I can describe the geometrical properties of strata (dip, strike, thickness, outcrop width) and apply the rule of V's.
\end{itemize}\end{multicols}}

\begin{sommaire}
\begin{multicols}{2}
\begin{enumerate}
\item Section 1: Types and Categories of Geological Maps
\item Section 2: Characteristics of a Geological Map
\item Section 3: Direction, Orientation and Scale on Maps
\item Section 4: Conventional Geological Symbols on Maps
\item Section 5: Geometrical Properties of Strata
\item Section 6: Strike Lines and the Three-Point Problem
\item Section 7: Construction of Geological Cross-Sections I — Principles and Simple Cases
\item Section 8: Construction of Geological Cross-Sections II — Folds, Faults, Unconformities and Intrusions
\item Section 9: Description of Geological Structures on Maps I — Strata, Dip and Outcrop Patterns
\item Section 10: Description of Geological Structures on Maps II — Folds, Faults, Unconformities, Intrusions and Geological History
\item Section 11: Interpretation of Geological Photographs
\item Section 12: Integration — Complete Map Interpretation and GCE Examination Technique
\item Integration activity
\item Methods and worked examples
\item Exercises
\item Solutions to the exercises
\item Summary sheet
\end{enumerate}
\end{multicols}
\end{sommaire}

\begin{objectifs}
\begin{itemize}
\item By the end of this chapter I can distinguish the types and categories of geological maps and state their uses.
\item By the end of this chapter I can describe the essential characteristics of a good geological map (title, key, scale, orientation, grid).
\item By the end of this chapter I can determine directions, bearings and grid references on a topographic or geological map.
\item By the end of this chapter I can use map scales to convert map distances to ground distances and vice versa.
\item By the end of this chapter I can recognise and use conventional geological symbols (dip, strike, faults, folds, minerals, boreholes).
\item By the end of this chapter I can describe the geometrical properties of strata (dip, strike, thickness, outcrop width) and apply the rule of V's.
\item By the end of this chapter I can draw and use strike lines to determine dip direction and amount, and solve the three-point problem.
\item By the end of this chapter I can construct accurate geological cross-sections from a map along a given line.
\item By the end of this chapter I can describe and interpret geological structures (horizontal, dipping, folded, faulted, unconformable strata) on maps.
\item By the end of this chapter I can interpret geological photographs and deduce the geological history of a mapped area.
\end{itemize}
\end{objectifs}

\begin{prerequis}
\begin{itemize}
\item Rock types and the three rock families (Lower Sixth: rocks and rock cycle)
\item Sedimentary structures: bedding, stratification, fossils (Lower Sixth)
\item Basic structural geology: folds, faults, joints, unconformities (Lower Sixth)
\item Elementary topographic map reading: contours, relief, grid lines (Lower Sixth / Geography)
\item Simple geometry and trigonometry: angles, tangent, similar triangles (Mathematics)
\end{itemize}
\end{prerequis}

\newpage

\coursec{Section 1: Types and Categories of Geological Maps}

During the 2019 Bafoussam landslide, rescue teams and engineers relied on geological maps to identify unstable slopes and safe evacuation zones. When a mining company explores for iron at Mbalam or when a borehole is sited in the Adamawa plateau, the first document consulted is not a machine but a geological map. Yet to most people, the coloured patches and strange symbols of such a map look like abstract art. How can a flat sheet of paper reveal rocks buried hundreds of metres deep, folded mountains, and hidden faults? This chapter turns you into a reader — and a builder — of geological maps.

\courssub{What is a Geological Map?}

\begin{definition}[Geological Map]
A \cle{geological map} is a representation of the distribution, nature, age and structural relationships of rock units and geological structures at the Earth's surface, plotted on a topographic base map. It shows \emph{what} rocks are present (lithology), \emph{how old} they are (stratigraphy), and \emph{how they are arranged} (structure: dip, strike, folds, faults, unconformities).
\end{definition}

Unlike a topographic map which shows only the shape of the land surface (contours, rivers, roads), a geological map interprets the \emph{subsurface geology} from surface observations. Each coloured polygon represents a \cle{rock unit} (formation or group) with a specific age and lithology. Lines represent boundaries between units (contacts) or structural features (faults, fold axes). Symbols give the orientation of planar features (bedding, foliation, cleavage) at specific localities.

\begin{attention}[Topographic Map vs Geological Map]
\textbf{Do not confuse the two.} A topographic map is the \emph{base}; the geological map is the \emph{interpretation} overlain on that base. In the GCE examination, you will be given a geological map extract \emph{with} topographic contours. You must use both: contours for cross-section construction and relief interpretation; geological colours and symbols for structure and history.
\end{attention}

\courssub{Classification by Scale}

The \cle{scale} of a map is the ratio of a distance on the map to the corresponding distance on the ground. It determines the level of detail, the area covered, and the geological problems the map can address.

\begin{center}
\begin{tabularx}{\linewidth}{|l|X|X|X|}
\hline
\rowcolor{popblueL}
\textbf{Scale Category} & \textbf{Representative Fraction (RF)} & \textbf{Typical Area Covered} & \textbf{Main Uses} \\
\hline
\cle{Small-scale} (Regional) & 1:500,000 to 1:1,000,000 & Whole country or large region (e.g. Cameroon) & Regional synthesis, tectonic overview, planning national infrastructure, teaching \\
\hline
\cle{Medium-scale} & 1:100,000 to 1:250,000 & Geological province or basin (e.g. Benue Trough sheet) & Basin analysis, groundwater surveys, mineral exploration targeting, road corridor planning \\
\hline
\cle{Large-scale} (Detailed) & 1:10,000 to 1:50,000 & Mining lease, dam site, engineering project & Mine planning, foundation engineering, detailed structural analysis, ore reserve estimation \\
\hline
\end{tabularx}
\end{center}

\begin{savaistu}[Cameroon's Geological Maps]
The \emph{Geological Map of Cameroon} at 1:1,000,000 (small-scale) was first published in 1968 and updated in 1991. The \emph{Bureau des Recherches Géologiques et Minières (BRGM)} and later the \emph{Institut de Recherche Géologique et Minière (IRGM)} produced 1:200,000 sheets (medium-scale) covering the whole country. For mining (e.g. Mbalam iron ore, Nkamouna cobalt-nickel) and engineering (e.g. Lom Pangar dam, Nachtigal hydroelectric project), 1:50,000 or 1:25,000 maps are produced by companies or consultants.
\end{savaistu}

\courssub{Classification by Purpose}

Geological maps are also categorised by \emph{what they emphasise}. The same topographic base can carry different geological themes.

\begin{definition}[Solid Geology Map vs Drift (Superficial) Map]
A \cle{solid geology map} (or bedrock map) shows the \emph{in-situ} bedrock formations, ignoring or only thinly representing the unconsolidated cover (soil, alluvium, laterite, volcanic ash). A \cle{drift map} (superficial deposits map) shows the distribution and thickness of \emph{unconsolidated Quaternary deposits} (alluvium, colluvium, aeolian sand, lateritic duricrust, volcanic ash) that overlie the bedrock. In tropical Cameroon, lateritic cover is thick; a solid map infers bedrock from sparse outcrops, while a drift map is essential for agriculture, groundwater recharge zones, and foundation conditions.
\end{definition}

\begin{center}
\begin{tabularx}{\linewidth}{|l|X|X|}
\hline
\rowcolor{popblueL}
\textbf{Map Type} & \textbf{Emphasis} & \textbf{Typical Application} \\
\hline
\cle{Lithological map} & Rock type (granite, sandstone, basalt, limestone) & Quarrying, aggregate search, soil parent material \\
\cle{Structural map} & Dip/strike, folds, faults, lineaments & Tectonic analysis, hydrocarbon trapping, seismic hazard \\
\cle{Stratigraphic map} & Age and correlation of units & Basin evolution, fossil correlation, sequence stratigraphy \\
\cle{Hydrogeological map} & Aquifers, recharge/discharge, water quality, piezometric surface & Borehole siting, water supply management (e.g. Adamawa plateau) \\
\cle{Engineering geology map} & Rock mass rating, landslide susceptibility, excavatability, foundation conditions & Road design (e.g. Bamenda–Enugu), dam foundations, urban planning \\
\cle{Mineral / Economic geology map} & Mineral occurrences, mines, alteration zones, geochemical anomalies & Exploration targeting (Mbalam Fe, Nkamouna Co-Ni, Mobilong diamond) \\
\cle{Geomorphological map} & Landforms, processes, erosion surfaces & Landscape evolution, hazard mapping (Lake Nyos area) \\
\hline
\end{tabularx}
\end{center}

\courssub{Classification by Production Method}

How the data were gathered affects the map's reliability and resolution.

\begin{center}
\begin{tabularx}{\linewidth}{|l|X|X|}
\hline
\rowcolor{popblueL}
\textbf{Method} & \textbf{Description} & \textbf{Strengths / Limitations} \\
\hline
\cle{Field survey map} & Geologist walks traverses, records outcrops, measures dip/strike, collects samples. Ground truth. & Highest reliability for lithology and structure. Slow, expensive, limited by access and vegetation (thick rainforest in South Cameroon). \\
\cle{Photogeological map} & Interpretation of aerial photographs or satellite imagery (stereo pairs, tone, texture, lineaments) with limited ground control. & Rapid coverage of large/inaccessible areas. Good for structural lineaments and geomorphology. Poor for lithology in thick soil/vegetation. \\
\cle{Geophysical / Remote-sensing derived map} & Airborne magnetics, radiometrics, EM, gravity; Landsat, Sentinel, ASTER, radar (SRTM, ALOS PALSAR). & Reveals buried structures, lithological contrasts under cover. Requires calibration with field data. Not a substitute for geological mapping. \\
\hline
\end{tabularx}
\end{center}

\begin{aretenir}[Modern Mapping Workflow]
In practice, a reliable geological map combines all three: \textbf{remote sensing} for regional structure and targeting, \textbf{photogeology} for preliminary unit delineation, and \textbf{field survey} for ground truth, detailed structure, and sampling. The GCE expects you to understand that a map's legend usually states the method and date of survey.
\end{aretenir}

\courssub{The Geological Map of Cameroon: A Case Study}

Cameroon's geology is a textbook example of African geology: a Pan-African basement, Mesozoic rift basins, and a Cenozoic volcanic line. A small-scale map (1:1,000,000) captures these first-order units.

\begin{popfigure}
\begin{tikzpicture}[scale=0.85]
% Cameroon outline simplified
\draw[popdark, thick] (0,0) -- (1,0) -- (1.2,1) -- (1.5,2) -- (1.3,3) -- (1.5,4) -- (1.2,5) -- (0.8,5.5) -- (0.5,5) -- (0.2,4.5) -- (0,4) -- (0.1,3) -- (0,2) -- (0.2,1) -- (0,0);
% Major units as filled regions
% 1. Pan-African Basement (South & East) - light gray
\fill[popmuted!30] (0.1,0.2) -- (0.9,0.2) -- (1.0,1.2) -- (1.1,2.5) -- (0.9,3.5) -- (0.5,4.2) -- (0.2,3.8) -- (0.1,2.5) -- cycle;
% 2. Congo Craton (Far SE) - light yellow
\fill[popgold!40] (0.8,0.2) -- (1.2,0.2) -- (1.2,1) -- (1.0,1.2) -- cycle;
% 3. Cretaceous Basins (North: Benue Trough, Logone-Birni) - light blue
\fill[popblueL!50] (0.2,3.8) -- (0.9,3.5) -- (1.1,2.5) -- (1.2,1.5) -- (1.5,2) -- (1.3,3) -- (1.5,4) -- (1.2,5) -- (0.8,5.5) -- (0.5,5) -- cycle;
% 4. Cameroon Volcanic Line (SW-NE chain) - pink
\fill[poppink!40] (0.5,0.5) -- (1.0,1.5) -- (1.2,2.5) -- (1.0,3.5) -- (0.6,3.0) -- (0.3,2.0) -- cycle;
% Mount Cameroon
\fill[poporange] (0.5,0.5) circle (0.15);
\node[popdark, font=\tiny\bfseries] at (0.5,0.5) {Mt C};
% Lake Nyos
\fill[popblue] (0.9,3.2) circle (0.07);
\node[popdark, font=\tiny] at (0.9,3.5) {L. Nyos};
% Legend
\begin{scope}[xshift=7cm, yshift=0cm]
\node[font=\bfseries, popdark] at (0,3.5) {Legend};
\fill[popmuted!30] (0,3) rectangle (0.6,3.4); \node[anchor=west, font=\small] at (0.7,3.2) {Pan-African Basement (600-500 Ma)};
\fill[popgold!40] (0,2.5) rectangle (0.6,2.9); \node[anchor=west, font=\small] at (0.7,2.7) {Congo Craton (Archean-Paleoproterozoic)};
\fill[popblueL!50] (0,2) rectangle (0.6,2.4); \node[anchor=west, font=\small] at (0.7,2.2) {Cretaceous Basins (Benue Trough, Logone-Birni)};
\fill[poppink!40] (0,1.5) rectangle (0.6,1.9); \node[anchor=west, font=\small] at (0.7,1.7) {Cameroon Volcanic Line (Cenozoic)};
\fill[poporange] (0,1) circle (0.1); \node[anchor=west, font=\small] at (0.2,1) {Mount Cameroon (active)};
\fill[popblue] (0,0.5) circle (0.07); \node[anchor=west, font=\small] at (0.2,0.5) {Lake Nyos (CO$_2$ hazard)};
\end{scope}
\end{tikzpicture}
\legende{Simplified geological sketch map of Cameroon showing four major units. Not to scale; for illustration only.}
\end{popfigure}

\begin{savaistu}[Cameroon's Four First-Order Units]
\begin{enumerate}
\item \textbf{Pan-African Basement (600–500 Ma)}: Gneisses, schists, granites covering most of South and East Cameroon. Hosts gold (Bétaré-Oya), uranium (Kitongo), and iron (Mbalam — actually in the Congo Craton margin).
\item \textbf{Congo Craton (Archean–Paleoproterozoic)}: Stable block in far SE (Ntem complex). Oldest rocks in Cameroon.
\item \textbf{Cretaceous Sedimentary Basins}: Benue Trough (North, oil/gas potential), Logone-Birni Basin (Far North, groundwater), Douala/Kribi-Campo Basins (coastal, hydrocarbons).
\item \textbf{Cameroon Volcanic Line (CVL)}: Y-shaped chain from Pagalu (Annobón) through Mt Cameroon, Manengouba, Bambouto, Oku, to Biu Plateau (Nigeria). Cenozoic to Recent. Mt Cameroon (4,095 m) is the only currently active volcano. Lake Nyos (1986 CO$_2$ disaster) lies on the CVL.
\end{enumerate}
\end{savaistu}

\courssub{Choosing the Right Map for the Job}

\begin{methode}[Selecting a Geological Map Scale and Type for a Project]
\begin{enumerate}
\item \textbf{Define the objective:} Regional reconnaissance? Detailed engineering? Mineral exploration? Groundwater?
\item \textbf{Determine the required ground resolution:} What is the smallest feature you must resolve? (e.g. a 10 m wide vein needs 1:10,000; a regional fold trend needs 1:200,000).
\item \textbf{Match scale to objective:}
\begin{itemize}
\item National/regional planning $\rightarrow$ 1:500,000–1:1,000,000 (solid geology, structural)
\item Basin evaluation / groundwater province $\rightarrow$ 1:200,000 (hydrogeological, stratigraphic)
\item Feasibility study (dam, highway, mine) $\rightarrow$ 1:25,000–1:50,000 (engineering geology, structural, lithological)
\item Mine design / ore reserve $\rightarrow$ 1:5,000–1:10,000 (detailed structural, grade control)
\end{itemize}
\item \textbf{Choose the thematic type:} Solid geology for bedrock structure; drift for foundation soils; hydrogeological for water; mineral for exploration.
\item \textbf{Check metadata:} Survey date, method (field vs photo), contour interval, coordinate system. Older maps may lack GPS control.
\end{enumerate}
\end{methode}

\begin{exemplebox}[Map Selection Scenarios]
\begin{itemize}
\item \textbf{Scenario A:} Siting a borehole for a village in the Adamawa plateau. \emph{Use:} 1:200,000 hydrogeological map (aquifer distribution, recharge zones) + 1:50,000 solid geology map (fracture zones from lineaments).
\item \textbf{Scenario B:} Designing foundations for the Lom Pangar dam. \emph{Use:} 1:25,000 engineering geology map (rock mass quality, fault zones, permeability) + large-scale drift map (alluvium thickness).
\item \textbf{Scenario C:} Exploring for cobalt-nickel laterite at Nkamouna. \emph{Use:} 1:50,000 regolith/geochemical map + 1:10,000 detailed drilling map (laterite profile, grade).
\end{itemize}
\end{exemplebox}

\begin{aretenir}[GCE Exam Tip: Uses of Geological Maps]
A classic GCE question: \emph{``State two uses of geological maps.''} Give \textbf{specific, geological} answers, not vague ones.
\begin{itemize}
\item \textbf{Good:} ``To identify aquifer units for borehole siting''; ``To locate fault zones for seismic hazard assessment''; ``To delineate mineralised zones for exploration targeting''; ``To reconstruct the geological history of an area''.
\item \textbf{Weak:} ``To know where rocks are''; ``For planning''; ``For mining''.
\end{itemize}
Learn 4–5 distinct uses with a Cameroon example each.
\end{aretenir}

\courssub{Worked Example: Interpreting Map Metadata}

\begin{exoresolu}[Reading a Map Legend]
\textbf{Statement:} You are given a geological map extract with the following legend entries:
\begin{itemize}
\item \emph{Surveyed by IRGM, 1985–1987; compiled from 1:50,000 field sheets.}
\item \emph{Topographic base: 1:200,000, UTM Zone 32N, WGS84.}
\item \emph{Geology: Solid geology. Drift >2 m not shown.}
\item \emph{Contour interval: 50 m.}
\end{itemize}
Answer the following:
\begin{enumerate}
\item What is the publication scale of this map?
\item What does ``Drift >2 m not shown'' imply for engineering geology?
\item Why is the coordinate system (UTM Zone 32N, WGS84) important?
\end{enumerate}

\tcblower
\textbf{Detailed solution:}
\begin{enumerate}
\item The map is published at \textbf{1:200,000} (medium-scale). The field sheets were 1:50,000 (large-scale), but they were \emph{compiled} (generalised) to 1:200,000. Small features visible at 1:50,000 (thin dykes, small faults) may be omitted.
\item \textbf{Implication:} Any unconsolidated cover thicker than 2 m (laterite, alluvium, colluvium) is \emph{not} represented. The map shows bedrock as if it crops out everywhere. For engineering (foundations, excavations), you \textbf{must} obtain a separate drift map or conduct site investigation (augering, trial pits) to know the thickness and properties of the overburden. This is critical in Cameroon where laterite can be 20–30 m thick.
\item The coordinate system allows \textbf{precise location} of any point (e.g. a borehole, an outcrop) using GPS. UTM Zone 32N covers Cameroon (approx. 6°E–12°E). WGS84 is the datum used by GPS. If you plot GPS coordinates on a map with a different datum (e.g. Clarke 1880 / Adindan), you could be off by 100–200 m.
\end{enumerate}
\end{exoresolu}

\courssub{Summary}

\begin{aretenir}[Key Points of Section 1]
\begin{itemize}
\item A geological map shows \textbf{lithology, age, and structure} on a topographic base.
\item \textbf{Scale} controls detail: small (regional), medium (basin/province), large (site/mine).
\item \textbf{Purpose} controls theme: solid vs drift, structural, lithological, hydrogeological, engineering, mineral, geomorphological.
\item \textbf{Method} controls reliability: field survey (best), photogeology (fast, structural), geophysics/remote sensing (buried features).
\item Cameroon's geology at 1:1,000,000 shows four first-order units: Pan-African Basement, Congo Craton, Cretaceous Basins, Cameroon Volcanic Line.
\item Always read the \textbf{legend metadata} (scale, date, method, drift treatment, coordinate system) before interpreting.
\end{itemize}
\end{aretenir}

\begin{exercice}[Practice Questions]
\begin{enumerate}
\item Define a geological map and distinguish it from a topographic map.
\item A mining company holds a 1:200,000 geological map of a concession. They need to plan drill holes to test a 50 m wide vein. Is this map adequate? Justify your answer and state what scale they should request.
\item Explain the difference between a solid geology map and a drift map. Why is this distinction particularly important in Cameroon?
\item List three uses of a hydrogeological map and three uses of an engineering geology map.
\item The legend of a map states: ``Geology compiled from Landsat 8 imagery with 10\% field check.'' What are the likely strengths and weaknesses of this map for structural analysis?
\end{enumerate}
\end{exercice}

\begin{corrige}[Exercise 1]
\begin{enumerate}
\item A geological map represents the distribution, nature, age and structural relationships of rock units at the Earth's surface on a topographic base. A topographic map shows only the shape of the land surface (relief, drainage, cultural features) without geological interpretation.
\item No. At 1:200,000, 1 mm = 200 m on the ground. A 50 m wide vein would be 0.25 mm wide — below the limit of graphical resolution (typically 0.5 mm). They need a \textbf{1:25,000 or 1:10,000} map where 50 m = 2–5 mm, clearly plottable.
\item Solid map = bedrock geology (in-situ formations). Drift map = unconsolidated superficial deposits (alluvium, laterite, colluvium). In Cameroon, thick lateritic weathering (often >10 m) masks bedrock; a solid map infers geology from sparse outcrops, while a drift map shows the actual surface material critical for agriculture, groundwater recharge, and foundation engineering.
\item Hydrogeological: (1) Aquifer delineation for borehole siting, (2) Recharge/discharge zone mapping for water resource management, (3) Groundwater quality zoning. Engineering: (1) Rock mass classification for slope/tunnel stability, (2) Landslide susceptibility mapping for route selection, (3) Excavatability and foundation condition assessment for dams/roads.
\item Strengths: Good regional structural lineaments (faults, dykes, fold axes) visible on imagery; rapid coverage. Weaknesses: Lithological boundaries often inaccurate under vegetation/soil; dip/strike measurements absent or inferred; small structures (< imagery resolution) missed; 10\% field check means 90\% unchecked — high risk of misinterpretation.
\end{enumerate}
\end{corrige}

\coursec{Section 2: Characteristics of a Geological Map}

In Section 1 we classified geological maps according to their purpose, their scale and the kind of information they carry. We saw that a small-scale reconnaissance map of the whole of Cameroon and a detailed 1/25\,000 sheet of the Mount Cameroon region are very different documents. Yet, whatever its category, every geological map must obey the same \emph{grammar}: it must contain a fixed set of elements, arranged according to well-established conventions, so that any geologist anywhere in the world can read it. This section studies that grammar — the \cle{characteristics} that make a document a genuine geological map rather than a mere coloured drawing.

\courssub{Essential Elements of a Geological Map}

Imagine you are handed a coloured sheet with no title, no scale and no north arrow. You see patches of green and pink, lines and symbols — but you cannot say where the area is, how big it is, or what the colours mean. The document is useless. This simple thought experiment shows why a geological map is more than its coloured face: it is a \emph{complete scientific document} whose marginal elements are as important as the map itself.

\begin{definition}[Geological map]
A \cle{geological map} is a cartographic document that represents, on a topographic base, the distribution at the Earth's surface of rock units (formations), their boundaries (\cle{contacts}) and their structures (dip, folds, faults, unconformities, intrusions), together with the marginal information needed to locate, measure and interpret them.
\end{definition}

Every complete geological map must possess the following \textbf{essential elements}:

\begin{itemize}
\item \textbf{Title}: states the name of the area mapped, the sheet number and the type of map (e.g. \emph{``Geological map of the Buea sheet, 1/50\,000''}). It immediately tells the reader what the document is about.
\item \textbf{Legend (key)}: the ordered list of all rock units shown, with their colours, symbols and ages (developed below).
\item \textbf{Scale}: the ratio between distances on the map and distances on the ground (studied in detail in Section 3). It may be given as a fraction ($1/50\,000$), in words, or graphically as a \cle{scale bar}.
\item \textbf{North arrow}: indicates the orientation of the sheet; usually geographic (true) north, sometimes with the magnetic declination stated.
\item \textbf{Grid or graticule}: a network of lines of longitude and latitude (graticule) or of rectangular coordinates (grid, e.g. UTM) allowing every point to be located numerically.
\item \textbf{Author and date of survey}: the geologist(s) or institution that carried out the mapping and the year of publication — essential for judging the reliability and the vintage of the data.
\end{itemize}

\begin{attention}[A map without a scale or a north arrow]
In the GCE examination, a map extract is sometimes presented without one of these elements. Never assume the top of the page is north, and never estimate distances without checking the scale bar. The absence of an essential element is itself a piece of information worth commenting on.
\end{attention}

\courssub{The Legend (Key) of a Geological Map}

The legend is the \emph{dictionary} of the map: without it, the colours are meaningless.

\begin{definition}[Legend of a geological map]
The \cle{legend} (or \cle{key}) of a geological map is the marginal panel that lists every rock unit represented on the map, each identified by a colour, an ornament (pattern) and a letter symbol, and arranged as a \cle{stratigraphic column} with the \textbf{oldest unit at the bottom} and the \textbf{youngest at the top}, reproducing the order of superposition of the strata.
\end{definition}

Three conventions govern the legend:

\begin{enumerate}
\item \textbf{Vertical ordering by age.} Units are stacked exactly as they occur in nature: the legend is a miniature stratigraphic column. Reading it from bottom to top is travelling forward through geological time.
\item \textbf{Colour code.} Each unit carries the colour internationally assigned to its geological period (see the property below). Within a period, different formations are distinguished by lighter or darker shades of the same colour.
\item \textbf{Ornament and letter symbols.} Lithology is suggested by patterns (dots for sandstone, brick pattern for limestone, crosses or ``v'' for igneous rocks), and each unit bears an abbreviation, usually the initial of the period plus a number (e.g. K1, K2 for Lower and Upper Cretaceous).
\end{enumerate}

\begin{propriete}[International chronostratigraphic colour code]
By international convention (Commission for the Geological Map of the World), each geological period is assigned a standard colour: \textbf{Quaternary} is yellow, \textbf{Neogene/Palaeogene} (Tertiary) yellow-orange to orange, \textbf{Cretaceous} green, \textbf{Jurassic} blue, \textbf{Triassic} violet-mauve, \textbf{Palaeozoic} systems range from blue-green (Permian, Carboniferous) to brown and olive (older systems), and \textbf{Precambrian basement} is shown in pink, red or dark brown. Igneous intrusions are conventionally shown in vivid red or pink with appropriate ornament. This code is not examinable as a proof, but you must be able to \emph{use} it to estimate the age of a unit at a glance.
\end{propriete}

\begin{savaistu}[Colours on Cameroonian maps]
On the geological map of Cameroon, the vast Precambrian basement of the Centre and South (the northern edge of the Congo craton) appears in pink and red tones, the Cretaceous sediments of the Benue trough and the coastal basin in green, and the young volcanics of the Cameroon Volcanic Line — Mount Cameroon, the Manengouba–Bambouto highlands, the Adamawa plateau — in dark red to brown. A trained eye can thus sketch the geological skeleton of the whole country from colours alone.
\end{savaistu}

\begin{attention}[Oldest at the bottom!]
The single most frequent mistake in map reading is to read the legend from top to bottom and to declare the uppermost unit the oldest. The legend is a stratigraphic column: the \textbf{oldest unit is at the BOTTOM}, the youngest at the TOP. Always check the legend's order before writing anything about relative ages or geological history.
\end{attention}

\begin{popfigure}
\begin{center}
\begin{tikzpicture}[scale=1.0, font=\small]
% map frame
\draw[thick, popink] (0,0) rectangle (12,8.6);
% map face
\draw[popline, fill=popblueL!40] (0.3,0.3) rectangle (8.2,6.6);
\node[popmuted] at (4.25,6.3) {\textbf{MAP FACE}};
% some geology on the face
\draw[thick, popgreen] (0.3,2.2) .. controls (2.5,2.8) and (4,1.6) .. (8.2,2.6);
\draw[thick, popgreen] (0.3,3.6) .. controls (3,4.4) and (5,3.2) .. (8.2,4.2);
\draw[thick, poppink] (0.3,5.0) .. controls (3,5.6) and (5.5,4.6) .. (8.2,5.4);
\draw[thick, poporange] (5.2,0.3) -- (4.6,6.6);
\node[poporange, rotate=84] at (4.75,3.4) {\scriptsize fault};
% north arrow
\draw[-latex, very thick, popdark] (8.7,7.2) -- (8.7,8.3);
\node[popdark] at (8.7,8.45) {\textbf{N}};
\node[popmuted, right] at (8.85,7.7) {\scriptsize north arrow};
% title
\draw[fill=popgold!30, draw=popline] (0.3,7.0) rectangle (7.6,8.3);
\node[align=center] at (3.95,7.65) {\textbf{TITLE}: Geological map of the X sheet\\ \scriptsize author, institution, date of survey, sheet number};
% legend
\draw[fill=white, draw=popline] (8.6,0.3) rectangle (11.7,6.6);
\node at (10.15,6.3) {\textbf{LEGEND}};
\fill[poppink] (8.75,5.5) rectangle (9.25,5.85); \node[right, font=\scriptsize] at (9.3,5.67) {granite (basement)};
\fill[popblue] (8.75,4.85) rectangle (9.25,5.2); \node[right, font=\scriptsize] at (9.3,5.02) {limestone (J)};
\fill[popgreen] (8.75,4.2) rectangle (9.25,4.55); \node[right, font=\scriptsize] at (9.3,4.37) {sandstone (K)};
\fill[popgold] (8.75,3.55) rectangle (9.25,3.9); \node[right, font=\scriptsize] at (9.3,3.72) {alluvium (Q)};
\draw[thick, poporange] (8.75,3.0) -- (9.25,3.0); \node[right, font=\scriptsize] at (9.3,3.0) {fault};
\draw[thick, popink] (8.75,2.5) -- (9.25,2.5) node[midway, sloped] {\tiny$\triangleright$};
\node[right, font=\scriptsize] at (9.3,2.5) {dip $30^{\circ}$};
\node[font=\scriptsize, align=left, text width=2.6cm] at (10.15,1.5) {oldest at bottom, youngest at top};
% scale bar
\draw[thick, popink] (0.5,0.75) rectangle (1.5,0.95);
\draw[thick, popink, fill=popink] (1.5,0.75) rectangle (2.5,0.95);
\node[font=\scriptsize] at (0.5,0.55) {0};
\node[font=\scriptsize] at (2.5,0.55) {10 km};
\node[font=\scriptsize, popmuted] at (3.6,0.85) {scale bar};
% cross-section inset
\draw[fill=white, draw=popline] (4.6,0.5) rectangle (8.0,2.0);
\draw[thick, popink] (4.8,1.5) -- (7.8,1.5);
\draw[popgreen, thick] (4.8,1.5) -- (6.3,0.8);
\draw[poppink, thick] (6.3,0.8) -- (7.8,0.8);
\node[font=\scriptsize, popmuted] at (6.3,1.8) {cross-section inset};
% grid ticks
\foreach \x in {2,4,6} {\draw[popmuted] (\x,0.3) -- (\x,0.45); \draw[popmuted] (\x,6.45) -- (\x,6.6);}
\node[font=\scriptsize, popmuted, rotate=90] at (0.15,3.4) {graticule};
\end{tikzpicture}
\end{center}
\legende{Mock-up of a geological map frame showing the essential marginal elements: title block, legend, scale bar, north arrow, graticule and cross-section inset.}
\end{popfigure}

\courssub{The Stratigraphic Legend Column in Practice}

The figure below shows a typical legend column for six units, using the standard colours. Note how the ages decrease upwards and how the ornaments hint at the lithology.

\begin{popfigure}
\begin{center}
\begin{tikzpicture}[scale=1.0, font=\small]
% unit 6 youngest - Quaternary
\fill[popgold] (0,5.0) rectangle (3,6.0);
\node[right] at (3.15,5.5) {6. Alluvium, sands and gravels --- \textbf{Quaternary}};
% unit 5 - Cretaceous sandstone
\fill[popgreen] (0,4.0) rectangle (3,5.0);
\foreach \x in {0.3,0.9,1.5,2.1,2.7} \foreach \y in {4.25,4.75} \fill[popdark] (\x,\y) circle (0.03);
\node[right] at (3.15,4.5) {5. Sandstone (dotted ornament) --- \textbf{Cretaceous}};
% unit 4 - Cretaceous limestone
\fill[popgreen!60] (0,3.0) rectangle (3,4.0);
\foreach \x in {0.4,1.4,2.4} {\draw[popdark] (\x,3.15) rectangle (\x+0.4,3.45); \draw[popdark] (\x+0.2,3.6) rectangle (\x+0.6,3.9);}
\node[right] at (3.15,3.5) {4. Limestone (brick ornament) --- \textbf{Cretaceous}};
% unit 3 - Jurassic shale
\fill[popblue] (0,2.0) rectangle (3,3.0);
\foreach \y in {2.2,2.5,2.8} \draw[white, thin] (0.1,\y) -- (2.9,\y);
\node[right] at (3.15,2.5) {3. Shale (thin bedding lines) --- \textbf{Jurassic}};
% unit 2 - dyke
\fill[poppink!85!black] (0,1.0) rectangle (3,2.0);
\foreach \x in {0.3,0.8,1.3,1.8,2.3,2.8} \node[white, font=\scriptsize] at (\x,1.5) {v};
\node[right] at (3.15,1.5) {2. Dolerite dyke (``v'' ornament) --- intrusive};
% unit 1 oldest - basement
\fill[poppink] (0,0.0) rectangle (3,1.0);
\foreach \x in {0.4,1.2,2.0,2.8} \node[popdark, font=\scriptsize] at (\x,0.5) {+};
\node[right] at (3.15,0.5) {1. Granite--gneiss basement (crosses) --- \textbf{Precambrian}};
% age arrow
\draw[-latex, very thick, poppurple] (-0.5,0.2) -- (-0.5,5.8);
\node[rotate=90, poppurple] at (-0.85,3.0) {decreasing age};
\draw[thick, popink] (0,0) rectangle (3,6.0);
\end{tikzpicture}
\end{center}
\legende{A stratigraphic legend column of six units with standard international colours and lithological ornaments. The oldest unit (Precambrian basement) lies at the bottom, the youngest (Quaternary alluvium) at the top.}
\end{popfigure}

\courssub{Marginal Information and Accompanying Documents}

Beyond the essential elements, professional map sheets carry additional marginal information:

\begin{itemize}
\item \textbf{Index to adjoining sheets}: a small diagram showing the position of the sheet among its neighbours, so that several sheets can be assembled to cover a larger region.
\item \textbf{Location map}: a small inset of the country or region with the mapped area outlined — for example, an inset of Cameroon situating the Buea sheet in the South-West Region.
\item \textbf{Reliability diagram}: a panel showing which parts of the sheet were mapped in detail on the ground, which were interpolated from aerial photographs, and which are uncertain. It warns the reader where contacts are observed and where they are merely inferred.
\item \textbf{Cross-sections printed on the margin}: one or more vertical sections along lines drawn on the map face (e.g. A--B), giving the third dimension of the geology (the construction of such sections is the subject of Sections 7 and 8).
\item \textbf{Explanatory memoir}: a booklet published with the sheet, describing each formation (lithology, thickness, fossils, age), the structure and the geological history of the area, together with the methods used during the survey.
\end{itemize}

\courssub{Quality Criteria of a Good Geological Map}

Not all maps are equally trustworthy. When you assess a geological map — a skill directly tested at the GCE — ask the following questions:

\begin{enumerate}
\item \textbf{Accuracy of contacts}: are the boundaries between units precisely drawn and consistent with the topography (do they form sensible V-shapes across valleys)? Are observed contacts distinguished from inferred (dashed) ones?
\item \textbf{Completeness of structural data}: are there enough dip and strike symbols, spread over the whole sheet, to constrain the geometry? A map with beautiful colours but no dip measurements is of little use for cross-section work.
\item \textbf{Clarity of symbols and legend}: are the colours distinct, the ornaments legible, the symbols standard? Is every colour and line on the map face explained in the legend?
\item \textbf{Internal consistency}: do the marginal cross-sections agree with the map face? Does the legend order agree with the contacts observed in the field?
\item \textbf{Provenance}: is the author identified, the survey dated, the base map acknowledged?
\end{enumerate}

\courssub{Reading Order: First Approach to an Unknown Map}

Faced with an unknown map in the examination hall, many candidates plunge straight into the colours and get lost. Follow a disciplined order instead.

\begin{methode}[The five-step first reading of a geological map]
\begin{enumerate}
\item \textbf{Read the title and marginal information}: identify the area, the author, the date, and note any cross-section lines (A--B) drawn on the face.
\item \textbf{Check the scale}: convert it into a ground distance (e.g. $1/50\,000$: 1 cm = 500 m) and get a feeling for the size of the area.
\item \textbf{Study the legend}: list the units from oldest (bottom) to youngest (top); note their lithologies and colours; locate each colour on the map face.
\item \textbf{Examine the structures}: find the dip symbols, fold axes, faults, unconformities and intrusions; note their orientation using the north arrow.
\item \textbf{Read the relief}: use contour lines, rivers and spot heights to see how the outcrop pattern responds to topography (rule of V's, developed in Section 9).
\end{enumerate}
Only after these five steps should you attempt any interpretation or geological history.
\end{methode}

\begin{exoresolu}[First reading of a map extract]
Statement: A map extract is entitled ``Geological map of the Lake Nyos area, 1/50\,000, surveyed 1984''. Its legend shows, from bottom to top: (1) Precambrian granite--gneiss (pink, crosses), (2) Cretaceous sandstone (green, dots), (3) Quaternary basalt flows and pyroclasts (dark red). A north arrow, a scale bar and a cross-section A--B are provided. (a) State the essential elements present. (b) Give the relative ages of the units. (c) What does the colour of unit 2 suggest, and why?
\tcblower
\textbf{Solution.}
\begin{enumerate}
\item \textbf{(a)} Essential elements present: title (area, scale, date of survey), legend (stratigraphic column of three units), scale (ratio $1/50\,000$ and graphic scale bar), north arrow, and a marginal cross-section A--B. The author/institution and the graticule are not mentioned, so they cannot be confirmed.
\item \textbf{(b)} The legend is a stratigraphic column read from bottom to top: the Precambrian granite--gneiss (unit 1) is the oldest, followed by the Cretaceous sandstone (unit 2), and the Quaternary basalt (unit 3) is the youngest. This is consistent with the geology of the area, where young volcanics of the Cameroon Volcanic Line rest on older rocks.
\item \textbf{(c)} Unit 2 is coloured green, which is the international colour for the Cretaceous period; the dotted ornament indicates a sandstone lithology. The colour therefore confirms, at a glance, the age stated in the legend.
\end{enumerate}
\textbf{Key point}: the relative ages come from the \emph{position in the legend} (oldest at the bottom), and the colours provide an independent international check.
\end{exoresolu}

\begin{exercice}[Testing your first reading]
On an unknown map sheet you observe: a title block ``Sheet NB-32-XI, 1/200\,000''; a legend with four units (bottom to top: pink gneiss, blue limestone, green marl, yellow alluvium); dip symbols of $25^{\circ}$ all pointing east; a red dashed line cutting across all units; and a scale bar. (a) List the essential elements you can confirm and any that appear to be missing. (b) Arrange the four units in chronological order. (c) Suggest what the red dashed line might represent and what you would check to confirm it. (d) Using the five-step method, state what you would examine next.
\end{exercice}

\begin{corrige}[Exercise 1]
\begin{enumerate}
\item \textbf{(a)} Confirmed: title, legend, scale (ratio and bar), structural symbols (dips). Not mentioned, hence to be verified: north arrow, grid/graticule, author and date of survey.
\item \textbf{(b)} Oldest to youngest: pink gneiss (Precambrian basement) $\rightarrow$ blue limestone (Jurassic colour) $\rightarrow$ green marl (Cretaceous colour) $\rightarrow$ yellow alluvium (Quaternary colour). The colours agree with the legend order.
\item \textbf{(c)} A red line cutting across all units is most probably a fault (or possibly a dyke); the dashed style suggests it is inferred rather than directly observed. Check the legend for the symbol, look for dip symbols or displaced contacts across it, and examine the marginal cross-section.
\item \textbf{(d)} Having done title, scale and legend (steps 1--3), proceed to step 4 (structures: the uniform $25^{\circ}$ easterly dips suggest homoclinal strata, plus the suspected fault) and step 5 (relief: contour lines and rivers to test the outcrop pattern against topography).
\end{enumerate}
\end{corrige}

\begin{aretenir}[Key points of Section 2]
\begin{itemize}
\item A geological map is a complete document: \cle{title}, \cle{legend}, \cle{scale}, \cle{north arrow}, \cle{grid/graticule}, \cle{author} and \cle{date} are essential elements.
\item The legend is a stratigraphic column: \textbf{oldest unit at the bottom, youngest at the top}; colours follow the international chronostratigraphic code (Quaternary yellow, Cretaceous green, Jurassic blue, basement pink/red).
\item Marginal information (index of adjoining sheets, location map, reliability diagram, cross-sections, explanatory memoir) completes the sheet and indicates its reliability.
\item Quality is judged by the accuracy of contacts, the completeness of structural data, the clarity of symbols and internal consistency.
\item First reading of any map follows five steps: title $\rightarrow$ scale $\rightarrow$ legend $\rightarrow$ structures $\rightarrow$ relief.
\end{itemize}
\end{aretenir}

With the anatomy of a geological map now clear, we can move to the tools that allow us to measure on it: Section 3 examines direction, orientation and scale — the quantitative skeleton beneath the colours.

\coursec{Section 3: Direction, Orientation and Scale on Maps}

In Section 2 we saw that a geological map is a \cle{scaled}, \cle{oriented} representation of the solid geology at the Earth's surface, complete with a key, a scale and a north arrow. Two of those characteristics — \textbf{orientation} and \textbf{scale} — deserve a full treatment of their own, because every practical exercise in map work depends on them. Before you can draw a cross-section, describe a fold or solve a three-point problem, you must be able to say precisely \emph{where} a point lies on the map, \emph{in what direction} one locality is from another, and \emph{how far apart} they really are on the ground. This section gives you those three skills: direction, grid references, and scale.

\courssub{The Three Norths and Magnetic Declination}

Stand on the slopes of Mount Cameroon with a compass in your hand and ask a simple question: \emph{which way is north?} Surprisingly, there are three different correct answers.

\begin{definition}[The three norths]
\begin{itemize}
\item \cle{True north} (geographic north): the direction towards the Earth's geographic North Pole, i.e. along a line of longitude (meridian). It is fixed.
\item \cle{Grid north}: the direction of the vertical grid lines (eastings) printed on the map. Because a flat map cannot perfectly represent a curved Earth, grid north may differ very slightly from true north.
\item \cle{Magnetic north}: the direction in which a compass needle points, towards the Earth's magnetic pole. This pole wanders slowly with time.
\end{itemize}
\end{definition}

The angle between magnetic north and true (or grid) north is called the \cle{magnetic declination}. It varies from place to place and changes slowly year by year, which is why good topographic maps print the declination, its date, and its annual change in the margin.

\begin{popfigure}
\begin{center}
\begin{tikzpicture}[scale=1.1]
% True north
\draw[-{Stealth},very thick,popblue] (0,0) -- (0,4.2) node[above,popblue]{\textbf{TN}};
% Grid north
\draw[-{Stealth},very thick,popgreen] (0,0) -- (0.35,4.15) node[above right,popgreen]{\textbf{GN}};
% Magnetic north
\draw[-{Stealth},very thick,poporange] (0,0) -- (-0.85,3.9) node[above left,poporange]{\textbf{MN}};
% declination arcs
\draw[popmuted,thick] (0,2.6) arc (90:85:2.6) node[midway,right,popmuted,font=\small]{GN--TN};
\draw[popdark,thick] (0,3.1) arc (90:102:3.1) node[midway,above,popdark,font=\small]{declination};
\fill (0,0) circle (2pt) node[below right]{observer};
% compass ticks
\draw[-{Stealth},gray] (0,0) -- (3.9,0) node[right,gray]{E};
\draw[-{Stealth},gray] (0,0) -- (0,-1.2) node[below,gray]{S};
\draw[-{Stealth},gray] (0,0) -- (-1.2,0) node[left,gray]{W};
\end{tikzpicture}
\end{center}
\legende{The three norths. TN = true north, GN = grid north, MN = magnetic north. The angle between MN and TN is the magnetic declination (here MN lies west of TN: a westerly declination).}
\end{popfigure}

\begin{attention}[Declination changes with time]
Never assume the declination printed on an old map is still valid. Magnetic north drifts by a few minutes of arc per year. In examinations, if a declination is given in the question, use \emph{that} value and state whether it is easterly or westerly.
\end{attention}

\courssub{Expressing Direction: Compass Points, Azimuths and Quadrants}

Geologists need to communicate direction unambiguously — for example the direction of dip of a bed, or the trend of a fault trace. Three systems are in use.

\textbf{1. Compass points.} The cardinal points (N, E, S, W), the intercardinal points (NE, SE, SW, NW) and finer subdivisions such as NNE or ENE. This is quick but coarse: each of the 16 classic points covers $22.5^{\circ}$.

\textbf{2. Whole-circle bearings (azimuths).} The direction is measured \emph{clockwise from north}, from $000^{\circ}$ to $360^{\circ}$. North is $000^{\circ}$ (or $360^{\circ}$), east $090^{\circ}$, south $180^{\circ}$, west $270^{\circ}$.

\begin{definition}[Bearing (azimuth)]
The \cle{bearing} (or \cle{azimuth}) of a line is the horizontal angle measured \textbf{clockwise from north} to the line, expressed in degrees from $000^{\circ}$ to $360^{\circ}$, always written with three figures (e.g. $045^{\circ}$, $270^{\circ}$).
\end{definition}

\textbf{3. Quadrant bearings.} The direction is stated as an angle east or west of north or south, e.g. N $45^{\circ}$ E, S $30^{\circ}$ W. This system is common in structural geology for recording strike (e.g. a bed striking N $40^{\circ}$ E).

\begin{exemplebox}[Converting between systems]
A dyke trends N $60^{\circ}$ E. In whole-circle bearing this is simply $060^{\circ}$. A fault trending S $30^{\circ}$ W corresponds to $180^{\circ}+30^{\circ}=210^{\circ}$. Conversely, a bearing of $315^{\circ}$ is NW (equivalently N $45^{\circ}$ W).
\end{exemplebox}

\courssub{Measuring Bearings on a Map; Back Bearings}

\begin{methode}[Measuring a bearing with a protractor]
\begin{enumerate}
\item Draw a light pencil line joining the two points A and B.
\item At the point \emph{from which} the bearing is required (A), draw a small north line parallel to the grid eastings.
\item Place the protractor centre on A with its $0^{\circ}$ on the north line.
\item Read the angle \textbf{clockwise} from north to the line AB. That is the bearing of B from A.
\end{enumerate}
\end{methode}

The \cle{back bearing} is the bearing of A from B — the reverse direction. Since the two directions differ by exactly $180^{\circ}$:
\[
\text{back bearing} =
\begin{cases}
\text{bearing} + 180^{\circ} & \text{if bearing} < 180^{\circ},\\[2pt]
\text{bearing} - 180^{\circ} & \text{if bearing} \geq 180^{\circ}.
\end{cases}
\]

\begin{exemplebox}[Back bearing]
The bearing of a trigonometrical station from a spring is $065^{\circ}$. The back bearing (spring from the station) is $065^{\circ}+180^{\circ}=245^{\circ}$.
\end{exemplebox}

\courssub{Grid References: Locating a Point Precisely}

Topographic and geological maps carry a printed kilometre grid. The \textbf{vertical} lines are the \cle{eastings} (their numbers increase eastward) and the \textbf{horizontal} lines are the \cle{northings} (increasing northward). A grid reference locates a point within this grid.

\begin{methode}[Giving a six-figure grid reference]
\begin{enumerate}
\item Read the \textbf{easting first}: identify the grid line to the \emph{left} (west) of the point, e.g. 34.
\item Estimate how many tenths of a grid square the point lies to the east of that line, e.g. 6 tenths. This gives the first three figures: 346.
\item Then read the \textbf{northing}: identify the grid line \emph{below} (south of) the point, e.g. 57, and estimate tenths northward, e.g. 2. This gives 572.
\item Write the six-figure reference as \textbf{346572} — eastings before northings.
\end{enumerate}
\end{methode}

A \cle{four-figure grid reference} (e.g. 3457) identifies only the whole $1\ \mathrm{km}$ square whose south-west corner is at easting 34, northing 57. A \cle{six-figure grid reference} (e.g. 346572) locates a point to within $100\ \mathrm{m}$.

\begin{popfigure}
\begin{center}
\begin{tikzpicture}[scale=1.6]
\draw[step=1,popline,thin] (0,0) grid (3,3);
\draw[very thick,popblue] (0,0) rectangle (1,1);
% labels
\node[below] at (0,0) {\small 34};
\node[below] at (1,0) {\small 35};
\node[below] at (2,0) {\small 36};
\node[left] at (0,0) {\small 57};
\node[left] at (0,1) {\small 58};
\node[left] at (0,2) {\small 59};
% point P
\fill[poporange] (0.6,0.2) circle (2.2pt);
\node[right,poporange] at (0.66,0.2) {\textbf{P}};
% tenths guides
\draw[dashed,popmuted] (0.6,0) -- (0.6,0.2);
\draw[dashed,popmuted] (0,0.2) -- (0.6,0.2);
\node[below,popmuted,font=\scriptsize] at (0.3,0) {6 tenths E};
\node[left,popmuted,font=\scriptsize] at (0,0.1) {};
\node[popdark,font=\small] at (1.5,-0.55) {P: four-figure \textbf{3457}; six-figure \textbf{346572}};
\end{tikzpicture}
\end{center}
\legende{Grid references. The shaded square has the four-figure reference 3457. Point P lies 6 tenths east and 2 tenths north within it, giving the six-figure reference 346572.}
\end{popfigure}

\begin{attention}[Eastings before northings]
The most frequent error is reversing the order. Remember: \emph{``along the corridor (eastings), then up the stairs (northings)''}. Swapping them — writing 572346 instead of 346572 — places the point in a completely different part of the map and loses all the marks for that question.
\end{attention}

\courssub{Map Scale: Relating Map Distance to Ground Distance}

A map is a reduced image of the ground, and the \cle{scale} states by how much it has been reduced.

\begin{definition}[Representative fraction (RF)]
The \cle{representative fraction} is the ratio of a distance on the map to the corresponding distance on the ground, both in the \textbf{same units}:
\[
\mathrm{RF}=\frac{\text{map distance}}{\text{ground distance}}.
\]
An RF of $1:25\,000$ means that 1 unit on the map represents $25\,000$ of the same units on the ground: $1\ \mathrm{cm}$ on the map $=25\,000\ \mathrm{cm}=250\ \mathrm{m}$ on the ground.
\end{definition}

Scales are expressed in three ways:
\begin{itemize}
\item \textbf{Representative fraction (ratio) scale}: e.g. $1:25\,000$, $1:50\,000$, $1:200\,000$. Unit-free and universal.
\item \textbf{Statement (verbal) scale}: e.g. ``$1\ \mathrm{cm}$ to $250\ \mathrm{m}$'' or ``$4\ \mathrm{cm}$ to $1\ \mathrm{km}$''.
\item \textbf{Graphical (linear) scale}: a drawn bar divided into ground distances. Its great advantage is that it enlarges or shrinks together with the map if the map is photocopied.
\end{itemize}

\begin{popfigure}
\begin{center}
\begin{tikzpicture}[scale=1]
% 1 cm = 250 m ; total 1 km = 4 cm, with first cm subdivided into 100 m
\fill[popdark] (0,0) rectangle (0.25,0.35);
\fill[white] (0.25,0) rectangle (0.5,0.35);
\fill[popdark] (0.5,0) rectangle (0.75,0.35);
\fill[white] (0.75,0) rectangle (1,0.35);
\draw[thick] (0,0) rectangle (4,0.35);
\foreach \x in {0.25,0.5,0.75} \draw[thick] (\x,0) -- (\x,0.35);
\foreach \x in {1,2,3} \draw[thick] (\x,0) -- (\x,0.35);
\node[below,font=\small] at (0,0) {0};
\node[below,font=\small] at (0.5,0) {100};
\node[below,font=\small] at (1,0) {200 m};
\node[below,font=\small] at (2,0) {400};
\node[below,font=\small] at (3,0) {600};
\node[below,font=\small] at (4,0) {800 m};
\node[above,font=\small] at (2,0.35) {Graphical scale for $1:25\,000$ ($1\ \mathrm{cm}=250\ \mathrm{m}$)};
\end{tikzpicture}
\end{center}
\legende{A graphical (linear) scale for a $1:25\,000$ map. The left-hand division is subdivided into hundreds of metres for precise readings.}
\end{popfigure}

\courssub{Scale Conversions: Worked Examples}

\begin{methode}[Converting map distance to ground distance]
\[
\text{ground distance} = \text{map distance} \times \text{scale denominator}
\]
then convert the result into convenient units ($100\,000\ \mathrm{cm}=1\ \mathrm{km}$).
\end{methode}

\begin{exoresolu}[Map distance to ground distance]
On a $1:50\,000$ geological map, the outcrop of a dolerite dyke measures $6.4\ \mathrm{cm}$ along a straight valley. What is the true length of the dyke on the ground?
\tcblower
\textbf{Solution.}
\begin{align*}
\text{ground distance} &= 6.4\ \mathrm{cm} \times 50\,000 = 320\,000\ \mathrm{cm}\\
&= \frac{320\,000}{100\,000}\ \mathrm{km} = 3.2\ \mathrm{km}.
\end{align*}
The dyke is $3.2\ \mathrm{km}$ long on the ground.
\end{exoresolu}

\begin{exoresolu}[Ground distance to map distance]
Two villages are $7.5\ \mathrm{km}$ apart. How far apart will they appear on a $1:25\,000$ map?
\tcblower
\textbf{Solution.} Convert the ground distance to centimetres, then divide by the denominator:
\begin{align*}
7.5\ \mathrm{km} &= 7.5 \times 100\,000 = 750\,000\ \mathrm{cm},\\
\text{map distance} &= \frac{750\,000}{25\,000} = 30\ \mathrm{cm}.
\end{align*}
The villages will be $30\ \mathrm{cm}$ apart on the map.
\end{exoresolu}

\begin{attention}[Exam tip: always carry the units]
The RF is unit-free only because \emph{both} distances are in the same unit. In your working, write every step with its units ($\mathrm{cm}$, $\mathrm{m}$, $\mathrm{km}$) and show the conversion $1\ \mathrm{km}=100\,000\ \mathrm{cm}$. Examiners award method marks for correct unit handling, and most lost marks in scale questions come from missing zeros.
\end{attention}

\courssub{Consequences of Scale: Detail, Enlargement and Reduction}

The scale controls how much detail a map can show. A \cle{large-scale} map (small denominator, e.g. $1:10\,000$) shows a small area in great detail — individual lava flows on Mount Cameroon, small dykes, minor faults. A \cle{small-scale} map (large denominator, e.g. $1:1\,000\,000$) shows a vast area — the whole Cameroon Volcanic Line — but only the broadest geological units.

\begin{center}
\begin{tabularx}{\linewidth}{|l|X|X|}
\hline
\rowcolor{popblueL} \textbf{Scale} & \textbf{$1\ \mathrm{cm}$ represents} & \textbf{Typical use} \\
\hline
$1:10\,000$ & $100\ \mathrm{m}$ & Detailed mine or site maps \\
\hline
$1:25\,000$ & $250\ \mathrm{m}$ & Detailed field mapping, GCE exercises \\
\hline
$1:50\,000$ & $500\ \mathrm{m}$ & Standard topographic/geological sheets \\
\hline
$1:200\,000$ & $2\ \mathrm{km}$ & Regional reconnaissance maps \\
\hline
$1:1\,000\,000$ & $10\ \mathrm{km}$ & National or continental compilations \\
\hline
\end{tabularx}
\end{center}

\textbf{Enlarging and reducing maps.} If a map is enlarged by a factor $k$, every distance is multiplied by $k$, so the new RF denominator is the old one \emph{divided} by $k$. Enlarging a $1:50\,000$ map by a factor of 2 produces a $1:25\,000$ map. Note that a photocopied RF or statement scale becomes \emph{wrong} after enlargement — only the \textbf{graphical scale} remains valid, which is why it is indispensable on examination maps.

\begin{exemplebox}[Effect of enlargement]
A $1:100\,000$ map is enlarged three times for a classroom exercise. New scale $=100\,000/3 \approx 1:33\,333$. A contact that measured $2\ \mathrm{cm}$ on the original now measures $6\ \mathrm{cm}$, but still represents the same $2\ \mathrm{km}$ on the ground.
\end{exemplebox}

\courssub{Gradient and Slope from Contour Spacing}

Scale also lets us measure \emph{vertical} relationships. The \cle{gradient} of a slope is the ratio of vertical rise to horizontal distance, usually expressed as a fraction ($1$ in $n$), a percentage, or an angle:
\[
\text{gradient}=\frac{\text{vertical interval (VI)}}{\text{horizontal equivalent (HE)}}, \qquad \tan\theta=\frac{\text{VI}}{\text{HE}}.
\]

\begin{exoresolu}[Calculating a gradient from a map]
On a $1:25\,000$ map, two successive $20\ \mathrm{m}$ contours crossing a spur are $8\ \mathrm{mm}$ apart. Calculate the gradient and the angle of slope.
\tcblower
\textbf{Solution.}
\begin{align*}
\text{HE} &= 0.8\ \mathrm{cm}\times 25\,000 = 20\,000\ \mathrm{cm}=200\ \mathrm{m},\\
\text{gradient} &= \frac{\text{VI}}{\text{HE}}=\frac{20}{200}=\frac{1}{10}=10\ \%,\\
\theta &= \tan^{-1}(0.1)\approx 5.7^{\circ}.
\end{align*}
The slope has a gradient of 1 in 10 (10 \%) and an angle of about $5.7^{\circ}$.
\end{exoresolu}

This skill is the bridge to the next part of the chapter: when you construct geological cross-sections (Sections 7 and 8), you will choose a vertical scale, compute true dips from apparent geometry, and lay off slopes exactly as we have done here. Mastering direction, grid references and scale now will make every later exercise faster and more accurate.

\begin{aretenir}[Key points of Section 3]
\begin{itemize}
\item Three norths exist: true, grid and magnetic; the angle between magnetic and true north is the \cle{magnetic declination}.
\item A \cle{bearing} is measured clockwise from north, $000^{\circ}$--$360^{\circ}$; back bearing $=$ bearing $\pm 180^{\circ}$.
\item Grid references: \cle{eastings before northings}; four figures locate a $1\ \mathrm{km}$ square, six figures locate a point to $100\ \mathrm{m}$.
\item $\mathrm{RF}=\dfrac{\text{map distance}}{\text{ground distance}}$; ground distance $=$ map distance $\times$ denominator.
\item Large scale $=$ small denominator $=$ much detail; only the \cle{graphical scale} survives enlargement.
\item Gradient $=$ VI/HE; slope angle $\theta=\tan^{-1}(\text{VI}/\text{HE})$.
\end{itemize}
\end{aretenir}

\begin{exercice}[Practice]
\begin{enumerate}
\item A fault trace bears $138^{\circ}$ from a village. What is the back bearing? Express $138^{\circ}$ as a quadrant bearing.
\item A spring lies 7 tenths east and 4 tenths north within grid square 4281. Give its six-figure grid reference.
\item On a $1:50\,000$ map, the outcrop width of a sandstone bed measures $1.6\ \mathrm{cm}$. Calculate the true outcrop width in metres.
\item A $1:200\,000$ map is enlarged by a factor of 4. What is the scale of the enlarged map?
\item Contours of $300\ \mathrm{m}$ and $340\ \mathrm{m}$ are $5\ \mathrm{mm}$ apart on a $1:25\,000$ map. Find the gradient and slope angle.
\end{enumerate}
\end{exercice}

\begin{corrige}[Exercise 1]
\begin{enumerate}
\item Back bearing $=138^{\circ}+180^{\circ}=318^{\circ}$. Quadrant form: $138^{\circ}-90^{\circ}=48^{\circ}$ east of south, i.e. S $42^{\circ}$ E (since $180^{\circ}-138^{\circ}=42^{\circ}$).
\item Easting 42 + 7 tenths $\rightarrow 427$; northing 81 + 4 tenths $\rightarrow 814$. Reference: \textbf{427814}.
\item $1.6\ \mathrm{cm}\times 50\,000=80\,000\ \mathrm{cm}=800\ \mathrm{m}$.
\item $200\,000/4=50\,000$: the enlarged map is at $1:50\,000$.
\item HE $=0.5\ \mathrm{cm}\times25\,000=12\,500\ \mathrm{cm}=125\ \mathrm{m}$; gradient $=40/125=1$ in $3.125$ ($32\ \%$); $\theta=\tan^{-1}(0.32)\approx 17.7^{\circ}$.
\end{enumerate}
\end{corrige}

\coursec{Section 4: Conventional Geological Symbols on Maps}

In Section 3 we learnt how to locate ourselves on a map: the north arrow, grid references, bearings and the scale bar tell us \emph{where} we are and \emph{how big} things are. But a geological map must do much more than locate points. It must compress a three-dimensional body of rock --- beds dipping underground, folds plunging beneath the surface, faults cutting across country --- onto a flat sheet of paper. This is achieved through a vocabulary of \cle{conventional symbols}: small, standardised signs whose meaning is fixed by international and national (Cameroon GCE Board) convention. A geologist in Bamenda, in Buea or in Birmingham reads the same symbol in the same way. Mastering this vocabulary is the key that unlocks every geological map you will meet in the GCE Advanced Level examination.

\courssub{Why Symbols? Representing 3D Information on a 2D Sheet}

A geological map is a \emph{plan view}: it shows what a vertical observer, looking straight down, would see at the Earth's surface (with soil and vegetation mentally removed). Yet the geologist's real interest is underground. The solution adopted worldwide is elegant:

\begin{itemize}
\item \textbf{Lines} (contacts, faults, fold axes) show \emph{where} geological boundaries meet the surface.
\item \textbf{Symbols} (strike-and-dip signs, plunge arrows, downthrow ticks) show \emph{how} the structures are oriented in three dimensions at those places.
\item \textbf{Colours and ornament} (studied in Section 2) show the \emph{age and nature} of each rock unit.
\end{itemize}

Thus a single strike-and-dip symbol placed on a contact tells the reader both the horizontal direction of the bed and the angle and direction at which it sinks into the ground. Symbols are the bridge between the flat map and the solid Earth.

\courssub{Boundary Lines: Contacts Between Rock Units}

\begin{definition}[Geological contact (boundary)]
A \cle{geological contact} (or boundary) is the surface that separates two distinct rock units --- two formations of different age or lithology, or a rock unit and an unconsolidated deposit. On a map, the contact is drawn as the \emph{line along which that surface intersects the ground surface}.
\end{definition}

Because contacts are rarely visible everywhere (they may be hidden by soil, forest, rivers or laterite, as is common in the Cameroonian humid zone), map-makers distinguish contacts according to how well they are known:

\begin{center}
\begin{tabularx}{\linewidth}{|l|X|X|}
\hline
\rowcolor{popblueL} \textbf{Line style} & \textbf{Meaning} & \textbf{When used} \\
\hline
Solid continuous line & \cle{Observed contact}: seen and located accurately in the field & Contact exposed in outcrops, road cuts, stream beds \\
\hline
Dashed line & \cle{Inferred (approximate) contact}: position deduced, not directly seen & Contact hidden by soil, vegetation or water \\
\hline
Dotted line & \cle{Concealed contact}: known to exist but completely covered & Contact buried beneath alluvium or a lava flow \\
\hline
Solid line with teeth/ticks & Fault or unconformity (ornament shows type) & Structural or erosional boundaries \\
\hline
\end{tabularx}
\end{center}

Contacts are also classified by their geological nature:

\begin{itemize}
\item \textbf{Conformable contact}: separates beds deposited in continuous succession, parallel to one another. On the map the contact runs roughly parallel to neighbouring contacts.
\item \textbf{Unconformable contact}: an eroded, ancient land surface on which younger beds rest. On the map the younger unit \emph{cuts across} (truncates) the contacts of the older units below --- a key recognition criterion.
\item \textbf{Fault contact}: the boundary is a fracture along which movement has occurred; it is drawn as a fault line (see below), not as an ordinary contact.
\item \textbf{Intrusive contact}: an igneous body cuts across pre-existing structures; the contact is typically irregular and discordant.
\end{itemize}

\begin{attention}[Dashed does not mean ``unimportant'']
A \textbf{dashed contact is an inferred contact}, not an observed one. In examination descriptions, writing ``the contact between units A and B is a fault'' when the map shows an ordinary dashed line is a serious error. Always state the line style's meaning: ``the contact is inferred (dashed), probably because it is concealed by alluvium along the valley''.
\end{attention}

\courssub{Attitude Symbols: Strike and Dip}

The most important symbol on any geological map is the \cle{strike-and-dip symbol}, which records the attitude (orientation in space) of a bedding plane measured at an outcrop.

\begin{definition}[Strike and dip]
The \cle{strike} of a bed is the direction of the horizontal line formed by the intersection of the bedding plane with a horizontal plane. The \cle{dip} is the angle, measured in a vertical plane perpendicular to the strike, between the bedding plane and the horizontal; the \cle{dip direction} is the compass direction in which the bed slopes downward.
\end{definition}

\begin{methode}[Reading a strike-and-dip symbol]
\begin{enumerate}
\item Identify the \textbf{long line}: it represents the \emph{strike} of the bed. Its orientation on the map is the true strike direction (read it with a protractor against the north arrow).
\item Identify the \textbf{short tick} attached at the midpoint of the long line: it points in the \emph{dip direction}.
\item Read the \textbf{number} written beside the tick: it is the \emph{angle of dip} in degrees (e.g.\ $30$ means the bed dips at $30^{\circ}$).
\item Combine: ``the bed strikes N--S and dips $30^{\circ}$ to the east''.
\end{enumerate}
\end{methode}

\begin{attention}[The tick points DOWN-dip]
The short tick of the strike-and-dip symbol always points \textbf{down-dip} --- in the direction in which water would run down the bedding plane --- \emph{never} up-dip. Reversing the dip direction is the single most common mistake in GCE map work: it inverts the whole geometry of the cross-section and ruins the geological history. Check the tick before writing anything.
\end{attention}

Special attitudes have special symbols:

\begin{itemize}
\item \textbf{Horizontal beds}: a cross ($+$) inside a small circle; dip $= 0^{\circ}$, strike undefined.
\item \textbf{Vertical beds}: a single long line (strike) crossed by a short bar on \emph{both} sides (or a tick on each side); dip $= 90^{\circ}$, no dip direction.
\item \textbf{Overturned beds}: the normal strike-and-dip symbol, but the tick is bent into a hook (a short ``U'' or arrow curving back towards the strike line), indicating that the bed has been tilted past vertical so that it now dips in the direction opposite to its original younging.
\end{itemize}

\begin{popfigure}
\begin{center}
\begin{tikzpicture}[scale=1.0, line cap=round]
% (a) dipping bed
\draw[very thick, popink] (0,-0.7) -- (0,0.7);
\draw[very thick, poporange] (0,0) -- (0.45,0);
\node[popink, anchor=west] at (0.55,0.15) {\small $30$};
\node[align=center, text width=3.2cm] at (0.3,-1.3) {\small (a) Dipping bed:\\ strike N--S, dip $30^{\circ}$ E};
% (b) horizontal bed
\begin{scope}[xshift=4.2cm]
\draw[very thick, popink] (0,0) circle (0.28);
\draw[very thick, popink] (-0.28,0) -- (0.28,0);
\draw[very thick, popink] (0,-0.28) -- (0,0.28);
\node[align=center, text width=3.2cm] at (0,-1.3) {\small (b) Horizontal bed\\ (dip $0^{\circ}$)};
\end{scope}
% (c) vertical bed
\begin{scope}[xshift=8.2cm]
\draw[very thick, popink] (0,-0.7) -- (0,0.7);
\draw[very thick, poporange] (-0.35,0) -- (0.35,0);
\node[align=center, text width=3.2cm] at (0,-1.3) {\small (c) Vertical bed\\ (dip $90^{\circ}$)};
\end{scope}
% (d) overturned bed
\begin{scope}[xshift=12.2cm]
\draw[very thick, popink] (0,-0.7) -- (0,0.7);
\draw[very thick, poporange] (0,0) -- (0.45,0);
\draw[very thick, poporange, ->] (0.45,0) arc (0:-95:0.22);
\node[popink, anchor=west] at (0.55,0.35) {\small $60$};
\node[align=center, text width=3.4cm] at (0.3,-1.3) {\small (d) Overturned bed:\\ hooked tick, dip $60^{\circ}$};
\end{scope}
\end{tikzpicture}
\end{center}
\legende{Conventional attitude symbols: (a) normal dipping bed, (b) horizontal bed, (c) vertical bed, (d) overturned bed. The long line is always the strike; the tick always points down-dip.}
\end{popfigure}

\courssub{Fold Symbols}

Folds are represented by their \cle{axial trace} --- the line along which the axial plane of the fold meets the ground surface --- ornamented to show the fold type:

\begin{itemize}
\item \textbf{Anticline}: a line along the fold axis with two short arrows pointing \emph{away} from the axis (limbs dip away from each other). Oldest rocks lie in the core.
\item \textbf{Syncline}: a line along the axis with two short arrows pointing \emph{towards} the axis (limbs dip towards each other). Youngest rocks lie in the core.
\item \textbf{Plunging folds}: an arrow placed \emph{on} the axis line points in the direction of \cle{plunge} (the direction in which the fold axis sinks below the horizontal); a number may give the plunge angle. Plunging anticlines produce V-shaped or U-shaped outcrop patterns closing in the plunge direction; plunging synclines close in the opposite direction.
\item \textbf{Overturned fold}: both limb arrows point in the \emph{same} general direction, showing that both limbs dip the same way (one limb is inverted).
\end{itemize}

\courssub{Fault Symbols}

A fault is drawn as a line (usually thicker or in a distinctive colour, often red or black) with ornaments describing the movement:

\begin{itemize}
\item \textbf{Fault trace}: solid where observed, dashed where inferred.
\item \textbf{Dip of the fault plane}: a short tick on the line points in the dip direction of the fault plane, with a number giving the dip angle (read exactly like a bedding dip).
\item \textbf{Downthrow side}: a short bar or tick (or the letter \textbf{D}, with \textbf{U} on the upthrow side) marks the side that has moved \emph{down}. On many maps a small arrow perpendicular to the fault shows the relative vertical movement.
\item \textbf{Strike-slip (wrench) faults}: two half-arrows parallel to the fault, one on each side, pointing in opposite directions to show horizontal relative movement (dextral or sinistral).
\item \textbf{Thrust faults}: triangles (teeth) on the line point towards the \emph{upper} (overriding) block.
\end{itemize}

\begin{popfigure}
\begin{center}
\begin{tikzpicture}[scale=1.0, line cap=round]
% anticline
\draw[very thick, poppurple] (0,-0.8) -- (0,0.8);
\draw[thick, poppurple, ->] (0,0.25) -- (-0.45,0.25);
\draw[thick, poppurple, ->] (0,-0.25) -- (0.45,-0.25);
\node[align=center, text width=2.6cm] at (0,-1.35) {\small Anticline axis\\ (arrows apart)};
% syncline
\begin{scope}[xshift=3.4cm]
\draw[very thick, poppurple] (0,-0.8) -- (0,0.8);
\draw[thick, poppurple, ->] (-0.45,0.25) -- (0,0.25);
\draw[thick, poppurple, ->] (0.45,-0.25) -- (0,-0.25);
\node[align=center, text width=2.6cm] at (0,-1.35) {\small Syncline axis\\ (arrows together)};
\end{scope}
% plunging anticline
\begin{scope}[xshift=6.8cm]
\draw[very thick, poppurple] (0,-0.8) -- (0,0.8);
\draw[thick, poppurple, ->] (0,0.25) -- (-0.45,0.25);
\draw[thick, poppurple, ->] (0,-0.25) -- (0.45,-0.25);
\draw[very thick, poporange, ->] (0,0.1) -- (0,0.75);
\node[align=center, text width=2.8cm] at (0,-1.35) {\small Plunging anticline (plunge arrow)};
\end{scope}
% dip-slip fault
\begin{scope}[xshift=10.2cm]
\draw[very thick, popdark] (0,-0.8) -- (0,0.8);
\draw[thick, popdark] (0,0.1) -- (0.4,0.1);
\node[popdark, anchor=west] at (0.42,0.35) {\small D};
\node[popdark, anchor=east] at (-0.08,0.35) {\small U};
\node[align=center, text width=2.8cm] at (0.1,-1.35) {\small Dip-slip fault (D = downthrow)};
\end{scope}
% strike-slip fault
\begin{scope}[xshift=13.4cm]
\draw[very thick, popdark] (0,-0.8) -- (0,0.8);
\draw[thick, poporange, ->] (-0.3,-0.5) -- (-0.3,0.5);
\draw[thick, poporange, ->] (0.3,0.5) -- (0.3,-0.5);
\node[align=center, text width=2.6cm] at (0,-1.35) {\small Strike-slip fault (half-arrows)};
\end{scope}
\end{tikzpicture}
\end{center}
\legende{Conventional fold and fault symbols: anticline and syncline axes, plunging anticline, dip-slip fault with downthrow tick, and strike-slip fault with opposed half-arrows.}
\end{popfigure}

\courssub{Other Common Symbols}

Geological maps carry a range of additional symbols recording points of economic or scientific interest:

\begin{center}
\begin{tabularx}{\linewidth}{|l|X|}
\hline
\rowcolor{popblueL} \textbf{Symbol} & \textbf{Meaning} \\
\hline
Small circle with central dot, labelled with depth & \cle{Borehole}: the number gives the depth drilled or the elevation of the bed reached \\
\hline
Small filled or open circle labelled ``well'' & Water well \\
\hline
Small circle with a short wavy tail & \cle{Spring}: natural emergence of groundwater (common on the flanks of Mount Cameroon) \\
\hline
Crossed hammers, or a crossed pick-and-shovel sign & Mine or quarry (active or abandoned) \\
\hline
Small crossed hammers with a fossil sketch or ``F'' & \cle{Fossil locality} \\
\hline
Star or asterisk, often labelled (Au, Fe, etc.) & Mineral occurrence / show of ore \\
\hline
Thick straight or sinuous band, often dark & \cle{Dyke}: tabular igneous intrusion cutting across bedding \\
\hline
Thin irregular line, sometimes with mineral label & Vein (e.g.\ quartz vein) \\
\hline
\end{tabularx}
\end{center}

\begin{savaistu}[Symbols and the Cameroon Volcanic Line]
On geological maps of the Cameroon Volcanic Line, dykes appear as dense swarms of short, thick, parallel lines trending roughly N$30^{\circ}$E around Mount Cameroon and the Bakossi highlands, while small circles with wavy tails mark the abundant springs fed by fractured basaltic aquifers. Reading these symbols lets a geologist predict groundwater occurrence --- a direct application of map work to water supply in towns such as Buea.
\end{savaistu}

\courssub{Reading a Symbol-Rich Map Fragment}

Let us now combine everything on a small map extract. Read it carefully before looking at the worked example: three sedimentary units (A, B, C) are separated by two contacts; a fault cuts unit A; a fold axis runs through unit B; three strike-and-dip symbols and a borehole complete the picture.

\begin{popfigure}
\begin{center}
\begin{tikzpicture}[scale=1.0, line cap=round]
% frame 5 cm square
\draw[thick, popink] (0,0) rectangle (5,5);
% north arrow
\draw[thick, popink, ->] (4.6,4.0) -- (4.6,4.7);
\node[popink] at (4.6,3.85) {\small N};
% unit colours
\fill[popblueL] (0,0) -- (0,5) -- (1.6,5) -- (2.2,0) -- cycle;
\fill[popgreenL] (2.2,0) -- (1.6,5) -- (3.4,5) -- (3.0,0) -- cycle;
\fill[popgold] (3.0,0) -- (3.4,5) -- (5,5) -- (5,0) -- cycle;
% contacts
\draw[very thick, popink] (2.2,0) -- (1.6,5);
\draw[very thick, popink, dashed] (3.0,0) -- (3.4,5);
% fault
\draw[very thick, popdark] (0.4,0) -- (1.2,5);
\draw[thick, popdark] (0.95,3.4) -- (1.35,3.4);
\node[popdark] at (1.5,3.55) {\small D};
\node[popdark] at (0.55,3.55) {\small U};
% fold axis (syncline) in central unit
\draw[very thick, poppurple] (2.6,0.2) -- (2.5,4.8);
\draw[thick, poppurple, ->] (2.15,2.8) -- (2.55,2.8);
\draw[thick, poppurple, ->] (2.95,2.2) -- (2.55,2.2);
% dip symbols: all strike N-S (vertical lines), ticks point towards the axis
% symbol in unit A: dips east 25
\draw[very thick, popink] (1.9,0.9) -- (1.9,1.5);
\draw[very thick, poporange] (1.9,1.2) -- (2.3,1.2);
\node[popink, anchor=west] at (2.32,1.05) {\small $25$};
% symbol in unit B (west of axis): dips east 40, towards axis
\draw[very thick, popink] (2.25,3.9) -- (2.25,4.5);
\draw[very thick, poporange] (2.25,4.2) -- (2.65,4.2);
\node[popink, anchor=south] at (2.25,4.55) {\small $40$};
% symbol in unit C: dips west 30, towards axis
\draw[very thick, popink] (4.2,2.1) -- (4.2,2.7);
\draw[very thick, poporange] (4.2,2.4) -- (3.8,2.4);
\node[popink, anchor=west] at (4.28,2.55) {\small $30$};
% borehole
\draw[popink] (4.3,4.3) circle (0.12);
\fill[popink] (4.3,4.3) circle (0.03);
\node[popink, anchor=west] at (4.42,4.3) {\small BH1};
% labels
\node at (0.9,0.35) {\small A};
\node at (2.55,0.35) {\small B};
\node at (4.1,0.35) {\small C};
\end{tikzpicture}
\end{center}
\legende{A symbol-rich map fragment ($5\ \mathrm{cm} \times 5\ \mathrm{cm}$): three units (A, B, C), one observed contact (solid), one inferred contact (dashed), a dip-slip fault in unit A (downthrow to the east), a synclinal axis trending N--S through unit B, three strike-and-dip symbols (all striking N--S) and a borehole (BH1).}
\end{popfigure}

\begin{exoresolu}[Decoding the map fragment]
Study the map fragment above. (a) State the nature of the contact between units A and B, and between B and C. (b) Describe the structure affecting units B and C, using the dip symbols as evidence. (c) Give the attitude of the beds at each of the three dip symbols. (d) Describe the fault.
\tcblower
\textbf{Solution.}
\begin{enumerate}
\item \textbf{Contact A/B}: drawn as a solid line, so it is an \emph{observed} contact, located accurately in the field. \textbf{Contact B/C}: drawn dashed, so it is an \emph{inferred} contact --- its exact position was deduced, probably because it is concealed by superficial deposits.
\item \textbf{Structure in B and C}: the axial trace in unit B carries two arrows pointing \emph{towards} the axis, the symbol for a \emph{syncline} trending roughly N--S. The dip symbols confirm this: the symbol in unit B (west of the axis) dips $40^{\circ}$ to the east, i.e.\ towards the axis, and the symbol in unit C (east of the axis) dips $30^{\circ}$ to the west, i.e.\ also towards the axis. Limbs dipping towards each other define a syncline; the youngest rocks are preserved along the axial zone.
\item \textbf{Attitudes}: symbol in unit A: strike N--S, dip $25^{\circ}$ to the east. Symbol in unit B: strike N--S, dip $40^{\circ}$ to the east. Symbol in unit C: strike N--S, dip $30^{\circ}$ to the \emph{west} (the tick points west --- always read the tick, never the position of the number).
\item \textbf{Fault}: a NNE--SSW dip-slip fault cutting unit A; the tick marked D on the eastern side shows that the \emph{eastern block is downthrown} relative to the western block (marked U); the fault plane dips eastwards.
\end{enumerate}
\end{exoresolu}

\begin{attention}[Number position is not dip direction]
The dip \emph{number} is written wherever there is space on the map; only the \emph{tick} gives the dip direction. In the fragment above, the symbol in unit C dips west even though the number ``30'' is written to the east of the symbol. Similarly, the symbol in unit B has its number written above the strike line, yet the bed dips east.
\end{attention}

\begin{aretenir}[Exam tip: learn the symbol table by heart]
Direct recall questions --- ``draw the symbol for an overturned bed'', ``what does this symbol mean?'' --- are frequent in the GCE Advanced Level paper and cost candidates easy marks. Memorise: the four attitude symbols, anticline/syncline axes, plunge arrow, downthrow tick, strike-slip half-arrows, thrust teeth, and the solid/dashed/dotted contact convention. Practise by sketching the full table from memory, then checking against your notes.
\end{aretenir}

\begin{exercice}[Symbol drill]
\begin{enumerate}
\item Draw the conventional symbols for: (a) a bed striking N$45^{\circ}$E and dipping $35^{\circ}$ to the south-east; (b) a vertical bed striking E--W; (c) a plunging anticline plunging north; (d) a thrust fault with the upper block to the west.
\item On the map fragment above, the beds of unit A dip $25^{\circ}$ to the east. A borehole is to be drilled vertically from a point in unit A, close to the A/B contact. Will it intersect the A/B contact at depth, or will it remain in unit A? Justify your answer using the dip direction.
\item Explain why a dashed contact between two lava flows on the flank of Mount Cameroon is geologically reasonable.
\end{enumerate}
\end{exercice}

\begin{corrige}[Exercise 1]
\begin{enumerate}
\item (a) Long line oriented NE--SW, tick pointing SE, number $35$ beside it. (b) Single E--W line with a short bar on both sides. (c) Axis line with limb arrows pointing away from the axis and an arrow on the axis pointing north. (d) Fault line with triangles on the western side.
\item The beds of unit A dip $25^{\circ}$ to the east, so the A/B contact surface also slopes down towards the east, i.e.\ \emph{away} from unit A and beneath unit B. A vertical borehole drilled in unit A close to the contact will therefore remain in unit A and will \emph{not} intersect the contact at depth; to cut the contact, the hole would have to be drilled in unit B, or inclined eastwards from unit A.
\item Successive lava flows weather rapidly and are quickly mantled by ash, soil and dense vegetation in the humid climate; the contact is therefore rarely exposed and must be inferred between outcrops --- hence the dashed line.
\end{enumerate}
\end{corrige}

With the symbol vocabulary mastered, we are ready for Section 5, where the geometrical properties of strata --- dip, strike, apparent dip and thickness --- are treated quantitatively, giving mathematical substance to the symbols you can now read fluently.

\coursec{Section 5: Geometrical Properties of Strata}

In Section 4 we learned to \emph{read} the conventional symbols printed on a geological map: the little ``T'' of dip and strike, the colours and ornament of each formation, the fault traces and the fold axes. But a symbol is only a compressed message. The dip-and-strike symbol, for instance, tells us that a bedding plane is \emph{inclined} --- yet what exactly is inclined, in which direction, and by how much? To answer these questions, and later to construct cross-sections and compute thicknesses, we must now study the \cle{geometry} of strata themselves: how a planar bed is oriented in space, how thick it really is, and how its outcrop pattern betrays its attitude. This section is the mathematical heart of map work; master it and the rest of the chapter becomes easy.

\courssub{Bedding Planes and the Orientation of a Plane in Space}

\begin{definition}[Bedding plane]
A \cle{bedding plane} is the surface that separates two successive beds (strata) deposited one above the other. It represents a pause or a change in sedimentation. Except in strongly deformed rocks, bedding planes are treated as \emph{planes} in the geometrical sense.
\end{definition}

Imagine a hard, flat slab of rock --- say a sandstone bed exposed in a quarry face near Bafoussam. To describe exactly how that slab lies in space, a geologist needs only \textbf{two numbers}: the direction of a horizontal line drawn on the slab, and the angle by which the slab slopes down from the horizontal. These two numbers are the \cle{strike} and the \cle{dip}.

\begin{definition}[Strike and dip]
\begin{itemize}
\item The \cle{strike} of a bedding plane is the direction (bearing) of any \textbf{horizontal line} drawn on the plane --- equivalently, the line of intersection of the bedding plane with a horizontal plane. It is expressed as a compass bearing, e.g.\ N 045$^{\circ}$ E (or simply 045$^{\circ}$).
\item The \cle{true dip} (or simply \cle{dip}) is the \textbf{maximum angle of inclination} of the plane below the horizontal, measured in a vertical plane \emph{perpendicular to the strike}. It is given as an angle and a direction, e.g.\ $30^{\circ}$ towards SE (135$^{\circ}$).
\item The \cle{dip direction} is the compass bearing of the line of greatest slope of the plane; it is always \textbf{perpendicular to the strike}, on the downhill side.
\item An \cle{apparent dip} is the angle of inclination of the plane measured in any vertical section that is \emph{not} perpendicular to the strike.
\end{itemize}
\end{definition}

\begin{popfigure}
\begin{tikzpicture}[scale=1.05, line join=round]
% ground plane
\fill[popgreenL] (-3.4,-1.6) -- (3.4,-1.6) -- (4.4,-0.2) -- (-2.4,-0.2) -- cycle;
\draw[popmuted] (-3.4,-1.6) -- (3.4,-1.6) -- (4.4,-0.2) -- (-2.4,-0.2) -- cycle;
\node[popmuted, font=\small] at (-2.9,-1.35) {horizontal ground};
% bedding plane (tilted)
\fill[popblueL, opacity=0.85] (-2.6,0.4) -- (2.2,1.6) -- (3.2,-0.6) -- (-1.6,-1.8) -- cycle;
\draw[popblue, thick] (-2.6,0.4) -- (2.2,1.6) -- (3.2,-0.6) -- (-1.6,-1.8) -- cycle;
\node[popblue, font=\small] at (2.6,1.35) {bedding plane};
% strike line (horizontal line on the plane)
\draw[popink, very thick] (-2.2,-0.55) -- (2.6,0.85);
\node[popink, font=\small, anchor=south west] at (-2.2,-0.5) {strike line (horizontal)};
% dip direction (line of greatest slope, perpendicular to strike on plane)
\draw[poporange, very thick, -latex] (0.2,-0.05) -- (1.15,-1.35);
\node[poporange, font=\small, anchor=west] at (1.2,-1.2) {dip direction};
% dip angle arc
\draw[popdark, thick] (0.2,-0.05) -- (1.55,-0.05);
\draw[popdark, thick] (1.15,-0.05) arc[start angle=0, end angle=-58, radius=0.95];
\node[popdark, font=\small] at (1.75,-0.55) {dip $\theta$};
% compass
\draw[popmuted, -latex] (-3.0,1.4) -- (-3.0,2.2) node[above, font=\small]{N};
\end{tikzpicture}
\legende{Block diagram of a tilted bedding plane. The strike is a horizontal line on the plane; the true dip is the maximum slope, measured at right angles to the strike; the dip direction points down the greatest slope.}
\end{popfigure}

\begin{attention}[Strike has no arrow; dip has]
A strike of 045$^{\circ}$ is the same line as 225$^{\circ}$ --- a line has no ``downhill'' end. The dip direction, by contrast, \emph{must} carry a sense (towards SE, not NW). On the map symbol, the long bar is the strike and the short tick points in the dip direction, with the dip angle written beside it.
\end{attention}

\courssub{True Dip and Apparent Dip}

Cut a vertical trench across a dipping bed. If your trench runs exactly perpendicular to the strike, the bed appears to slope at its steepest: you see the \textbf{true dip}. If your trench runs obliquely, the bed appears to slope more gently: you see an \textbf{apparent dip}. A trench cut exactly \emph{along} the strike shows the bed horizontal --- an apparent dip of $0^{\circ}$.

\begin{propriete}[True dip is the maximum dip]
The true dip $\theta$ is always greater than or equal to the apparent dip $\alpha$ measured in any other vertical section:
\[
\theta \;\ge\; \alpha .
\]
If $\beta$ is the angle between the section direction and the true dip direction, then
\[
\tan \alpha \;=\; \tan \theta \, \cos \beta .
\]
(The derivation from the geometry of two vertical sections is instructive and may be asked as a short structured question.)
\end{propriete}

\textbf{Why is this true?} Consider a bed dipping at true dip $\theta$. Take a horizontal distance $d$ travelled exactly in the dip direction: the bed drops by $h = d\tan\theta$. Now travel the same horizontal distance $d$ along a direction making angle $\beta$ with the dip direction. Your progress \emph{in the dip direction} is only $d\cos\beta$, so the drop is $h' = d\cos\beta \tan\theta$, and the apparent dip satisfies $\tan\alpha = h'/d = \tan\theta\cos\beta$. Since $\cos\beta \le 1$, we get $\tan\alpha \le \tan\theta$, hence $\alpha \le \theta$. \hfill$\square$

\begin{exemplebox}[Apparent dip in a quarry face]
A bed has true dip $40^{\circ}$ towards the south. A quarry face trends N\,060$^{\circ}$\,E -- S\,240$^{\circ}$\,W, i.e.\ its trace makes $\beta = 60^{\circ}$ with the south dip direction. Then
\[
\tan\alpha = \tan 40^{\circ}\times \cos 60^{\circ} = 0.839 \times 0.5 = 0.420
\quad\Rightarrow\quad \alpha \approx 22.8^{\circ}.
\]
The bed appears to dip at only about $23^{\circ}$ in the quarry face --- barely more than half the true dip.
\end{exemplebox}

\courssub{Thickness of a Bed: True, Vertical and Outcrop Width}

When a geologist says a bed is ``20 m thick'', this means the \cle{true thickness}: the \emph{perpendicular} distance between its two bounding bedding planes. But on a map we never measure thickness directly --- we measure the \cle{outcrop width} $w$, the horizontal distance on the map between the two boundary traces of the bed. Confusing the two is the most common error in the GCE practical paper.

\begin{definition}[Three different ``thicknesses'']
\begin{itemize}
\item \cle{True thickness} $t$: perpendicular distance between the footwall and hanging-wall bedding planes of the bed.
\item \cle{Vertical thickness} $v$: thickness measured along a vertical line (e.g.\ in a borehole or well).
\item \cle{Outcrop width} $w$: horizontal map distance between the two contacts of the bed, measured \emph{perpendicular to the strike} on flat ground.
\end{itemize}
\end{definition}

\begin{popfigure}
\begin{tikzpicture}[scale=1.15]
% dipping bed between two parallel planes, flat ground
\fill[popblueL] (0,2.6) -- (4.6,0.0) -- (4.6,-1.3) -- (0,1.3) -- cycle;
\draw[popblue, thick] (0,2.6) -- (4.6,0.0);
\draw[popblue, thick] (0,1.3) -- (4.6,-1.3);
\node[popblue, font=\small] at (3.9,0.9) {bed};
% ground line
\draw[popink, very thick] (-0.3,2.6) -- (0,2.6) (0,2.6) -- (4.6,2.6) (4.6,2.6) -- (4.9,2.6);
\draw[popink, very thick] (-0.3,2.6) -- (4.9,2.6);
\node[popink, font=\small, anchor=south] at (0.6,2.65) {flat ground};
% outcrop width w
\draw[poporange, very thick, latex-latex] (0.62,2.75) -- (3.15,2.75);
\node[poporange, font=\small, anchor=south] at (1.9,2.8) {outcrop width $w$};
% true thickness t (perpendicular to bedding)
\draw[popgreen, very thick, latex-latex] (2.3,1.19) -- (2.82,2.24);
\node[popgreen, font=\small, anchor=south east] at (2.45,1.85) {$t$};
% vertical thickness v
\draw[poppurple, very thick, latex-latex] (3.8,2.6) -- (3.8,-0.68);
\node[poppurple, font=\small, anchor=west] at (3.85,1.0) {vertical $v$};
% dip angle
\draw[popdark] (1.6,2.6) -- (1.6,2.05);
\draw[popdark] (1.6,2.6) arc[start angle=180, end angle=209.7, radius=1.0];
\node[popdark, font=\small] at (1.05,2.28) {$\theta$};
% right angle marker for t
\draw[popgreen] (2.44,1.34) -- (2.58,1.29) -- (2.63,1.43);
\end{tikzpicture}
\legende{Cross-section perpendicular to strike, on horizontal ground. The outcrop width $w$ is measured along the surface; the true thickness $t$ is perpendicular to the bedding; the vertical thickness $v$ is what a borehole would record. The dip is $\theta$.}
\end{popfigure}

\begin{propriete}[Thickness formulae]
For a bed dipping at angle $\theta$ on \textbf{horizontal ground}, with outcrop width $w$ measured perpendicular to strike:
\[
\boxed{\,t = w\,\sin\theta\,} \qquad\text{and}\qquad v = w\,\tan\theta , \qquad\text{so also}\qquad v = \dfrac{t}{\cos\theta}.
\]
These relations follow directly from the right-angled triangle formed by the ground surface, the bedding plane and the perpendicular between the two contacts. They are examinable and must be known with a labelled sketch.
\end{propriete}

\begin{attention}[Outcrop width is NOT thickness]
A vertical bed ($\theta = 90^{\circ}$) has $t = w$: width equals thickness. But a gently dipping bed can outcrop over a huge width while being quite thin: at $\theta = 10^{\circ}$, a bed only $t = 17$ m thick outcrops over $w = 100$ m. Conversely, on a steep \emph{dip slope} a thick bed may show a narrow outcrop. \textbf{Never quote an outcrop width as a thickness.} Width depends on dip \emph{and} on topography.
\end{attention}

\begin{methode}[Computing true thickness --- exam technique]
\begin{enumerate}
\item \textbf{Draw a small sketch} first: a cross-section perpendicular to strike showing ground, the two contacts, the dip angle $\theta$ and the width $w$.
\item Check the ground is horizontal (same contour on both contacts). If not, correct for slope first (Section 8).
\item Identify the right-angled triangle: $w$ is the hypotenuse-side along the ground, $t$ is opposite to $\theta$.
\item Apply $t = w\sin\theta$; give the answer with units and a sensible number of significant figures.
\item Sanity check: $t$ must be \emph{smaller} than $w$ (for $\theta < 90^{\circ}$).
\end{enumerate}
\end{methode}

\begin{exoresolu}[Thickness of a limestone bed near the Benue trough]
On a 1/50\,000 map of flat terrain, the two contacts of a limestone bed are $4.0$ mm apart, measured perpendicular to the strike. The dip symbol reads $30^{\circ}$. Calculate the true thickness and the vertical thickness of the bed.
\tcblower
\textbf{Step 1 --- real width.} Scale 1/50\,000: 1 mm on the map $= 50$ m on the ground, so
\[
w = 4.0 \times 50 = 200\ \mathrm{m}.
\]
\textbf{Step 2 --- sketch.} Cross-section perpendicular to strike: horizontal ground, bed dipping $30^{\circ}$, width $w = 200$ m along the surface.

\textbf{Step 3 --- true thickness.}
\[
t = w\sin\theta = 200 \times \sin 30^{\circ} = 200 \times 0.5 = 100\ \mathrm{m}.
\]
\textbf{Step 4 --- vertical thickness.}
\[
v = w\tan\theta = 200 \times \tan 30^{\circ} = 200 \times 0.577 \approx 115\ \mathrm{m}.
\]
\textbf{Check:} $t = 100\ \mathrm{m}$, $v \approx 115\ \mathrm{m}$, $w = 200\ \mathrm{m}$. Since $\theta = 30^{\circ} < 45^{\circ}$, we have $\sin\theta < \tan\theta < 1$, hence $t < v < w$. Indeed $100 < 115 < 200$. \checkmark{}
\end{exoresolu}

\courssub{Effect of Topography: the Rule of V's}

So far the ground was flat. In the field --- the hills around Dschang, the valleys of the Adamawa Plateau --- contacts cross valleys and spurs, and their map traces bend. The way a contact bends as it crosses a valley tells us the dip, sometimes even quantitatively. This is summarised in the \cle{Rule of V's}.

\begin{definition}[Rule of V's]
When a contact (or any planar structure) crosses a valley, its outcrop trace forms a ``V'' whose shape and pointing direction depend on the dip of the plane relative to the valley gradient:
\begin{itemize}
\item \textbf{Horizontal bed:} the contact follows the contour lines; the V points \emph{upstream} (up-valley), exactly like the contours.
\item \textbf{Vertical bed:} the contact is a \emph{straight line}; it crosses the valley with no deflection at all.
\item \textbf{Bed dipping downstream more steeply than the valley gradient:} the V points \emph{downstream}.
\item \textbf{Bed dipping upstream, or downstream but more gently than the valley gradient:} the V points \emph{upstream}, but the contact V is ``sharper'' or ``blunter'' than the contour V depending on the dip.
\end{itemize}
\end{definition}

\begin{popfigure}
\begin{tikzpicture}[scale=0.92]
% ---- Panel 1: horizontal bed ----
\begin{scope}
\draw[popmuted, dashed] (-1.6,1.6) -- (0,0.2) -- (1.6,1.6);
\draw[popmuted, dashed] (-1.6,0.6) -- (0,-0.8) -- (1.6,0.6);
\draw[popmuted, dashed] (-1.6,-0.4) -- (0,-1.8) -- (1.6,-0.4);
\draw[popblue, very thick] (-1.6,0.6) -- (0,-0.8) -- (1.6,0.6);
\draw[popink, -latex] (0,-2.1) -- (0,-1.2);
\node[font=\small, align=center] at (0,-2.5) {(a) horizontal bed:\\ V points upstream};
\node[popmuted, font=\scriptsize] at (1.5,-1.6) {contours};
\node[popblue, font=\scriptsize] at (1.35,0.95) {contact};
\end{scope}
% ---- Panel 2: vertical bed ----
\begin{scope}[xshift=5.2cm]
\draw[popmuted, dashed] (-1.6,1.6) -- (0,0.2) -- (1.6,1.6);
\draw[popmuted, dashed] (-1.6,0.6) -- (0,-0.8) -- (1.6,0.6);
\draw[popmuted, dashed] (-1.6,-0.4) -- (0,-1.8) -- (1.6,-0.4);
\draw[popgreen, very thick] (-0.5,1.9) -- (-0.5,-2.0);
\draw[popink, -latex] (0,-2.1) -- (0,-1.2);
\node[font=\small, align=center] at (0,-2.5) {(b) vertical bed:\\ straight trace, no V};
\node[popgreen, font=\scriptsize] at (-0.95,1.6) {contact};
\end{scope}
% ---- Panel 3: dipping bed (downstream dip) ----
\begin{scope}[xshift=10.4cm]
\draw[popmuted, dashed] (-1.6,1.6) -- (0,0.2) -- (1.6,1.6);
\draw[popmuted, dashed] (-1.6,0.6) -- (0,-0.8) -- (1.6,0.6);
\draw[popmuted, dashed] (-1.6,-0.4) -- (0,-1.8) -- (1.6,-0.4);
\draw[poporange, very thick] (-1.6,-0.9) -- (0,0.5) -- (1.6,-0.9);
\draw[popink, -latex] (0,-2.1) -- (0,-1.2);
\node[font=\small, align=center] at (0,-2.5) {(c) bed dips downstream:\\ V points downstream};
\node[poporange, font=\scriptsize] at (1.35,-1.25) {contact};
\end{scope}
\end{tikzpicture}
\legende{Plan view of the Rule of V's. Dashed lines: contours bending upstream across a valley (stream flows downwards, arrow). (a) A horizontal contact mimics the contours. (b) A vertical contact is straight. (c) A contact dipping downstream (steeper than the valley) makes a V pointing downstream.}
\end{popfigure}

\begin{methode}[Applying the Rule of V's to deduce dip direction]
\begin{enumerate}
\item Locate where a contact crosses a stream valley on the map.
\item Note the direction of stream flow (contour V's point upstream).
\item Observe the contact: straight line $\Rightarrow$ vertical bed; V identical to contours $\Rightarrow$ horizontal bed; V pointing downstream $\Rightarrow$ bed dips downstream; V pointing upstream but not matching contours $\Rightarrow$ bed dips upstream, or downstream with a dip gentler than the valley slope.
\item State the dip direction as a compass bearing, and justify it in one sentence quoting the V pattern.
\end{enumerate}
\end{methode}

\courssub{Reading Strata Geometry from Map Patterns Alone}

Even without dip symbols, the outcrop pattern reveals the attitude of strata:

\begin{center}
\begin{tabularx}{\linewidth}{|l|X|X|}
\hline
\rowcolor{popblueL} \textbf{Attitude} & \textbf{Map evidence} & \textbf{Width--thickness relation} \\
\hline
Horizontal beds & Contacts parallel to contours; V's point upstream; outcrops cap hills and follow valley sides & $w$ varies with slope; $t$ = contour interval between contacts \\
\hline
Vertical beds & Straight contacts, independent of topography & $w = t$ exactly \\
\hline
Dipping beds & Contacts cut across contours; V's governed by the Rule of V's & $t = w\sin\theta$ on flat ground \\
\hline
\end{tabularx}
\end{center}

In \cle{cuesta} landscapes --- typical of gently dipping sedimentary country such as parts of the Benue trough --- an inclined resistant bed forms an asymmetric ridge: a gentle \cle{dip slope} following the bedding surface, and a steep \cle{scarp slope} cutting across the bedding. The outcrop width of the bed is \textbf{wide on the gentle dip slope} and \textbf{narrow on the steep scarp slope}, even though the true thickness is the same everywhere. This is a direct field confirmation that width is controlled by the angle between the ground surface and the bedding, not by thickness alone.

\begin{savaistu}[Cuestas and dip slopes in Cameroon]
The sedimentary plateaux and escarpments of the Adamawa and the Benue sedimentary basin owe their asymmetric profiles to gently dipping sandstone and limestone beds: long gentle dip slopes on one side, abrupt scarp slopes on the other. Recognising this pattern on a topographic map is often the first step of the GCE map-interpretation question.
\end{savaistu}

\begin{exercice}[Thickness and dip from map data]
On flat ground at 1/25\,000, a sandstone bed outcrops over a width of 8.0 mm measured perpendicular to strike, with a dip of $45^{\circ}$. (a) Calculate the true thickness. (b) Calculate the vertical thickness. (c) The same bed crosses a valley whose stream flows towards the north; the contact makes a V pointing north. State the dip direction of the bed relative to the valley.
\end{exercice}

\begin{corrige}[Exercise 1]
(a) Real width: $w = 8.0 \times 25 = 200\ \mathrm{m}$. Then $t = w\sin 45^{\circ} = 200 \times 0.707 \approx 141\ \mathrm{m}$.
(b) $v = w\tan 45^{\circ} = 200 \times 1 = 200\ \mathrm{m}$.
(c) The V points downstream (north, the flow direction), so by the Rule of V's the bed dips downstream, i.e.\ towards the north, at an angle steeper than the valley gradient.
\end{corrige}

\begin{aretenir}[Key points of Section 5]
\begin{itemize}
\item A plane in space is fixed by its \cle{strike} (horizontal line on the plane) and its \cle{true dip} (maximum inclination, perpendicular to strike, with a dip direction).
\item True dip is always the maximum: $\tan\alpha = \tan\theta\cos\beta$, so apparent dip $\alpha \le \theta$.
\item \cle{True thickness} $t$ is perpendicular to bedding; on flat ground $t = w\sin\theta$ and $v = w\tan\theta$.
\item Outcrop width is \emph{not} thickness: it depends on dip and topography.
\item Rule of V's: horizontal contacts follow contours; vertical contacts are straight; dipping contacts make V's whose pointing direction reveals the dip direction.
\item Always sketch a cross-section before any thickness calculation.
\end{itemize}
\end{aretenir}

With the geometry of single beds mastered, we are ready in Section 6 to use \emph{strike lines} (structure contours) to reconstruct the full attitude of a plane from map data, and to solve the celebrated three-point problem.

\coursec{Section 6: Strike Lines and the Three-Point Problem}

In the previous section we mastered the \cle{geometrical properties of strata} — strike, dip, apparent dip and thickness. We learned how to measure these attitudes in the field and how they appear on a map. Now we turn to a powerful cartographic tool that allows us to \emph{predict} the subsurface geometry of planar beds from a few scattered outcrop measurements: the \cle{strike line} (also called a \cle{structure contour}). This concept is the bridge between isolated spot heights on a map and the full three‑dimensional orientation of a geological plane. It is also the key to solving the classic \cle{three‑point problem}, a near‑annual feature of the GCE Advanced Level practical examination.

\courssub{What Is a Strike Line?}

\begin{definition}[Strike Line / Structure Contour]
A \textbf{strike line} (or \textbf{structure contour}) is a line drawn on a geological map that joins all points of \emph{equal elevation} (height above a datum, usually sea level) on a specific planar geological surface — typically a bedding plane, a fault plane, or an unconformity surface.
\end{definition}

\begin{aretenir}[Key Properties of Strike Lines on a Planar Surface]
\begin{itemize}
\item Strike lines on a single planar surface are \textbf{parallel} to each other.
\item They are \textbf{equally spaced} (horizontally) when the contour interval is constant.
\item The \textbf{strike direction} of the plane is exactly the direction of the strike lines.
\item The \textbf{dip direction} is perpendicular to the strike lines, pointing \emph{towards decreasing height values}.
\item The \textbf{dip amount} ($\delta$) is related to the horizontal spacing ($S$) between adjacent strike lines of a known contour interval ($\Delta h$) by:
\[
\tan\delta = \frac{\Delta h}{S}
\]
\end{itemize}
\end{aretenir}

\begin{attention}[Planar Assumption]
The properties above (parallel, equally spaced) hold \emph{only} if the geological surface is truly planar (constant strike and dip). Folded or curved surfaces produce strike lines that curve and vary in spacing — a vital clue to non‑planar geometry.
\end{attention}

\courssub{Constructing Strike Lines from Spot Heights}

In the field or on a map, we rarely have a continuous line of equal height. Instead we have \cle{spot heights}: the elevation of a bedding plane measured at discrete outcrops or boreholes. To draw strike lines we \emph{interpolate} between these known points.

\begin{methode}[Drawing Strike Lines from Spot Heights]
\begin{enumerate}
\item Plot all known spot heights of the target plane on the map (outcrop points, borehole intersections).
\item Choose a contour interval (e.g., $10\ \mathrm{m}$, $20\ \mathrm{m}$) appropriate to the relief and map scale.
\item For each pair of adjacent points with heights $h_1$ and $h_2$ ($h_1 < h_2$), locate the point where the plane reaches a chosen contour value $H$ ($h_1 \le H \le h_2$) by linear interpolation:
\[
\text{distance from } h_1 = \frac{H - h_1}{h_2 - h_1} \times \text{map distance between the two points}
\]
\item Mark the interpolated points for the same contour value $H$ and join them with a smooth line — this is the strike line for height $H$.
\item Repeat for other contour values. Verify that the resulting lines are parallel and equally spaced (if the plane is planar).
\item Draw a \textbf{dip arrow} perpendicular to the strike lines, pointing from higher to lower contour values. Label the dip amount calculated from the spacing.
\end{enumerate}
\end{methode}

\begin{exemplebox}[Constructing Strike Lines — Simple Example]
Two outcrops of the same sandstone bed are mapped:
\begin{itemize}
\item Point A at elevation $120\ \mathrm{m}$
\item Point B at elevation $160\ \mathrm{m}$
\item Horizontal map distance AB = $200\ \mathrm{m}$
\end{itemize}
We want to draw the $140\ \mathrm{m}$ strike line (contour interval $20\ \mathrm{m}$).
\[
\text{Distance from A} = \frac{140-120}{160-120} \times 200 = \frac{20}{40} \times 200 = 100\ \mathrm{m}
\]
Measure $100\ \mathrm{m}$ from A towards B along line AB, mark the point, and draw the strike line through it. The dip direction is perpendicular to this line, pointing from B ($160\ \mathrm{m}$) towards A ($120\ \mathrm{m}$). The horizontal spacing between the $120\ \mathrm{m}$ and $140\ \mathrm{m}$ strike lines is $100\ \mathrm{m}$, so
\[
\tan\delta = \frac{20}{100} = 0.2 \quad\Rightarrow\quad \delta \approx 11.3^\circ
\]
\end{exemplebox}

\courssub{Uses of Strike Lines}

Strike lines are not just an academic exercise; they are practical tools for the geologist:

\begin{itemize}
\item \textbf{Predicting outcrop position:} Extend strike lines across the map to see where a bed \emph{should} outcrop (useful for mapping covered areas).
\item \textbf{Locating a bed at depth:} In borehole planning, strike lines tell you at what depth a target horizon will be encountered at any map location.
\item \textbf{Borehole planning:} Choose a drill site where the target bed is shallowest (lowest overburden) or where structural complications (faults, folds) are minimal.
\item \textbf{Volume estimation:} Combined with isopach maps, strike lines help estimate reserves of coal, limestone, or aquifer thickness.
\item \textbf{Structural analysis:} Curving or converging/diverging strike lines immediately reveal folds, domes, basins or faults.
\end{itemize}

\begin{savaistu}[Cameroon Application — Limestone Quarrying at Figuil]
The Figuil limestone quarry (North Region) exploits a gently dipping Cretaceous limestone layer. Strike lines constructed from borehole data allow the quarry engineers to plan the pit progression: they follow the strike direction to keep a constant floor elevation, and the dip angle (calculated from strike‑line spacing) determines the maximum safe bench height. This is a direct application of the methods you are learning.
\end{savaistu}

\courssub{The Three‑Point Problem}

\begin{definition}[The Three‑Point Problem]
Given \textbf{three points} of known map position (X, Y) and known elevation (Z) that lie on the \emph{same planar geological surface}, determine the \textbf{strike} and \textbf{dip} of that plane.
\end{definition}

\begin{aretenir}[Why Three Points?]
A plane in 3D space is uniquely defined by three non‑collinear points. In geology, these are typically:
\begin{itemize}
\item Three outcrop points of the same bed.
\item One outcrop + two borehole intersections.
\item Three borehole intersections.
\end{itemize}
The three‑point problem is the \textbf{fundamental graphical method} for finding the attitude of a plane from map data. It is a \emph{core practical skill} for the GCE exam.
\end{aretenir}

\courssub{Graphical Solution of the Three‑Point Problem — Step by Step}

\begin{methode}[5‑Step Graphical Solution of the Three‑Point Problem]
\begin{enumerate}
\item \textbf{Plot and label} the three points (A, B, C) with their elevations ($h_A, h_B, h_C$) on the map. Identify the \textbf{highest (H)}, \textbf{lowest (L)} and \textbf{intermediate (I)} points.
\item \textbf{Join H and L} with a straight line. This line lies in the plane and its projection is the line of \emph{maximum slope} direction (not the strike).
\item \textbf{Interpolate} the point $I'$ on line HL that has the \emph{same elevation as the intermediate point I}. Using linear interpolation:
\[
\frac{HI'}{HL} = \frac{h_H - h_I}{h_H - h_L}
\]
Mark $I'$ on line HL.
\item \textbf{Draw the strike line} through I and $I'$. This line joins two points of equal elevation ($h_I$), so it is a \textbf{strike line} (structure contour) for that elevation. The \textbf{strike direction} is the bearing of this line.
\item \textbf{Determine dip:}
\begin{itemize}
\item Dip direction is perpendicular to the strike line, pointing from H towards L (i.e., from higher to lower elevation).
\item Measure the horizontal distance $S$ (perpendicular distance) between the strike line through I/$I'$ and the next strike line (e.g., through H or L) for a known contour interval $\Delta h$.
\item Calculate dip: $\tan\delta = \frac{\Delta h}{S}$.
\end{itemize}
\end{enumerate}
\end{methode}

\begin{attention}[Common Pitfalls in the Three‑Point Problem]
\begin{itemize}
\item \textbf{Collinear points:} If the three points lie on a straight line on the map, the problem has infinite solutions (the plane can rotate about that line). The exam will never give collinear points for a three‑point problem.
\item \textbf{Wrong interpolation:} Always interpolate \emph{on the line joining highest and lowest}. Interpolating on the other lines gives the wrong strike.
\item \textbf{Forgetting the contour interval:} The dip calculation requires the \emph{vertical difference} between two adjacent strike lines you have drawn. Use the actual elevation difference, not an arbitrary interval.
\item \textbf{Units:} Map distance $S$ must be in the \emph{same units} as $\Delta h$ (usually metres). Convert using the map scale before calculating $\tan\delta$.
\end{itemize}
\end{attention}

\courssub{Worked Example — Complete Numerical Solution}

\begin{exoresolu}[Three‑Point Problem: Worked Example]
\textbf{Statement:} Three boreholes intersect the top of the \emph{Basement Granite} at the following map coordinates (local grid, metres) and elevations (m above sea level):
\begin{itemize}
\item BH1: (1000, 2000), elevation $850\ \mathrm{m}$
\item BH2: (1400, 2200), elevation $780\ \mathrm{m}$
\item BH3: (1200, 2500), elevation $820\ \mathrm{m}$
\end{itemize}
Map scale: 1:10,000. Determine the strike and dip of the granite surface.

\tcblower
\textbf{Solution:}

\textbf{Step 1: Identify H, L, I.}
\begin{itemize}
\item Highest (H): BH1, $h_H = 850\ \mathrm{m}$
\item Lowest (L): BH2, $h_L = 780\ \mathrm{m}$
\item Intermediate (I): BH3, $h_I = 820\ \mathrm{m}$
\end{itemize}

\textbf{Step 2: Join H and L.}
Draw line BH1–BH2 on the map. Map distance HL:
\[
HL = \sqrt{(1400-1000)^2 + (2200-2000)^2} = \sqrt{400^2 + 200^2} = \sqrt{200000} \approx 447.2\ \mathrm{m}
\]

\textbf{Step 3: Interpolate $I'$ on HL at elevation $820\ \mathrm{m}$.}
\[
\frac{HI'}{HL} = \frac{h_H - h_I}{h_H - h_L} = \frac{850-820}{850-780} = \frac{30}{70} = \frac{3}{7}
\]
\[
HI' = \frac{3}{7} \times 447.2 \approx 191.7\ \mathrm{m}
\]
Measure $191.7\ \mathrm{m}$ from BH1 towards BH2 along line HL, mark point $I'$.

\textbf{Step 4: Draw strike line through I (BH3) and $I'$.}
This is the $820\ \mathrm{m}$ structure contour. Measure its bearing on the map:
Using coordinates:
BH3 (1200, 2500), $I'$ lies on BH1–BH2 line.
Parametric coordinates of $I'$:
\[
x_{I'} = 1000 + \frac{3}{7}(400) \approx 1171.4,\quad y_{I'} = 2000 + \frac{3}{7}(200) \approx 2085.7
\]
Strike line direction vector: $(1200-1171.4,\ 2500-2085.7) = (28.6,\ 414.3)$.
Bearing = $\arctan(28.6/414.3) \approx 3.9^\circ$ east of north $\Rightarrow$ \textbf{Strike $\approx$ N04°E} (or $004^\circ$ / $184^\circ$).

\textbf{Step 5: Determine dip.}
Draw the $850\ \mathrm{m}$ strike line through H (BH1) parallel to the $820\ \mathrm{m}$ line.
Measure perpendicular horizontal distance $S$ between the $820\ \mathrm{m}$ and $850\ \mathrm{m}$ strike lines.
From geometry, the perpendicular distance between two parallel lines through H and I is the component of HI perpendicular to strike.
Vector HI = $(1200-1000,\ 2500-2000) = (200,\ 500)$.
Strike direction unit vector $\hat{s} \approx (0.068,\ 0.998)$.
Perpendicular component magnitude:
\[
S = |\vec{HI} \times \hat{s}| = |200\times0.998 - 500\times0.068| \approx |199.6 - 34| = 165.6\ \mathrm{m}
\]
(Alternatively, measure directly on the scaled map.)

Contour interval $\Delta h = 850 - 820 = 30\ \mathrm{m}$.
\[
\tan\delta = \frac{\Delta h}{S} = \frac{30}{165.6} \approx 0.181 \quad\Rightarrow\quad \delta \approx 10.3^\circ
\]
Dip direction is perpendicular to strike (N04°E), pointing from H ($850\ \mathrm{m}$) to L ($780\ \mathrm{m}$). The dip direction bearing is $004^\circ + 90^\circ = 094^\circ$ (or $184^\circ - 90^\circ = 094^\circ$).
\textbf{Result: Strike $004^\circ$, Dip $10^\circ$ towards $094^\circ$ (or simply $10^\circ$ East).}

\textbf{Verification:} Check that the third point (BH2, $780\ \mathrm{m}$) lies on the $780\ \mathrm{m}$ strike line parallel to the others. The perpendicular distance from $820\ \mathrm{m}$ to $780\ \mathrm{m}$ should be $\frac{40}{30} \times 165.6 \approx 220.8\ \mathrm{m}$. If the map shows this, the plane is planar.
\end{exoresolu}

\begin{popfigure}
\begin{tikzpicture}[scale=0.8]
% Map view: three boreholes, interpolated point, strike lines, dip arrow
% Scale: 1 unit = 50 m
% Real coords: BH1(1000,2000), BH2(1400,2200), BH3(1200,2500)
% TikZ coords:
\coordinate (BH1) at (20, 40);
\coordinate (BH2) at (28, 44);
\coordinate (BH3) at (24, 50);
% Interpolated point I' on BH1-BH2 at 3/7 from BH1 (requires calc library)
\coordinate (Ip) at ($(BH1)!3/7!(BH2)$); % (23.4286, 41.7143)
% Strike vector (from Ip to BH3) in TikZ units: (0.5714, 8.2857)
% Strike lines: through BH3 & Ip (solid), through BH1 (dashed, 850m), through BH2 (dashed, 780m)
% Extend strike lines using strike vector
\coordinate (S1) at ($(BH3) + (-1.5*0.5714, -1.5*8.2857)$);
\coordinate (S2) at ($(BH3) + (2.5*0.5714, 2.5*8.2857)$);
\coordinate (S1_H) at ($(BH1) + (-1.5*0.5714, -1.5*8.2857)$);
\coordinate (S2_H) at ($(BH1) + (2.5*0.5714, 2.5*8.2857)$);
\coordinate (S1_L) at ($(BH2) + (-1.5*0.5714, -1.5*8.2857)$);
\coordinate (S2_L) at ($(BH2) + (2.5*0.5714, 2.5*8.2857)$);
% Draw strike lines
\draw[thick, popblue] (S1) -- (S2) node[midway, above, sloped, popblue, text width=3cm, align=center] {Strike line 820 m};
\draw[thick, popblue, dashed] (S1_H) -- (S2_H) node[midway, above, sloped, popblue, text width=3cm, align=center] {850 m};
\draw[thick, popblue, dashed] (S1_L) -- (S2_L) node[midway, below, sloped, popblue, text width=3cm, align=center] {780 m};
% Points
\fill[poporange] (BH1) circle (3pt) node[above left] {BH1 (850 m)};
\fill[poporange] (BH2) circle (3pt) node[below right] {BH2 (780 m)};
\fill[poporange] (BH3) circle (3pt) node[above] {BH3 (820 m)};
\fill[popgreen] (Ip) circle (3pt) node[below left] {$I'$ (820 m)};
% Line H-L
\draw[popdark, thick] (BH1) -- (BH2) node[midway, above, sloped] {H–L line};
% Dip arrow: perpendicular to strike, from high to lower elevation.
% Strike angle ~86 deg from x-axis. Dip direction = strike - 90 = -4 deg (or 356 deg).
% Unit vector approx (0.9976, -0.0698). Length 1.5 units.
\draw[popred, very thick, ->] (BH1) -- ++(356:1.5) node[midway, right] {Dip 10°};
% North arrow
\draw[popdark, ->] (35,55) -- ++(90:1) node[above] {N};
\draw (35,55) circle (0.3);
% Scale bar: 200 m = 4 units
\draw[popdark, thick] (15,35) -- ++(4,0) node[midway, below] {200 m};
\end{tikzpicture}
\legende{Map‑view solution of the three‑point problem. The three boreholes (orange) define the plane. The line joining highest (BH1) and lowest (BH2) is used to interpolate point $I'$ at the intermediate elevation (820 m). The strike line through BH3 and $I'$ (solid blue) gives the strike (N04°E). Parallel strike lines through BH1 and BH2 (dashed) allow dip measurement. The dip arrow (red) is perpendicular to strike, pointing towards lower elevations (bearing 094°).}
\end{popfigure}

\begin{popfigure}
\begin{tikzpicture}[scale=1.2]
% Cross-section inset: shows how strike-line spacing converts to dip angle
% Horizontal axis: distance perpendicular to strike
% Vertical axis: elevation
\draw[->] (0,0) -- (6,0) node[right] {Horizontal distance (m)};
\draw[->] (0,0) -- (0,3.5) node[above] {Elevation (m)};
% Strike lines as horizontal lines on the cross-section (perpendicular to strike)
\draw[popblue, thick] (1,2) -- (5,2) node[right] {850 m};
\draw[popblue, thick] (1,1.4) -- (5,1.4) node[right] {820 m};
\draw[popblue, thick] (1,0.8) -- (5,0.8) node[right] {780 m};
% Vertical spacing
\draw[popmuted, <->] (0.5,2) -- (0.5,1.4) node[midway, left] {$\Delta h = 30$ m};
\draw[popmuted, <->] (0.5,1.4) -- (0.5,0.8) node[midway, left] {$\Delta h = 40$ m};
% Horizontal spacing S (perpendicular distance between strike lines)
\draw[popmuted, <->] (1,1.2) -- (3.5,1.2) node[midway, below] {$S = 165.6$ m};
% Dip angle representation
\draw[popred, thick] (1,2) -- (3.5,1.4);
\draw[popred, thick] (3.5,1.4) -- (5.5,0.8);
% Angle arc
\draw[popred, thick] (1.5,1.9) arc (180:165:0.5) node[above left, popred] {$\delta$};
% Formula
\node[popdark, align=center] at (3,3) {$\tan\delta = \dfrac{\Delta h}{S}$};
\end{tikzpicture}
\legende{Cross‑section perpendicular to strike showing the geometric relationship: the dip angle $\delta$ is the angle the plane makes with the horizontal. The tangent of the dip equals the contour interval $\Delta h$ divided by the horizontal spacing $S$ between adjacent strike lines.}
\end{popfigure}

\courssub{Verification and Quality Control}

After solving the three‑point problem, always perform these checks:

\begin{enumerate}
\item \textbf{Parallelism:} Are the strike lines you drew through H, I, L parallel? If not, the surface is not planar (or you made a construction error).
\item \textbf{Equal spacing:} For a constant contour interval, the perpendicular spacing between strike lines should be constant. Measure at several places.
\item \textbf{Third point check:} The lowest point (L) must lie exactly on its predicted strike line. If it falls off, either the plane is non‑planar or the interpolation/measurement has an error.
\item \textbf{Dip consistency:} Calculate dip from two different contour intervals (e.g., 850–820 and 820–780). They should give the same $\delta$ within measurement tolerance.
\end{enumerate}

\begin{experience}[Practical Exercise: Three‑Point Problem on a Map Extract]
\textbf{Materials:} Map extract (scale 1:25,000) with three spot heights on the \emph{Mamfe Sandstone}: Point X ($450\ \mathrm{m}$), Y ($380\ \mathrm{m}$), Z ($410\ \mathrm{m}$). Coordinates given.
\textbf{Task:}
\begin{enumerate}
\item Plot the three points and label elevations.
\item Solve the three‑point problem graphically (use the 5‑step method).
\item State the strike (as a bearing) and dip (amount and direction).
\item Draw the $400\ \mathrm{m}$ and $350\ \mathrm{m}$ strike lines and predict where the sandstone outcrops on the eastern edge of the map.
\item A borehole is planned at grid reference (123456, 789012). At what depth will it hit the top of the sandstone? (Ground elevation at borehole = $520\ \mathrm{m}$).
\end{enumerate}
\textbf{Teacher's note:} This mirrors GCE Practical Question 2 (Map Work). Work neatly; show all construction lines. Marks are awarded for correct interpolation, parallel strike lines, correct dip arrow, and correct depth calculation.
\end{experience}

\courssub{Exam Technique: Mastering the Three‑Point Problem}

\begin{aretenir}[GCE Examination Checklist for Three‑Point Problem]
\begin{itemize}
\item \textbf{Read the question carefully:} Are the points outcrops or boreholes? What is the contour interval required?
\item \textbf{Scale conversion:} Always convert map distances to ground distances (metres) before calculating $\tan\delta$. $\delta = \arctan(\frac{\Delta h}{S_{\text{ground}}})$.
\item \textbf{Strike notation:} Give strike as a \textbf{three‑figure bearing} (e.g., $004^\circ$, $135^\circ$) or as \textbf{N04°E / S04°W}. Both are accepted; be consistent.
\item \textbf{Dip notation:} Dip amount in degrees, dip direction as a bearing (e.g., $10^\circ$ towards $094^\circ$). Never just "$10^\circ$ East" unless the strike is exactly N‑S.
\item \textbf{Construction lines:} Do \emph{not} erase your interpolation lines (H–L line, $I'$ point). Examiners award method marks for visible construction.
\item \textbf{Neatness:} Use a sharp pencil, fine lines, clear labels. A messy map loses communication marks.
\end{itemize}
\end{aretenir}

\begin{attention}[Trap: Apparent Dip vs True Dip in Cross‑Sections]
When you later construct a cross‑section (Sections 7–8), the dip you measure \emph{on the cross‑section line} is the \textbf{apparent dip} unless the section line is perpendicular to strike. The three‑point problem gives you the \textbf{true dip}. Always check the angle between your section line and the strike direction before plotting the bed in section.
\end{attention}

\courssub{Summary}

\begin{aretenir}[Section 6 — Key Takeaways]
\begin{itemize}
\item A \textbf{strike line} joins points of equal elevation on a planar surface. On a plane, strike lines are parallel, equally spaced, and their direction = strike of the plane.
\item \textbf{Dip direction} $\perp$ strike lines, towards lower values. \textbf{Dip amount}: $\tan\delta = \frac{\Delta h}{S}$.
\item \textbf{Three‑point problem:} Three known points on a plane $\rightarrow$ unique strike and dip.
\item \textbf{Graphical solution (5 steps):} Plot $\rightarrow$ Join H–L $\rightarrow$ Interpolate $I'$ at $h_I$ $\rightarrow$ Strike line through I and $I'$ $\rightarrow$ Measure dip from spacing.
\item \textbf{Uses:} Outcrop prediction, borehole planning, depth estimation, structural analysis.
\item \textbf{Exam alert:} Three‑point problem appears almost every year in GCE Practical. Practise until the 5 steps are automatic. Show all construction lines.
\end{itemize}
\end{aretenir}

You now possess the essential tool to read the third dimension from a two‑dimensional map. In the next sections we will use strike lines extensively to construct geological cross‑sections and to unravel the full geological history of complex map areas. Keep your ruler, protractor and calculator sharp — they are as much a geologist’s tools as the hammer and hand lens!

\coursec{Section 7: Construction of Geological Cross-Sections I — Principles and Simple Cases}

In Section 6 we mastered the \cle{three-point problem}: from three points where a planar bed crops out we deduced its true \cle{strike} and true \cle{dip}. That skill is the key that unlocks the third dimension. A geological map is a plan view; it shows \emph{where} rocks outcrop at the surface. A \cle{geological cross-section} shows \emph{how} those rocks are arranged beneath the surface along a chosen line. Think of it as slicing a cake and looking at the cut face — except the cake is the Earth's crust, the layers are sedimentary beds, lava flows or intrusive bodies, and the knife is a vertical plane we call the \cle{section line}.

\begin{definition}[Geological Cross-Section]
A \textbf{geological cross-section} is a representation of the geology on a vertical plane defined by a chosen straight line (the \textbf{section line}) drawn on the map. It displays the topography along that line, the attitude (strike and dip) of geological contacts, and the inferred geometry of rock units at depth, and it must be consistent with \emph{all} the surface data shown on the map.
\end{definition}

\courssub{Why Do We Draw Cross-Sections?}

A map alone cannot answer questions such as:
\begin{itemize}
  \item How thick is the \cle{basalt} pile on the flank of \cle{Mount Cameroon}?
  \item At what depth would a borehole near Buea encounter the \cle{granitic basement}?
  \item Is a sedimentary sequence in the \cle{Benue Trough} folded into an anticline or a syncline at depth?
\end{itemize}
A cross-section converts the two-dimensional map pattern into a vertical slice that reveals the three-dimensional architecture. It is the principal tool for:
\begin{itemize}
  \item estimating \cle{stratigraphic thicknesses} (true thickness versus apparent thickness);
  \item visualising \cle{fold} hinges, \cle{fault} displacements and \cle{unconformity} surfaces;
  \item planning drilling, tunnelling and hydrogeological investigations;
  \item reconstructing a coherent \cle{geological history} (developed in Section 10).
\end{itemize}

\courssub{Choosing the Section Line}

The section line (labelled \cle{A--B}, \cle{X--Y}, etc.) must be placed strategically:
\begin{enumerate}
  \item \textbf{Perpendicular to the strike} wherever possible. The dip then seen in the section is the \cle{true dip}. If the line is oblique to the strike, the section shows an \cle{apparent dip}, which is always \emph{smaller} than the true dip (see the warning box below).
  \item \textbf{Cross the key structures}: fold axes, major faults, unconformities and the contacts between the main units.
  \item \textbf{Use the topographic contours}: a line running along a ridge crest gives a nearly flat profile and is far less informative than one crossing valleys and hills.
  \item \textbf{Keep it straight}. Curved or kinked section lines are an advanced technique and are not required at GCE Advanced Level.
\end{enumerate}

\begin{attention}[The Apparent Dip Trap]
If the section line is \textbf{not perpendicular to the strike} of a planar surface, the dip seen in the cross-section is the \textbf{apparent dip} $\alpha$, which is \emph{always smaller} than the true dip $\delta$. If $\beta$ is the angle between the section line and the \textbf{strike direction}, then
\[
\tan\alpha = \tan\delta \times \sin\beta .
\]
Check the two extremes: if the section is perpendicular to strike, $\beta = 90^{\circ}$, $\sin\beta = 1$, so $\alpha = \delta$ (true dip); if the section runs \emph{along} the strike, $\beta = 0^{\circ}$, $\sin\beta = 0$, so $\alpha = 0^{\circ}$ (the contact appears horizontal). \textbf{In GCE examinations, always draw your section line perpendicular to the dominant strike unless the question explicitly states otherwise.}
\end{attention}

\courssub{Transferring Topography: From Contours to Profile}

\begin{methode}[Constructing the Topographic Profile]
\begin{enumerate}
  \item On the map, draw the section line \cle{A--B}. Wherever it crosses a topographic contour, make a small tick mark and note the contour elevation.
  \item Draw a horizontal baseline on graph paper. Using the \textbf{horizontal scale} of the map (e.g. 1:50,000), mark the positions of the tick marks along this baseline.
  \item Choose a \textbf{vertical scale} (see \S 7.4) and draw a vertical axis at A (and, for neatness, at B).
  \item At each tick position, plot the corresponding elevation using the vertical scale.
  \item Join the plotted points with a smooth curve — this is the \cle{topographic profile}. Never join them with straight segments: the ground surface is smooth between contours.
\end{enumerate}
\end{methode}

\begin{popfigure}
\begin{tikzpicture}[scale=0.85]
% Map frame
\draw[popline, thick] (0,0) rectangle (8,6);
% Contours
\draw[popblue] (0.5,0.5) .. controls (2,1) and (4,1.5) .. (7.5,2);
\draw[popblue] (0.5,1.5) .. controls (2,2.5) and (4,3) .. (7.5,3.5);
\draw[popblue] (0.5,2.5) .. controls (2,4) and (4,4.5) .. (7.5,5);
\node[popblue, font=\tiny, anchor=west] at (7.55,2) {200 m};
\node[popblue, font=\tiny, anchor=west] at (7.55,3.5) {300 m};
\node[popblue, font=\tiny, anchor=west] at (7.55,5) {400 m};
% Section line A-B
\draw[poporange, very thick] (1,0.5) -- (7,5.5);
\node[poporange, font=\small\bfseries, anchor=east] at (0.9,0.4) {A};
\node[poporange, font=\small\bfseries, anchor=south] at (7.1,5.6) {B};
% Tick marks at contour crossings
\foreach \x/\y/\elev in {2.0/1.25/200, 4.0/2.75/300, 6.0/4.25/400} {
  \draw[popdark, thick] (\x-0.12,\y-0.12) -- (\x+0.12,\y+0.12);
  \draw[popdark, thick] (\x-0.12,\y+0.12) -- (\x+0.12,\y-0.12);
  \node[font=\tiny, popdark, anchor=north west] at (\x+0.05,\y-0.05) {\elev};
}
% Geological contact
\draw[popgreen, thick, dashed] (0,3.2) -- (8,1.2);
\node[popgreen, font=\tiny, anchor=west] at (8.05,1.2) {contact};
\node[draw, fill=white, font=\tiny, anchor=north west] at (0.15,5.85) {Map view};
\end{tikzpicture}
\legende{Step 1: map with section line A--B, topographic contours (blue) and a geological contact (green, dashed). Tick marks show where the section line crosses the contours; their elevations are transferred to the profile.}
\end{popfigure}

\begin{popfigure}
\begin{tikzpicture}[scale=0.85]
% Axes
\draw[popline, thick] (0,0) -- (10,0);
\draw[popline, thick] (0,0) -- (0,5);
\draw[popline, thick] (10,0) -- (10,5);
% Horizontal scale marks
\foreach \x in {0,2,4,6,8,10} {
  \draw (\x,-0.08) -- (\x,0.08);
  \node[font=\tiny, anchor=north] at (\x,-0.15) {\x};
}
\node[font=\tiny, anchor=north] at (10.9,-0.15) {km};
% Vertical scale marks (1 cm = 100 m)
\foreach \y/\e in {1/100, 2/200, 3/300, 4/400, 5/500} {
  \draw (-0.08,\y) -- (0.08,\y);
  \node[font=\tiny, anchor=east] at (-0.15,\y) {\e};
}
\node[font=\tiny, anchor=east] at (-0.15,5.6) {m};
% Profile points and smooth curve
\draw[poporange, very thick] plot[smooth] coordinates {(0,1.5) (2,2.0) (4,3.0) (6,4.0) (8,4.5) (10,4.8)};
\foreach \x/\y/\elev in {0/1.5/150, 2/2.0/200, 4/3.0/300, 6/4.0/400, 8/4.5/450, 10/4.8/480} {
  \fill[popdark] (\x,\y) circle (1.5pt);
  \node[font=\tiny, anchor=south] at (\x,\y+0.15) {\elev};
}
\node[font=\small\bfseries, anchor=east] at (-0.4,0) {A};
\node[font=\small\bfseries, anchor=west] at (10.3,0) {B};
\node[font=\tiny, anchor=north] at (5,-0.7) {Horizontal scale: 1 cm = 1 km \quad Vertical scale: 1 cm = 100 m \quad (VE = 10)};
\end{tikzpicture}
\legende{Step 2: the topographic profile. Contour crossings are plotted at their elevations and joined by a smooth curve. Here the vertical scale is ten times the horizontal scale (VE = 10) so that the gentle relief becomes visible.}
\end{popfigure}

\courssub{Vertical Exaggeration (VE)}

\begin{definition}[Vertical Exaggeration]
\textbf{Vertical Exaggeration (VE)} is the ratio of the vertical scale to the horizontal scale of a cross-section, both expressed as representative fractions:
\[
\boxed{\mathrm{VE} = \frac{\text{vertical scale}}{\text{horizontal scale}}}
\]
Example: vertical scale 1:10,000 and horizontal scale 1:50,000 give $\mathrm{VE} = \dfrac{1/10\,000}{1/50\,000} = 5$.
\end{definition}

\begin{propriete}[Choosing the Vertical Exaggeration]
\begin{itemize}
  \item \textbf{VE = 1 (no exaggeration)}: true geometric relationships are preserved; dip angles measured on the section equal the true dips. Use it for structural analysis, for thickness calculations, and whenever the relief is strong (e.g. the flanks of Mount Cameroon).
  \item \textbf{VE $>$ 1 (commonly 2 to 10)}: vertical distances are stretched. Use it when the relief is gentle (e.g. the sedimentary plains of the \cle{Logone Basin} or the \cle{Benue Trough}) and a true-scale section would look almost flat.
  \item \textbf{VE $<$ 1}: rare; it compresses the vertical dimension.
\end{itemize}
\textbf{Critical rule}: if $\mathrm{VE} \neq 1$, the dips drawn on the section are \textbf{not} the true dips. Either draw the section at VE = 1, or convert each true dip $\delta$ into the dip to be plotted, $\delta_{\text{plot}}$, using
\[
\tan\delta_{\text{plot}} = \mathrm{VE} \times \tan\delta .
\]
(Proof not examinable; the formula follows directly from stretching all vertical distances by the factor VE.)
\end{propriete}

\begin{exemplebox}[Calculating VE and a Plotting Dip]
A section is drawn with horizontal scale 1:50,000 and vertical scale 1:10,000, so $\mathrm{VE} = 5$. A bed has true dip $\delta = 10^{\circ}$. The dip to plot is given by $\tan\delta_{\text{plot}} = 5 \times \tan 10^{\circ} = 5 \times 0.176 = 0.882$, hence $\delta_{\text{plot}} \approx 41^{\circ}$. A gentle $10^{\circ}$ bed therefore appears to dip at about $41^{\circ}$ on the exaggerated section — a dramatic illustration of why dips must never be read directly off an exaggerated section.
\end{exemplebox}

\begin{popfigure}
\begin{tikzpicture}[scale=0.7]
% VE = 1 (left)
\begin{scope}
\draw[popline, thick] (0,0) -- (8,0);
\draw[popline, thick] (0,0) -- (0,3);
\draw[poporange, very thick] plot[smooth] coordinates {(0,0.5) (2,0.7) (4,1.2) (6,1.8) (8,2.2)};
\draw[popgreen, thick] (2,0.7) -- (6,0.24);
\node[font=\small\bfseries, anchor=north] at (4,-0.4) {VE = 1 (true geometry)};
\node[font=\tiny, anchor=north] at (4,-0.8) {dip on section = true dip};
\end{scope}
% VE = 5 (right)
\begin{scope}[xshift=10cm]
\draw[popline, thick] (0,0) -- (8,0);
\draw[popline, thick] (0,0) -- (0,3);
\draw[poporange, very thick] plot[smooth] coordinates {(0,2.5) (2,3.5) (4,6.0) (6,9.0) (8,11.0)};
\begin{scope}
\clip (0,0) rectangle (8,3);
\draw[poporange, very thick] plot[smooth] coordinates {(0,0.5) (2,0.75) (4,1.4) (6,2.2) (8,2.7)};
\draw[popgreen, thick] (2,0.75) -- (4.5,0.09);
\end{scope}
\draw[poporange, very thick] plot[smooth] coordinates {(0,0.5) (2,0.75) (4,1.4) (6,2.2) (8,2.7)};
\draw[popgreen, thick] (2,0.75) -- (4.5,0.09);
\node[font=\small\bfseries, anchor=north] at (4,-0.4) {VE = 5 (exaggerated)};
\node[font=\tiny, anchor=north] at (4,-0.8) {dip on section $>$ true dip};
\end{scope}
\end{tikzpicture}
\legende{The same profile and the same dipping contact (green) drawn at VE = 1 (left) and VE = 5 (right). Vertical exaggeration steepens both slopes and bedding dips; angles measured on an exaggerated section are apparent values, not true dips.}
\end{popfigure}

\courssub{Plotting Contacts and Projecting Boundaries}

Once the topographic profile is drawn, the geology is added:
\begin{enumerate}
  \item \textbf{Mark every contact} where the section line crosses a geological boundary on the map. Transfer its horizontal position (using the map scale) and plot it on the profile at the elevation of the ground surface at that point.
  \item \textbf{Project each contact downward} using its dip:
      \begin{itemize}
        \item \textbf{Horizontal strata} (dip $= 0^{\circ}$): contacts are horizontal lines on the section.
        \item \textbf{Uniformly dipping strata} (constant true dip $\delta$): from each contact point draw a straight line at angle $\delta$ below the horizontal, in the dip direction, extending it down into the subsurface.
      \end{itemize}
  \item \textbf{Ornament each unit} exactly as in the map legend (colour, pattern, formation symbol).
  \item \textbf{Stop the projections} at the base of the section (an arbitrary depth, e.g. sea level) or where they meet another structure (fault, unconformity, intrusion).
\end{enumerate}

\begin{aretenir}[Key Principle — Projection by Dip]
For a planar contact with true dip $\delta$ (section perpendicular to strike), the contact appears on the cross-section as a straight line making the angle $\delta$ with the horizontal. The vertical drop $V$ over a horizontal distance $H$ along the section is
\[
V = H \tan\delta .
\]
This relation lets you compute exactly where a contact will be at any depth — for example, where a borehole would cut it.
\end{aretenir}

\courssub{Six-Step Method for Constructing a Cross-Section}

\begin{methode}[Complete Cross-Section Construction (6 Steps)]
\begin{enumerate}
  \item \textbf{Select and draw the section line} on the map; label its ends (A, B). Make it perpendicular to the regional strike.
  \item \textbf{Construct the topographic profile}: transfer the contour intersections, choose the vertical scale, and state the VE.
  \item \textbf{Mark all geological contacts} where the section line crosses boundaries on the map; plot them at the correct horizontal position and at the surface elevation on the profile.
  \item \textbf{Determine the dip} of each contact (from map dip symbols, from the three-point problem, or from data given in the question). If $\mathrm{VE} \neq 1$, convert each true dip into the plotting dip with $\tan\delta_{\text{plot}} = \mathrm{VE} \times \tan\delta$.
  \item \textbf{Project each contact downward} from its surface point at the appropriate angle, in the dip direction, until the base of the section or until it meets another structure.
  \item \textbf{Finalise}: ornament the units according to the legend; add the title (\emph{Geological section along A--B}), the horizontal and vertical scales, the VE, the orientation of the ends, and dip arrows with values.
\end{enumerate}
\end{methode}

\courssub{Worked Example 1: Horizontal Strata (VE = 1)}

\begin{exoresolu}[Cross-Section Across Horizontal Beds]
\textbf{Statement:} A geological map at 1:50,000 shows three horizontal sedimentary units: \cle{Sandstone} (top), \cle{Shale} (middle), \cle{Limestone} (base). The section line A--B, 4 km long, runs N--S across a valley. Contour crossings along A--B give the elevations: A = 350 m, then 300 m, 250 m (valley floor), 300 m, B = 350 m. On the map, the Sandstone/Shale contact follows the 300 m contour and the Shale/Limestone contact follows the 250 m contour. Draw the cross-section at VE = 1 (vertical scale 1:50,000, i.e. 1 cm = 0.5 km horizontally and 1 cm = 500 m vertically).
\tcblower
\textbf{Solution:}
\begin{enumerate}
  \item \textbf{Topographic profile}: horizontal scale 1 cm = 0.5 km, so 4 km = 8 cm; vertical scale 1 cm = 500 m. Plot A (0 km, 350 m), the 300 m crossings, the 250 m valley floor, and B (4 km, 350 m); join with a smooth valley-shaped curve.
  \item \textbf{Contacts}: the Sandstone/Shale contact crops out wherever the topography is at 300 m — two points, one on each valley side. The Shale/Limestone contact crops out where the topography is at 250 m — at the valley floor. Mark all these points on the profile.
  \item \textbf{Projection}: the beds are horizontal (dip $= 0^{\circ}$), so draw a horizontal line at 300 m across the whole section (Sandstone/Shale) and another at 250 m (Shale/Limestone). Each horizontal contact meets the surface exactly at the plotted outcrop points.
  \item \textbf{Ornament}: Sandstone (stippled), Shale (dashes), Limestone (brick pattern), as in the legend.
  \item \textbf{Label}: title \emph{Section A--B}, both scales, VE = 1, direction A $\rightarrow$ B.
\end{enumerate}
\textbf{Result:} the section shows a valley eroded into flat-lying layers. The contacts are horizontal and parallel; each unit keeps a constant true thickness (here 50 m for the Shale). Note how horizontal contacts \emph{follow the contours} on the map — this is why their outcrop pattern mimics the topography.
\end{exoresolu}

\begin{popfigure}
\begin{tikzpicture}[scale=0.75]
% Axes
\draw[popline, thick] (0,0) -- (10,0);
\draw[popline, thick] (0,0) -- (0,4);
\draw[popline, thick] (10,0) -- (10,4);
% Vertical scale: 1 cm = 100 m for drawing; label accordingly
\foreach \y/\e in {1/100, 2/200, 3/300, 4/400} {
  \draw (-0.08,\y) -- (0.08,\y);
  \node[font=\tiny, anchor=east] at (-0.15,\y) {\e};
}
\node[font=\tiny, anchor=east] at (-0.15,4.5) {m};
% Topography: valley, A=350, floor=250, B=350
\draw[poporange, very thick] plot[smooth] coordinates {(0,3.5) (2,3.0) (3.5,2.5) (5,2.5) (6.5,2.5) (8,3.0) (10,3.5)};
% Horizontal contacts at 300 m and 250 m
\draw[popgreen, thick] (0,3.0) -- (10,3.0);
\draw[popblue, thick] (0,2.5) -- (10,2.5);
% Outcrop points
\fill[popdark] (2,3.0) circle (1.5pt);
\fill[popdark] (8,3.0) circle (1.5pt);
\fill[popdark] (3.9,2.5) circle (1.5pt);
\fill[popdark] (6.1,2.5) circle (1.5pt);
% Ornament: sandstone dots above 300, shale dashes 250-300, limestone bricks below 250
\foreach \x in {0.5,1.5,2.5,3.5,4.5,5.5,6.5,7.5,8.5,9.5} {
  \fill[popmuted] (\x,3.2) circle (0.7pt);
  \draw[popmuted] (\x-0.15,2.75) -- (\x+0.15,2.75);
}
\foreach \x in {0,1,2,3,4,5,6,7,8,9} {
  \draw[popmuted] (\x,0.3) -- (\x+0.6,0.3);
  \draw[popmuted] (\x+0.3,0.3) -- (\x+0.3,0.7);
  \draw[popmuted] (\x,0.7) -- (\x+0.6,0.7);
  \draw[popmuted] (\x,1.1) -- (\x+0.6,1.1);
  \draw[popmuted] (\x+0.3,1.1) -- (\x+0.3,1.5);
  \draw[popmuted] (\x,1.5) -- (\x+0.6,1.5);
  \draw[popmuted] (\x,1.9) -- (\x+0.6,1.9);
  \draw[popmuted] (\x+0.3,1.9) -- (\x+0.3,2.3);
  \draw[popmuted] (\x,2.3) -- (\x+0.6,2.3);
}
% Labels
\node[font=\small\bfseries] at (5,4.3) {Section A--B: horizontal strata (VE = 1)};
\node[font=\tiny, anchor=east] at (-0.3,3.5) {A};
\node[font=\tiny, anchor=west] at (10.3,3.5) {B};
\node[font=\tiny, popgreen, anchor=west] at (10.15,3.0) {Ss/Sh (300 m)};
\node[font=\tiny, popblue, anchor=west] at (10.15,2.5) {Sh/Ls (250 m)};
\node[font=\tiny] at (1.2,3.35) {Sandstone};
\node[font=\tiny] at (1.2,2.72) {Shale};
\node[font=\tiny] at (1.2,2.1) {Limestone};
\node[font=\tiny, anchor=north] at (5,-0.4) {Horizontal scale 1:50,000 \quad Vertical scale 1:50,000 \quad VE = 1};
\end{tikzpicture}
\legende{Completed cross-section for horizontal strata. The contacts are horizontal lines that meet the surface exactly at the outcrop points transferred from the map; every unit keeps a constant true thickness.}
\end{popfigure}

\courssub{Worked Example 2: Uniformly Dipping Strata (VE = 1)}

\begin{exoresolu}[Cross-Section Across Dipping Beds]
\textbf{Statement:} Map scale 1:25,000. The section line X--Y, 4 km long, runs N--S, perpendicular to the strike (strike $090^{\circ}$, dip $30^{\circ}$ to the South). The ground rises uniformly from X (100 m) to Y (200 m). Two contacts cross the line: Basalt/Tuff at the point where the ground elevation is 150 m, and Tuff/Sandstone where it is 180 m. Draw the cross-section at VE = 1.
\tcblower
\textbf{Solution:}
\begin{enumerate}
  \item \textbf{Topographic profile}: horizontal scale 1 cm = 0.25 km (4 km = 16 cm on paper); vertical scale identical (VE = 1). The profile is a straight slope from (0 km, 100 m) to (4 km, 200 m).
  \item \textbf{Locate the contacts}: the ground rises 100 m over 4 km, i.e. 25 m per km.
      \begin{itemize}
        \item Basalt/Tuff at 150 m: $50$ m above X $\Rightarrow$ $50/25 = 2$ km from X. Plot P1 at (2 km, 150 m).
        \item Tuff/Sandstone at 180 m: $80$ m above X $\Rightarrow$ $80/25 = 3.2$ km from X. Plot P2 at (3.2 km, 180 m).
      \end{itemize}
  \item \textbf{Project the contacts}: the beds dip $30^{\circ}$ towards the South, i.e. towards Y. From P1 and from P2 draw straight lines descending at $30^{\circ}$ towards Y (use a protractor). The two lines must be \textbf{parallel}, since the dip is uniform.
  \item \textbf{Extend to the base}: take sea level (0 m) as the base of the section. From P1 (150 m), the vertical drop available is 150 m, so the contact reaches the base after a horizontal run $H = V/\tan\delta = 150/\tan 30^{\circ} = 150/0.577 \approx 260$ m beyond P1.
  \item \textbf{Ornament and label}: Basalt (top), Tuff (middle), Sandstone (bottom); title \emph{Section X--Y}, scales, VE = 1, dip arrow $30^{\circ}$ S.
\end{enumerate}
\textbf{Check:} because the two contacts are parallel, the Tuff unit has a constant \emph{true thickness}, measured perpendicular to the contacts. The \emph{vertical} thickness ($H\tan\delta$ style separation) is larger than the true thickness by the factor $1/\cos\delta$ — a classic examination trap.
\end{exoresolu}

\begin{popfigure}
\begin{tikzpicture}[scale=0.75]
% Axes
\draw[popline, thick] (0,0) -- (10,0);
\draw[popline, thick] (0,0) -- (0,4);
\draw[popline, thick] (10,0) -- (10,4);
\foreach \y/\e in {1/50, 2/100, 3/150, 4/200} {
  \draw (-0.08,\y) -- (0.08,\y);
  \node[font=\tiny, anchor=east] at (-0.15,\y) {\e};
}
\node[font=\tiny, anchor=east] at (-0.15,4.5) {m};
% Topography: X=100 m (y=2), Y=200 m (y=4)
\draw[poporange, very thick] (0,2) -- (10,4);
% Contact points: P1 at 2 km (x=5), 150 m (y=3); P2 at 3.2 km (x=8), 180 m (y=3.6)
\coordinate (P1) at (5,3);
\coordinate (P2) at (8,3.6);
\fill[popdark] (P1) circle (1.5pt);
\fill[popdark] (P2) circle (1.5pt);
\node[font=\tiny, anchor=south] at (P1) {P1 (150 m)};
\node[font=\tiny, anchor=south] at (P2) {P2 (180 m)};
% Projected contacts dipping 30 deg toward Y (right): slope -tan30 = -0.577
\draw[popgreen, thick] (P1) -- ++(5.2,-3.0);
\draw[popgreen, thick] (P1) -- ++(-3.46,2.0);
\draw[popblue, thick] (P2) -- ++(2.2,-1.27);
\draw[popblue, thick] (P2) -- ++(-6.93,4.0);
% Unit labels
\node[font=\tiny] at (2.2,3.6) {Basalt};
\node[font=\tiny] at (4.6,2.2) {Tuff};
\node[font=\tiny] at (2.0,0.6) {Sandstone};
% Dip arrow
\draw[->, poporange, thick] (9.2,4.4) -- ++(0.8,-0.46);
\node[font=\tiny, anchor=west] at (10.05,4.0) {dip $30^{\circ}$ S};
% Titles and scales
\node[font=\small\bfseries] at (5,4.9) {Section X--Y: uniform dip $30^{\circ}$ S (VE = 1)};
\node[font=\tiny, anchor=east] at (-0.3,2) {X (N)};
\node[font=\tiny, anchor=west] at (10.3,4) {Y (S)};
\node[font=\tiny, popgreen, anchor=west] at (10.15,0.4) {Basalt/Tuff};
\node[font=\tiny, popblue, anchor=west] at (10.15,2.6) {Tuff/Sandstone};
\node[font=\tiny, anchor=north] at (5,-0.4) {Horizontal scale 1:25,000 \quad Vertical scale 1:25,000 \quad VE = 1};
\end{tikzpicture}
\legende{Completed cross-section for uniformly dipping strata (dip $30^{\circ}$ S, VE = 1). The contacts are parallel straight lines drawn at the true dip; each unit keeps a constant true thickness measured perpendicular to the contacts.}
\end{popfigure}

\courssub{Labelling Conventions — What the Examiner Expects}

Every cross-section submitted in the GCE Advanced Level examination must carry:
\begin{itemize}
  \item \textbf{Title}: \emph{Geological cross-section along A--B} (or X--Y, etc.).
  \item \textbf{Horizontal scale}: e.g. 1:50,000 or 1 cm = 0.5 km.
  \item \textbf{Vertical scale}: e.g. 1:50,000 (VE = 1) or 1:10,000 (VE = 5). State the \textbf{VE} explicitly.
  \item \textbf{Orientation}: label the ends \cle{A} and \cle{B} with compass directions (N, S, E, W) where known.
  \item \textbf{Legend / ornamentation}: each unit filled with the \emph{same} pattern or colour as on the map, with formation symbols where given.
  \item \textbf{Dip arrows}: show the dip amount and direction on the projected contacts (e.g. $30^{\circ}$ S).
  \item \textbf{Topographic profile}: the ground surface drawn as a bold line.
  \item \textbf{Baseline}: a horizontal datum, usually sea level (0 m).
\end{itemize}

\begin{savaistu}[Cameroon Context — Sections Across the Cameroon Volcanic Line]
The \cle{Cameroon Volcanic Line} provides classic cross-section material. A section drawn perpendicular to the line across Mount Cameroon shows young, nearly horizontal \cle{basalt} flows and \cle{pyroclastic} deposits resting on an older volcanic pile, itself lying on the \cle{Precambrian basement}. Around \cle{Lake Nyos}, a section would show a crater depression filled with horizontal lake sediments overlying fractured granitic basement — the fractures through which the magmatic \ce{CO2} involved in the 1986 disaster migrated upwards. In the \cle{Benue Trough} of northern Cameroon, sections reveal folded Cretaceous sandstones and shales, commonly cut by faults — ideal practice material for Section 8.
\end{savaistu}

\courssub{Common Mistakes and How to Avoid Them}

\begin{attention}[GCE Cross-Section Checklist — Avoid These Errors]
\begin{enumerate}
  \item \textbf{Forgetting the VE}: drawing at VE = 5 but reading dip angles off the section as if they were true dips. Always state the VE; if $\mathrm{VE} \neq 1$, convert the dips before plotting.
  \item \textbf{Wrong dip direction}: projecting beds \emph{up}-dip instead of down-dip. Read the map dip symbol carefully and project \emph{towards} the dip direction on the section.
  \item \textbf{Non-parallel contacts} for uniformly dipping strata: if strike and dip are constant, the contacts must be parallel on the section. Non-parallel contacts imply changing dip or thickness — draw that only if the map evidence requires it.
  \item \textbf{Ignoring the topography}: a contact must meet the ground surface \emph{exactly} at the point where the section line crosses it on the map — not above, not below.
  \item \textbf{Careless ornamentation}: using patterns that are not in the legend, or leaving units blank. Copy the legend exactly.
  \item \textbf{Missing labels}: no title, no scales, no VE, no end labels. These are easy marks — never throw them away.
\end{enumerate}
\end{attention}

\begin{exercice}[Practice: Horizontal versus Dipping Sections]
A map at 1:50,000 shows two units: \cle{Granite} (basement) and \cle{Sandstone} (cover), separated by a planar contact that strikes $045^{\circ}$ and dips $20^{\circ}$ SE.
\begin{enumerate}
  \item Draw a cross-section along a 3 km line perpendicular to the strike (SW--NE), with the ground rising uniformly from 100 m (SW) to 200 m (NE); the contact crosses the line where the ground is at 150 m. Use VE = 1.
  \item Explain, without drawing, what a section drawn \emph{parallel} to the strike (NW--SE) would show.
\end{enumerate}
\end{exercice}

\begin{corrige}[Exercise 7.1]
\begin{enumerate}
  \item \textbf{SW--NE section (perpendicular to strike)}: the profile is a straight slope from (0 km, 100 m) to (3 km, 200 m). The contact crosses at 150 m, i.e. halfway along (1.5 km). Plot this point on the profile, then project the contact downward at $20^{\circ}$ towards the SE — which, on a SW--NE section viewed with NE to the right, is \emph{into} the section; the apparent dip in the plane of the section is found from $\tan\alpha = \tan 20^{\circ} \times \sin\beta$. Since the dip direction (SE) is at $45^{\circ}$ to the section line, $\beta = 45^{\circ}$ and $\tan\alpha = 0.364 \times 0.707 = 0.257$, so $\alpha \approx 14.4^{\circ}$ towards the NE end of the drawing. Granite lies below the contact, Sandstone above. Label the dip arrow $20^{\circ}$ SE (true dip) and note the apparent dip plotted.
  \item \textbf{NW--SE section (parallel to strike)}: the section runs \emph{along} the strike, so $\beta = 0^{\circ}$ and the apparent dip is $0^{\circ}$: the contact appears \textbf{horizontal}. The section shows Sandstone resting on Granite with a flat contact, whatever the true dip — a striking demonstration that a section parallel to strike hides all dipping structure.
\end{enumerate}
\end{corrige}

\courssub{Examination Technique — Cross-Section Questions}

\begin{aretenir}[GCE Cross-Section Protocol]
\begin{enumerate}
  \item \textbf{Read the question}: is the section line given or must you choose it? What VE is required (assume VE = 1 unless told otherwise)?
  \item \textbf{Draw the line} on the map extract and mark its ends.
  \item \textbf{Extract the data}: contour crossings $\rightarrow$ elevations; contact crossings $\rightarrow$ positions; dip symbols or the three-point problem $\rightarrow$ dip amount and direction.
  \item \textbf{Draw the profile} on graph paper with a sharp pencil, ruler and protractor.
  \item \textbf{Plot the contacts} as precise points.
  \item \textbf{Project with the protractor}: measure each dip from the horizontal, draw lightly first, then firm up the lines.
  \item \textbf{Ornament neatly} with pencil shading or stippling; never obscure the contact lines.
  \item \textbf{Label everything} — title, scales, VE, orientation, dips — before handing in.
\end{enumerate}
\end{aretenir}

\begin{experience}[Guided Practice: A Section up the Flank of Mount Cameroon]
\textbf{Materials:} an extract of the Buea 1:50,000 topographic/geological map (basalt flows, pyroclastic cones, basement gneiss), graph paper, ruler, protractor, coloured pencils.
\textbf{Task:}
\begin{enumerate}
  \item Draw a section line from the coast near Limbe (0 m) up the SW flank of Mount Cameroon towards the summit area (about 4,100 m) — a line roughly 25 km long.
  \item Choose the scales: at VE = 1 the drawing would need to be over 4 m high, so adopt a vertical scale of 1:5,000, giving $\mathrm{VE} = 10$. Construct the topographic profile.
  \item Plot the contacts between the recent basalts, the older volcanic formations and the basement gneiss.
  \item Project the contacts using the dips read from the map (gentle for the lava flows, steeper on the cone flanks), converting each true dip with $\tan\delta_{\text{plot}} = 10\tan\delta$.
  \item Discuss: does the profile have the broad, gently sloping form of a \cle{shield volcano}? Where might groundwater accumulate in such a section?
\end{enumerate}
\textbf{Teacher's note:} this exercise integrates topography, volcanology and hydrogeology — exactly the style of the \emph{Integration} questions in the GCE examination.
\end{experience}

We have now covered the \emph{principles and simple cases} of cross-section construction: horizontal strata and uniformly dipping strata, drawn at VE = 1 or with a controlled vertical exaggeration. In Section 8 we tackle the more complex structures that make geological maps fascinating — \textbf{folds, faults, unconformities and intrusions} — and learn how to project each of them correctly into the third dimension.

\coursec{Section 8: Construction of Geological Cross-Sections II — Folds, Faults, Unconformities and Intrusions}

In Section 7 we mastered the basic grammar of the geological cross-section: choosing the line of section, transferring topographic heights, plotting contacts, and completing a simple sequence of uniformly dipping beds. Real GCE maps, however, rarely show such simplicity. The examiner's favourite maps contain \cle{folds}, \cle{faults}, \cle{unconformities} and \cle{intrusions} — often two or more combined. This section gives you the tools to construct a correct section across each of these structures, and to check your finished work like a professional geologist.

\courssub{Constructing a Section Across Folded Strata}

\textbf{The intuition.} Imagine a layered cake that has been gently arched upwards. If you walk across it, the layers on one side tilt one way, and on the other side they tilt the opposite way. On a geological map, a fold is recognised by \emph{dip readings that point towards each other} (syncline) or \emph{away from each other} (anticline), and by the \emph{mirror-image repetition of outcrops} on either side of a central zone.

\begin{definition}[Axial trace and limbs]
The \cle{limbs} of a fold are the two flanks of dipping strata. The \cle{axial trace} is the line on the map along which the axial plane (the plane that divides the fold into two mirror halves) meets the ground surface. On the section, the fold axis lies where the dip directions reverse.
\end{definition}

\begin{methode}[Drawing fold profiles by projecting dips tangentially from both limbs]
\begin{enumerate}
\item Transfer the contacts and the topographic profile to your section exactly as in Section 7.
\item At each dip symbol on a limb, draw a short guide line at the measured dip angle, starting at the contact point on the surface.
\item Draw the contacts as smooth curves that are \emph{tangent} to these guide lines: the dip reading fixes the slope of the contact \emph{at the surface only}; underground, the contact curves smoothly towards the axis.
\item Work inwards from both limbs. Where the two sets of curving contacts meet, you have located the \cle{axis}. In an anticline the beds arch up; in a syncline they sag down.
\item Keep each bed at a constant true thickness, measured perpendicular to its boundaries, all the way around the fold.
\end{enumerate}
\end{methode}

\begin{attention}[Dip is a surface measurement]
A very common error is to draw fold contacts as straight lines at the measured dip all the way down. This makes beds either impossibly thick or forces them to cross. The dip symbol only fixes the \emph{tangent} at the outcrop; the contact must curve smoothly through the fold core.
\end{attention}

\courssub{Plunging Folds}

\textbf{Observation.} On many maps — and this is a classic GCE situation — the outcrops of a fold are not parallel bands but form V-shaped or U-shaped patterns that \emph{converge and close} in one direction. The fold is said to \cle{plunge}: its axis is not horizontal but tilted, like the keel of a sinking ship.

\begin{propriete}[Map expression of a plunging fold]
In a plunging anticline, contacts converge (close) in the \emph{direction of plunge}; the oldest rocks occupy the core. In a plunging syncline, contacts also converge in the direction of plunge, but the youngest rocks occupy the core. On a cross-section drawn \emph{perpendicular to the fold axis}, a plunging fold looks exactly like a non-plunging fold — the plunge only reveals itself on the map or in a longitudinal section (parallel to the axis).
\end{propriete}

Why do the contacts converge? Because the folded surface is tilted, erosion slices through it obliquely: the outcrop of each contact sweeps around the nose of the fold, producing the characteristic closure. Think of slicing a tilted, bent banana — the cut face shows a closed, curved pattern.

\begin{popfigure}
\begin{tikzpicture}[scale=0.9]
% MAP VIEW: plunging anticline, plunge North. Axial trace vertical (N-S). Section line A-B horizontal (E-W).
\begin{scope}
\draw[thick,->] (-0.2,0) -- (6.2,0) node[below left]{\small E};
\draw[thick,->] (0,-0.2) -- (0,4.4) node[below left]{\small N};
% converging contacts (plunge to the North)
\draw[thick,popblue] (0.6,4) .. controls (1.5,3.2) and (2.5,3.2) .. (3.4,4) -- (6,4);
\draw[thick,popblue] (0.6,0) -- (3.4,0) .. controls (4.3,0.8) and (4.3,3.2) .. (3.4,4);
\draw[thick,popgreen] (1.2,4) .. controls (1.8,3.4) and (2.8,3.4) .. (3.6,4) -- (6,3.8);
\draw[thick,popgreen] (1.2,0) -- (3.6,0) .. controls (4.2,0.6) and (4.2,3.4) .. (3.6,3.8);
\draw[thick,poporange] (1.8,4) .. controls (2.2,3.6) and (3.0,3.6) .. (3.8,4) -- (6,3.6);
\draw[thick,poporange] (1.8,0) -- (3.8,0) .. controls (4.1,0.4) and (4.1,3.6) .. (3.8,3.6);
% dip symbols on west limb (dip West)
\draw[thick] (1.0,3.5) -- (1.35,3.5); \draw[thick] (1.0,3.5) -- (1.0,3.2);
\node[left] at (0.65,3.35) {\scriptsize 40$^{\circ}$};
\draw[thick] (1.0,1.0) -- (1.35,1.0); \draw[thick] (1.0,1.0) -- (1.0,0.7);
\node[left] at (0.65,0.85) {\scriptsize 40$^{\circ}$};
% dip symbols on east limb (dip East)
\draw[thick] (5.0,3.5) -- (5.35,3.5); \draw[thick] (5.35,3.5) -- (5.35,3.2);
\node[right] at (5.7,3.35) {\scriptsize 40$^{\circ}$};
\draw[thick] (5.0,1.0) -- (5.35,1.0); \draw[thick] (5.35,1.0) -- (5.35,0.7);
\node[right] at (5.7,0.85) {\scriptsize 40$^{\circ}$};
% axial trace
\draw[dashed,popdark,thick] (3.5,0) -- (3.5,4) node[midway,right]{\scriptsize axial trace};
\node at (3,4.3) {\small\textbf{Map: plunging anticline (plunges N)}};
% Section line A-B (E-W)
\draw[thick,popdark] (0,2) node[left]{\scriptsize A} -- (6,2) node[right]{\scriptsize B};
\end{scope}
% SECTION (perpendicular to axis, i.e. along A-B)
\begin{scope}[yshift=-5.0cm]
\draw[thick,->] (-0.2,0) -- (6.2,0);
\draw[thick] (0,0) -- (0,2.6);
% surface
\draw[thick] (0,2.2) -- (6,2.2);
% folded contacts (symmetric arch)
\draw[thick,popblue] (0,1.2) .. controls (1.5,1.2) and (2.5,2.1) .. (3.0,2.18) .. controls (3.5,2.1) and (4.5,1.2) .. (6,1.2);
\draw[thick,popgreen] (0,0.7) .. controls (1.5,0.7) and (2.6,1.6) .. (3.0,1.68) .. controls (3.4,1.6) and (4.5,0.7) .. (6,0.7);
\draw[thick,poporange] (0,0.25) .. controls (1.5,0.25) and (2.7,1.1) .. (3.0,1.16) .. controls (3.3,1.1) and (4.5,0.25) .. (6,0.25);
% dip ticks at surface (matching map)
\draw[thick] (1.0,2.2) -- (1.0,1.75); \draw[thick] (1.0,2.2) -- (1.4,2.2);
\draw[thick] (5.0,2.2) -- (5.0,1.75); \draw[thick] (4.6,2.2) -- (5.0,2.2);
\node[below] at (3.0,0) {\small\textbf{Section along A--B (perpendicular to axis)}};
\node at (0.6,2.45) {\scriptsize A};
\node at (5.7,2.45) {\scriptsize B};
\end{scope}
\end{tikzpicture}
\legende{Map and cross-section pair for a plunging anticline (plunging North). On the map the contacts converge northward (direction of plunge) and dips point away from the axial trace; on the cross-section (perpendicular to the axis) the beds arch smoothly with constant thickness, identical to a non-plunging fold.}
\end{popfigure}

\courssub{Constructing a Section Across a Fault}

\textbf{Observation.} On a map, a fault appears as a line (usually drawn heavier than ordinary contacts) along which the outcrop pattern is suddenly \emph{offset}: a contact that you trace across the map stops at the fault and reappears, displaced, on the other side.

\begin{definition}[Downthrow side]
The \cle{downthrow side} (downthrown block) of a normal fault is the block that has moved \emph{down} relative to the other. On the section, the fault plane dips, and the same contacts are found at \emph{lower} levels on the downthrow side. The vertical displacement is the \cle{throw}.
\end{definition}

\begin{methode}[Recognising the downthrow side from repetition or omission of beds]
\begin{enumerate}
\item Identify the same contact (e.g.\ sandstone--limestone boundary) on both sides of the fault.
\item Compare its position relative to the topography: the side where the contact occurs at a \emph{lower} elevation (or where younger beds are preserved against older ones) is the downthrow side.
\item For a \emph{dipping} sequence cut by a normal fault: \cle{repetition} of beds along a traverse perpendicular to the fault indicates a fault whose dip opposes the bed dip; \cle{omission} of beds indicates a fault dipping in the same direction as the beds. On a map with flat topography, simply note which beds are cut out or duplicated at the fault trace.
\item Draw the fault plane on the section at its measured dip (if no dip is given, GCE convention allows a steep plane, about $60^{\circ}$--$70^{\circ}$), then draw the contacts on each side so that they match the surface outcrops and show the correct vertical displacement (\cle{throw}).
\end{enumerate}
\end{methode}

\begin{attention}[A fault offsets contacts even on flat ground]
Students trained on the ``rule of V's'' sometimes think offsets only happen in valleys. Wrong: even with perfectly flat topography, a fault \emph{shifts} the contacts horizontally on the map. Never draw the contacts continuing smoothly across a fault trace — the shift is the very signature of the fault.
\end{attention}

\begin{popfigure}
\begin{tikzpicture}[scale=0.9]
% MAP
\begin{scope}
\draw[thick] (0,0) rectangle (6,4);
\draw[thick,popblue] (0,1.0) -- (2.8,1.0); \draw[thick,popblue] (3.2,1.6) -- (6,1.6);
\draw[thick,popgreen] (0,2.2) -- (2.8,2.2); \draw[thick,popgreen] (3.2,2.8) -- (6,2.8);
\draw[thick,poporange] (0,3.4) -- (2.8,3.4); \draw[thick,poporange] (3.2,4.0) -- (6,4.0);
\draw[very thick,popdark] (2.8,-0.1) -- (3.2,4.1) node[above]{\scriptsize fault};
\node at (1.2,0.5) {\scriptsize shale};
\node at (1.2,1.6) {\scriptsize sst};
\node at (1.2,2.8) {\scriptsize lst};
\node at (4.8,1.15) {\scriptsize shale};
\node at (4.8,2.2) {\scriptsize sst};
\node at (4.8,3.4) {\scriptsize lst};
\node at (3,4.5) {\small\textbf{Map: faulted dipping sequence}};
\end{scope}
% SECTION
\begin{scope}[yshift=-4.8cm]
\draw[thick] (0,0) -- (6,0); \draw[thick] (0,0) -- (0,3);
\draw[thick] (0,2.6) -- (6,2.6);
% fault plane
\draw[very thick,popdark] (2.9,2.6) -- (3.5,0);
% left block contacts (upthrown)
\draw[thick,popblue] (0,1.9) -- (2.9,1.9);
\draw[thick,popgreen] (0,1.1) -- (2.97,1.1);
\draw[thick,poporange] (0,0.3) -- (3.03,0.3);
% right block contacts (downthrown)
\draw[thick,popblue] (3.18,1.3) -- (6,1.3);
\draw[thick,popgreen] (3.3,0.5) -- (6,0.5);
\node at (1.4,2.3) {\scriptsize upthrown block};
\node at (4.8,2.3) {\scriptsize downthrown block};
\node at (0.7,2.05) {\scriptsize lst};
\node at (0.7,1.5) {\scriptsize sst};
\node at (0.7,0.7) {\scriptsize shale};
\node at (5.3,1.5) {\scriptsize lst};
\node at (5.3,0.9) {\scriptsize sst};
\node at (3,-0.5) {\small\textbf{Section: normal fault, throw to the east}};
\end{scope}
\end{tikzpicture}
\legende{Map and section across a normal fault. On the map the contacts are offset at the fault trace; on the section the same contacts lie lower on the downthrown (eastern) block.}
\end{popfigure}

\begin{exoresolu}[Faulted sequence: finding the downthrow side]
On a flat map, a north--south fault cuts a sequence that dips $30^{\circ}$ east. West of the fault, a traverse from west to east crosses: shale, sandstone, limestone. East of the fault, the same traverse crosses: sandstone, limestone, conglomerate. Which side is downthrown?
\tcblower
\textbf{Solution.} The normal upward order of the sequence is shale $\rightarrow$ sandstone $\rightarrow$ limestone $\rightarrow$ conglomerate (oldest to youngest, since the beds dip and young eastward). West of the fault, the \emph{oldest} unit (shale) is exposed; east of the fault, the shale is missing and the \emph{youngest} unit (conglomerate) appears. Younger beds preserved against older beds means the eastern block has moved \emph{down}, bringing younger strata to the same erosion level. \textbf{The east side is downthrown.} On the section, draw the fault plane steeply, with every contact lower on the east side by the same vertical amount (the throw).
\end{exoresolu}

\courssub{Unconformities on Section}

\textbf{Observation.} On the map, an \cle{angular unconformity} is seen where one set of parallel outcrop bands is cut across by the outcrops of an overlying series with a \emph{different} strike. The boundary between them is the unconformity surface.

\begin{definition}[Angular unconformity on a section]
An angular unconformity is drawn as an uneven (often gently undulating) surface. Below it, the older beds are \cle{truncated}: their contacts end abruptly against the surface. Above it, the younger beds \emph{rest on} the surface and may \cle{overstep} — each successively younger bed overlapping further across the buried landscape, resting directly on progressively older basement.
\end{definition}

\textbf{Construction rules:}
\begin{itemize}
\item Draw the unconformity surface first, following its outcrop on the topographic profile.
\item Below: draw the older (often steeply dipping or folded) beds and \emph{stop them at the surface} — never continue them above it.
\item Above: draw the younger series (usually gently dipping or horizontal) resting on the surface; the lowest bed of the cover is typically a \cle{basal conglomerate}.
\end{itemize}

\courssub{Intrusions on Section}

\begin{definition}[Dyke, sill, batholith]
A \cle{dyke} is a sheet-like intrusion that \emph{cuts across} the bedding (discordant). A \cle{sill} is a sheet-like intrusion that \emph{follows} the bedding (concordant), injected between two layers. A \cle{batholith} is a very large, irregular granitic body; on a section its margins cut across country-rock structures and typically widen downwards.
\end{definition}

\begin{attention}[Sills follow bedding; dykes cut it]
Before drawing any intrusion, check the map: if the igneous outcrop runs \emph{parallel} to the contacts of the sedimentary beds, it is probably a sill — draw it sandwiched between beds, following every fold and tilt. If it slices \emph{across} the contacts (often as a straight, steep band), it is a dyke — draw it steep, crossing beds and even the unconformity if its outcrop does so. A sill can never cross an unconformity into the basement; a dyke can.
\end{attention}

\begin{popfigure}
\begin{tikzpicture}[scale=1.0]
\draw[thick] (0,0) -- (8,0); \draw[thick] (0,0) -- (0,4.2);
% surface
\draw[thick] (0,3.6) -- (8,3.6);
% basement tilted beds (truncated)
\draw[thick,popblue] (0,0.6) -- (2.1,3.1);
\draw[thick,popgreen] (0.6,0) -- (3.1,3.15);
\draw[thick,poporange] (1.7,0) -- (4.1,3.2);
\draw[thick,poppurple] (2.8,0) -- (5.1,3.25);
% unconformity surface
\draw[very thick,popdark] (0,3.1) .. controls (2,3.0) and (4,3.35) .. (6,3.2) .. controls (7,3.15) and (7.6,3.25) .. (8,3.2);
% cover beds
\draw[thick,popgold] (0,3.3) .. controls (2,3.2) and (4,3.55) .. (6,3.4) .. controls (7,3.35) and (7.6,3.45) .. (8,3.4);
% dyke cutting basement and unconformity, stopping in cover
\draw[very thick,popink,fill=popink!15] (5.6,0) -- (6.0,0) -- (5.5,3.45) -- (5.15,3.4) -- cycle;
\node at (1.1,3.85) {\scriptsize cover (conglomerate)};
\node at (2.2,1.4) {\scriptsize truncated basement beds};
\node at (7.3,3.0) {\scriptsize unconformity};
\node[rotate=78] at (5.55,1.7) {\scriptsize dyke};
\node at (4,-0.5) {\small\textbf{Angular unconformity with a cross-cutting dyke}};
\end{tikzpicture}
\legende{Section showing tilted basement beds truncated by an angular unconformity, overlain by a gently undulating cover series. A dyke cuts across the basement bedding and the unconformity, so it is younger than both.}
\end{popfigure}

\courssub{Combined Structures: A Full Worked Section}

GCE examiners love compound maps. Let us work through the logic of the classic combination: \emph{a folded series, cut by a fault, then planed off and overlain by an unconformable horizontal series}.

\begin{exoresolu}[Compound section: fold + fault + unconformity]
A flat map shows, from west to east: bands of mudstone--sandstone--limestone--sandstone--mudstone (dips of $45^{\circ}$ pointing \emph{towards} the centre of the pattern), all truncated by a N--S fault; east of the fault the same bands reappear, shifted northwards; the whole assemblage is overlain in the south of the map by horizontal beds of conglomerate then shale resting on a surface that cuts across all the older outcrops. Describe how to construct the section along an E--W line.
\tcblower
\textbf{Solution, step by step.}
\begin{enumerate}
\item \textbf{Topography:} flat profile at the map's elevation.
\item \textbf{Fold:} dips point towards each other and the limestone occupies the centre with the youngest position in the sequence — a \emph{syncline}. Draw the contacts as smooth U-shaped curves, tangent to the $45^{\circ}$ dips at the surface, keeping bed thickness constant.
\item \textbf{Fault:} draw the fault plane steeply (say $70^{\circ}$). The map shows the fold pattern shifted, so offset every folded contact at the fault by the same throw; the side where younger beds (mudstone margin) are preserved against older is downthrown.
\item \textbf{Unconformity:} draw a horizontal surface above everything. \emph{Stop} the folded beds \emph{and the fault} at this surface — the fault does not displace the cover, so it is older than the unconformity.
\item \textbf{Cover:} draw the conglomerate then shale as horizontal beds above the surface; the conglomerate rests on different older units along the section (overstep).
\item \textbf{History check:} deposition and folding of the old series $\rightarrow$ faulting $\rightarrow$ erosion (unconformity) $\rightarrow$ deposition of the cover. Every line on your section must agree with this order.
\end{enumerate}
\end{exoresolu}

\begin{savaistu}[Sections and the Cameroon Volcanic Line]
Geologists mapping the Cameroon Volcanic Line — the chain of volcanoes running from Mount Cameroon through the Manengouba, Bambouto and Oku massifs — rely on exactly these skills. Cross-sections drawn across the volcanic centres show lava flows and pyroclastic beds resting unconformably on much older granitic and gneissic basement, pierced by dykes that fed the eruptions. Reading such sections tells us the order of events over tens of millions of years.
\end{savaistu}

\courssub{Checking a Finished Section and Common Errors}

\begin{methode}[The five-point check of a finished section]
\begin{enumerate}
\item \textbf{Surface agreement:} every contact, fault, dyke and dip symbol on the line of section must appear on your profile at the correct position.
\item \textbf{Thickness consistency:} each bed keeps the same true thickness along the whole section, around folds and across the profile.
\item \textbf{Symbol agreement:} dips drawn on the section match the dip values and directions shown on the map; contacts never cross each other.
\item \textbf{Structural logic:} faults offset \emph{all} older features but stop at unconformities; dykes cut everything they cross; sills never leave their bedding horizon.
\item \textbf{Labelling:} title, line of section (A--B), horizontal and vertical scales, and a key to the ornament used for each rock unit.
\end{enumerate}
\end{methode}

\begin{attention}[Frequent construction errors]
\begin{itemize}
\item Continuing folded contacts as straight lines until they cross — contacts must curve and never intersect.
\item Forgetting that a fault offsets \emph{every} contact, including the axial trace of a fold.
\item Letting a fault cut through an unconformity when the map shows the cover undisturbed.
\item Drawing a sill crossing from one bed into another, or a dyke bending to follow bedding.
\item Changing bed thickness from one side of a fold to the other.
\end{itemize}
\end{attention}

\begin{exercice}[Practice: compound section]
On a flat map at elevation \SI{200}{m}, an E--W line of section crosses: a band of schist (vertical contact with granite to the west), a steep dyke cutting the schist, then an unconformity surface above which lie horizontal beds of basal conglomerate and sandstone. Draw the section, state the relative ages of the granite, schist, dyke and cover, and justify each step.
\end{exercice}

\begin{corrige}[Practice: compound section]
\textbf{Section:} Draw a horizontal topographic line at 200 m. From west to east: (1) Granite with a vertical contact against schist. (2) Schist with steeply dipping foliation/bedding. (3) A steep dyke (near vertical) cutting the schist, terminating at the unconformity. (4) An unconformity surface (horizontal line) truncating the schist, dyke, and granite contact. (5) Horizontal basal conglomerate resting on the unconformity, overlain by horizontal sandstone.

\textbf{Relative ages (oldest to youngest) with justification:}
\begin{enumerate}
\item \textbf{Schist:} Country rock, cut by both the granite contact and the dyke.
\item \textbf{Granite:} Intrudes the schist (vertical intrusive contact), therefore younger than schist.
\item \textbf{Dyke:} Cuts the schist (and possibly the granite if the dyke crosses the contact; the map says "dyke cutting the schist", so at minimum younger than schist). The dyke is truncated by the unconformity, so it is older than the unconformity.
\item \textbf{Unconformity:} Erosion surface cutting all older rocks.
\item \textbf{Conglomerate and Sandstone:} Deposited on the unconformity, youngest.
\end{enumerate}
\textbf{Order:} Schist $\rightarrow$ Granite $\rightarrow$ Dyke $\rightarrow$ Unconformity $\rightarrow$ Conglomerate $\rightarrow$ Sandstone.
\end{corrige}

\begin{aretenir}[Key points of Section 8]
\begin{itemize}
\item Draw fold contacts as smooth curves tangent to the measured dips; locate the axis where dip directions reverse; keep bed thickness constant.
\item Plunging folds show converging contacts on the map in the direction of plunge; anticlines have oldest rocks in the core, synclines the youngest. Cross-sections perpendicular to the axis show the true fold profile.
\item Faults offset contacts on map and section; the downthrow side is where younger beds are preserved or contacts lie lower. Use repetition/omission criteria for dipping sequences.
\item Unconformities truncate beds below and carry overstepping beds above; faults and dykes stop at the unconformity if the cover is undisturbed.
\item Dykes cut bedding; sills follow it; batholith margins cross-cut and widen downwards.
\item Always finish with the five-point check: surface agreement, thickness, symbols, structural logic, labelling.
\end{itemize}
\end{aretenir}

\coursec{Section 9: Description of Geological Structures on Maps I — Strata, Dip and Outcrop Patterns}

In Sections 7 and 8 you learnt to \emph{construct} geological cross-sections, first for simple dipping strata, then for folded, faulted, unconformable and intruded terrains. Every construction began in the same way: by \emph{reading} the map carefully. We now formalise that reading. The GCE examiner does not only ask you to draw sections; very often the first question on a map paper is simply: \emph{``Describe the geology of the area shown on the map.''} This section teaches you how to describe \cle{strata}, \cle{dip} and \cle{outcrop patterns} in a rigorous, ordered and mark-earning way. Section 10 will extend the method to folds, faults, unconformities, intrusions and the full geological history.

\courssub{A Framework for Describing Any Geological Map}

Imagine you are flying over the region in a small aircraft. What would you notice first? The \emph{shape of the land} (hills, valleys, plateaus). Then, looking closer, the \emph{rocks} exposed at the surface. Then the \emph{way the rock layers lie} (flat, tilted, vertical). Then the \emph{structures} that deform them (folds, faults). Finally you would ask \emph{how all this came about} — the history. A good map description follows exactly this natural order, from the most obvious to the most interpretative.

\begin{methode}[The standard five-point plan for describing a geological map]
\begin{enumerate}
\item \textbf{Relief and drainage.} State the highest and lowest points, the direction of slopes, the positions of valleys, ridges, scarps and rivers, using the contour lines and spot heights.
\item \textbf{Rock units.} List the rock types (or numbered/lettered units) present, using the legend, and state where each outcrops on the map.
\item \textbf{Attitude of the strata.} State whether the beds are horizontal, vertical or dipping; give the direction and amount of dip (from symbols or from outcrop patterns).
\item \textbf{Structures.} Describe folds, faults, unconformities and intrusions, with their locations, orientations and effects on the outcrops.
\item \textbf{Geological history.} Reconstruct the order of events from the oldest to the youngest, justifying each step (superposition, cross-cutting relations, unconformities).
\end{enumerate}
\end{methode}

\begin{attention}[Do not describe structures before stating the rock units]
A very common mistake is to jump straight to ``there is a fault in the north-east'' before saying what rocks are present. The examiner cannot credit a structure whose affected rocks you have never named. Similarly, never state a dip direction without first establishing that the beds dip at all. Follow the five-point plan \emph{in order}; each step depends on the previous one.
\end{attention}

\begin{attention}[Exam technique: description is marked point by point]
Description questions are marked with a point-by-point scheme: one mark per correct, relevant, factual statement. Therefore:
\begin{itemize}
\item Write \textbf{short, factual sentences}, one idea per sentence (``The beds dip $30^\circ$ towards the east.'' earns a mark; a long vague paragraph earns nothing).
\item Always give \textbf{locations} (north, south-west corner, along the river\ldots) and \textbf{values} (dip angles, heights, thicknesses) whenever the map allows.
\item Never invent information not shown on the map, and never repeat the same idea in different words hoping for two marks.
\end{itemize}
\end{attention}

\courssub{Describing Horizontal Strata}

\begin{definition}[Horizontal strata]
Strata are \cle{horizontal} when their bounding surfaces (the \cle{contacts}) are horizontal planes, i.e.\ every point on a given contact lies at the same altitude. Their dip is $0^\circ$.
\end{definition}

On a map, this simple geometry has a spectacular consequence. A horizontal contact is a surface of constant altitude — exactly like a contour line, which also joins points of equal altitude. Wherever the land surface cuts the contact, the outcrop of the contact must therefore follow the contour of that altitude.

\begin{propriete}[Contacts of horizontal beds follow contour lines]
The outcrop trace of the contact between two horizontal beds is everywhere \textbf{parallel to the contour line} whose altitude equals that of the contact. Consequently, the outcrop pattern of horizontal strata \emph{mimics the topography}: contacts bend into valleys and around hills exactly as the contours do.
\end{propriete}

This property is not examinable as a proof, but you must be able to \emph{use} it in both directions: if contacts follow contours, the beds are horizontal; if beds are horizontal, you can predict their outcrops from the contours.

How to describe horizontal strata on a map:
\begin{itemize}
\item State that the contacts are \textbf{parallel to contour lines}, and conclude that the beds are \textbf{horizontal} (dip $0^\circ$).
\item Note that each unit forms \textbf{irregular, sinuous outcrops} that follow the relief: the oldest unit outcrops in the \textbf{valley bottoms} (lowest ground), the youngest caps the \textbf{hill tops}.
\item State that the \textbf{thickness} of each unit is constant and equals the difference in altitude between its lower and upper contacts (readable from the contours).
\end{itemize}

\begin{exemplebox}[Horizontal strata in a valley]
On a map, the contact between a shale (below) and a sandstone (above) follows the \SI{400}{m} contour into every valley and around every spur. The contact between the sandstone and an overlying limestone follows the \SI{500}{m} contour. Description: ``The beds are horizontal. The shale outcrops in the valley floors below \SI{400}{m}, the sandstone forms a band between \SI{400}{m} and \SI{500}{m} on the hill flanks, and the limestone caps the hills above \SI{500}{m}. The sandstone is \SI{100}{m} thick.'' This is exactly the situation of the sandstone and limestone plateaus around parts of the Adamawa and Western Highlands of Cameroon, where flat-lying beds cap the hills.
\end{exemplebox}

\begin{savaistu}[Horizontal beds in Cameroon]
The Cretaceous sandstones of the Benue Trough (e.g., around Garoua) and the basaltic plateaus of the Adamawa often display near-horizontal attitudes. Their outcrop patterns on geological maps closely follow the topographic contours, producing the characteristic ``layer-cake'' topography of cuestas and mesas.
\end{savaistu}

\courssub{Describing Vertical Strata}

\begin{definition}[Vertical strata]
Strata are \cle{vertical} when their contacts are vertical planes; their dip is $90^\circ$.
\end{definition}

A vertical plane cuts the land surface along a \textbf{straight line}, whatever the relief, because the intersection of a vertical plane with the ground depends only on the horizontal position, not on the altitude. Hence:

\begin{propriete}[Outcrops of vertical beds are straight]
The contacts of vertical strata appear on the map as \textbf{straight lines} that cross valleys, ridges and contour lines without deflecting. Their outcrop width is the \emph{true thickness} of the bed, and it is constant along strike.
\end{propriete}

How to describe vertical strata: state that the contacts are \textbf{straight and independent of the relief}, conclude that the beds are \textbf{vertical} (dip $90^\circ$), give the \textbf{strike direction} of the bands (e.g.\ ``the bands trend N--S'') and their \textbf{true thickness} measured directly on the map using the scale. Note that for vertical beds, \emph{no younging direction can be deduced from the outcrop pattern alone} — you need the legend or sedimentary structures (cross-bedding, graded bedding) or facing indicators.

\begin{savaistu}[Vertical beds in Cameroon]
Vertical or near-vertical strata are common in the intensely deformed Precambrian basement of the Congo Craton (e.g., the Ntem Complex in southern Cameroon) and in dyke swarms associated with the Cameroon Volcanic Line. On maps, these appear as straight, parallel bands cutting across contours.
\end{savaistu}

\courssub{Describing Dipping Strata}

Most real strata are inclined. Three kinds of map evidence reveal a dip: the \textbf{V-rule pattern} of contacts crossing valleys, the \textbf{dip symbols}, and the \textbf{asymmetry of the landscape} (scarp and dip slopes).

\textbf{1. The V-rule.} When a dipping contact crosses a valley, its outcrop trace makes a ``V'' whose shape and direction depend on the dip:

\begin{center}
\begin{tabularx}{\linewidth}{|l|X|X|}
\hline
\rowcolor{popblueL} \textbf{Attitude} & \textbf{Pattern in the valley} & \textbf{Conclusion} \\
\hline
Horizontal beds & Contact follows contour; no distinct V & Dip $= 0^\circ$ \\
\hline
Vertical beds & No V; contact is straight & Dip $= 90^\circ$ \\
\hline
Dipping downstream, dip $<$ valley gradient & V points \emph{upstream} & Gentle downstream dip \\
\hline
Dipping downstream, dip $>$ valley gradient & V points \emph{downstream} & Steep downstream dip \\
\hline
Dipping upstream & V points \emph{downstream} & Upstream dip \\
\hline
\end{tabularx}
\end{center}

A useful memory aid: \emph{``the V tells you which way the water and the beds disagree.''} In GCE maps the two cases you will meet most are: V pointing \textbf{upstream} for gently dipping beds (the commonest), and \textbf{no V at all} for vertical beds.

\textbf{2. Dip symbols.} The conventional symbol is a T-shape: the long bar gives the \cle{strike}, the short tick points in the \cle{dip direction}, and the number beside it is the \cle{dip angle} in degrees. When symbols are present, simply report them: ``The beds dip $25^\circ$ towards the south-east.''

\textbf{3. Scarp and dip-slope asymmetry.} Dipping resistant beds produce an asymmetric ridge (a \cle{cuesta}): a steep \cle{scarp slope} cutting across the beds, and a gentle \cle{dip slope} following the top of the resistant bed. On the map, the contours are \emph{closely spaced} on the scarp side and \emph{widely spaced} on the dip-slope side. The gentle slope indicates the \textbf{dip direction}. This is how dip can be recognised even without symbols.

\begin{propriete}[Determining dip from map evidence alone]
Without dip symbols, the dip direction and amount can be obtained from the map by:
\begin{enumerate}
\item Drawing \textbf{strike lines} (structure contours): join points where the same contact crosses the same contour value. The strike lines are parallel to the strike.
\item The \textbf{dip direction} is perpendicular to the strike lines, towards \emph{decreasing} altitude of the contact.
\item The \textbf{amount of dip} $\theta$ is given by
\[ \tan\theta = \frac{\text{vertical interval between strike lines}}{\text{horizontal spacing of the strike lines (true distance)}}. \]
\end{enumerate}
This is the three-point-problem technique of Section 6 applied to map description. (The derivation follows directly from the definition of dip as the angle of maximum slope of the plane.)
\end{propriete}

\begin{exoresolu}[Reading a dip from outcrop pattern]
On a $1:25\,000$ map, the contact between a sandstone and an overlying shale crosses the \SI{300}{m} contour at point $A$ and the \SI{200}{m} contour at point $B$. $A$ and $B$ lie on the same line perpendicular to the strike, \SI{2}{cm} apart on the map. Determine the dip.
\tcblower
\textbf{Solution.}
\begin{enumerate}
\item The contact drops from \SI{300}{m} to \SI{200}{m} between $A$ and $B$: vertical interval $h = 100\ \mathrm{m}$.
\item True horizontal distance: $d = 2\ \mathrm{cm} \times 25\,000 = 50\,000\ \mathrm{cm} = 500\ \mathrm{m}$.
\item $\tan\theta = \dfrac{h}{d} = \dfrac{100}{500} = 0.2$, so $\theta = \arctan(0.2) \approx 11.3^\circ$.
\end{enumerate}
The beds dip approximately $11^\circ$ in the direction from $A$ towards $B$ (the direction in which the contact falls in altitude).
\end{exoresolu}

\begin{exoresolu}[Three-point problem on a map]
Three outcrops of the same limestone bed are found at: $P_1$ (easting 420, northing 150, altitude \SI{320}{m}), $P_2$ (420, 152, \SI{280}{m}), $P_3$ (422, 151, \SI{300}{m}). Map scale $1:50\,000$. Determine the strike and dip of the bed.
\tcblower
\textbf{Solution.}
\begin{enumerate}
\item Plot the three points. The bed passes through all three. The strike line is the line joining the two points of equal altitude. Here, no two have the same altitude, so we interpolate: find point $Q$ on $P_1P_2$ at \SI{300}{m}. $P_1$ (320) to $P_2$ (280) is 40 m difference over 2 km (2000 m). To drop 20 m from $P_1$, go halfway: $Q$ is at (420, 151, 300).
\item Strike line joins $Q$ (420, 151) and $P_3$ (422, 151). This line is exactly E--W. So \textbf{strike = E--W} (or N90°E).
\item Dip direction is perpendicular to strike, towards lower altitude. Altitude decreases from 300 m at $Q$ and $P_3$ to 280 m at $P_2$ (south). So \textbf{dip direction = South}.
\item Horizontal distance between strike line (through $Q$ and $P_3$) and the 280 m structure contour (through $P_2$) is 1 km (from northing 151 to 152). Vertical interval = 20 m.
\item $\tan\theta = 20/1000 = 0.02$, $\theta \approx 1.1^\circ$. \textbf{Dip = $1^\circ$ South}.
\end{enumerate}
\end{exoresolu}

\begin{attention}[Common traps with dipping strata]
\begin{itemize}
\item Do not confuse \textbf{dip direction} (downhill of the bed) with \textbf{strike} (along the bed); they are always perpendicular.
\item A V pointing upstream does \emph{not} always mean the beds dip upstream — gently dipping beds that dip downstream also make an upstream V. Check the dip symbol or the scarp/dip-slope asymmetry before concluding.
\item Outcrop \textbf{width} of a dipping bed is \emph{not} its true thickness: width increases as the ground slope decreases and as the dip decreases. Never quote outcrop width as thickness.
\item True thickness $t = w \sin\theta$ where $w$ is the map width measured perpendicular to strike. Only for vertical beds ($\theta=90^\circ$) does $w = t$.
\end{itemize}
\end{attention}

\begin{popfigure}
\begin{tikzpicture}[scale=0.9]
% ---- Map 1: horizontal ----
\begin{scope}
\draw[popline] (0,0) rectangle (4,3);
% contours
\draw[gray] (0,0.5) .. controls (1.5,1.1) and (2.5,0.2) .. (4,0.8);
\draw[gray] (0,1.5) .. controls (1.5,2.1) and (2.5,1.2) .. (4,1.8);
\draw[gray] (0,2.5) .. controls (1.5,3.1) and (2.5,2.2) .. (4,2.8);
% contacts parallel to contours
\draw[very thick,popblue] (0,0.5) .. controls (1.5,1.1) and (2.5,0.2) .. (4,0.8);
\draw[very thick,poporange] (0,1.5) .. controls (1.5,2.1) and (2.5,1.2) .. (4,1.8);
% river
\draw[thick,popblue!60] (2,3) -- (2,0);
\node[align=center,font=\small] at (2,-0.5) {Horizontal beds};
\node[align=center,font=\scriptsize] at (2,-1.0) {contacts follow contours};
\end{scope}
% ---- Map 2: vertical ----
\begin{scope}[xshift=5.2cm]
\draw[popline] (0,0) rectangle (4,3);
\draw[gray] (0,0.5) .. controls (1.5,1.1) and (2.5,0.2) .. (4,0.8);
\draw[gray] (0,1.5) .. controls (1.5,2.1) and (2.5,1.2) .. (4,1.8);
\draw[gray] (0,2.5) .. controls (1.5,3.1) and (2.5,2.2) .. (4,2.8);
\draw[very thick,popblue] (1.3,0) -- (1.3,3);
\draw[very thick,poporange] (2.7,0) -- (2.7,3);
\draw[thick,popblue!60] (2,3) -- (2,0);
\node[align=center,font=\small] at (2,-0.5) {Vertical beds};
\node[align=center,font=\scriptsize] at (2,-1.0) {straight contacts cross relief};
\end{scope}
% ---- Map 3: dipping ----
\begin{scope}[xshift=10.4cm]
\draw[popline] (0,0) rectangle (4,3);
\draw[gray] (0,0.5) .. controls (1.5,1.1) and (2.5,0.2) .. (4,0.8);
\draw[gray] (0,1.5) .. controls (1.5,2.1) and (2.5,1.2) .. (4,1.8);
\draw[gray] (0,2.5) .. controls (1.5,3.1) and (2.5,2.2) .. (4,2.8);
% dipping contacts with V at river
\draw[very thick,popblue] (0,0.3) -- (1.7,0.9) -- (2,0.55) -- (2.3,0.9) -- (4,1.5);
\draw[very thick,poporange] (0,1.3) -- (1.7,1.9) -- (2,1.55) -- (2.3,1.9) -- (4,2.5);
\draw[thick,popblue!60] (2,3) -- (2,0);
\node[align=center,font=\small] at (2,-0.5) {Dipping beds};
\node[align=center,font=\scriptsize] at (2,-1.0) {contacts make V's in the valley};
\end{scope}
\end{tikzpicture}
\legende{The same valley (blue stream, grey contours) crossed by horizontal, vertical and dipping strata. Horizontal contacts mimic the contours; vertical contacts are straight; dipping contacts deflect into V's where they cross the valley.}
\end{popfigure}

\courssub{Stratigraphic Order and Relative Ages}

The description of strata is incomplete without their \textbf{order of succession} and \textbf{relative ages}. Two sources of evidence are used together:

\begin{itemize}
\item \textbf{The legend.} On a geological map the legend (key) is conventionally arranged with the \textbf{youngest unit at the top and the oldest at the bottom} (some maps use the reverse — always check for an arrow or age labels). The legend also gives the name or age of each unit (e.g.\ Cretaceous sandstone, Precambrian gneiss).
\item \textbf{Superposition on the map.} For horizontal beds, the unit capping the hills is the youngest, the unit in the valley floors the oldest. For dipping beds, the beds become younger \emph{in the dip direction} (provided the succession is not overturned). For vertical beds, the order must be read from the legend alone.
\end{itemize}

State the succession explicitly, e.g.: ``From oldest to youngest the units are: mica schist, quartzite, sandstone, shale, limestone.'' This single sentence often carries several marks.

\begin{attention}[Overturned strata]
If sedimentary structures (cross-bedding, ripple marks) or fossil evidence show that younging is \emph{opposite} to the dip direction, the strata are \cle{overturned}. In that case, the order of superposition on the map is reversed: the oldest beds are in the direction of dip. Always check the legend for overturned symbols or explicit notes.
\end{attention}

\courssub{Guided Description of a Worked Map Extract}

\begin{popfigure}
\begin{tikzpicture}[scale=1.0]
% map frame
\draw[popline] (0,0) rectangle (9,6);
% contours (hill in SW, valley draining NE)
\draw[gray] (0,1.2) .. controls (2,2.0) and (3,0.6) .. (5,1.4) .. controls (7,2.2) and (8,1.6) .. (9,2.2);
\draw[gray] (0,2.6) .. controls (2,3.4) and (3,2.0) .. (5,2.8) .. controls (7,3.6) and (8,3.0) .. (9,3.6);
\draw[gray] (0,4.0) .. controls (2,4.8) and (3,3.4) .. (5,4.2) .. controls (7,5.0) and (8,4.4) .. (9,5.0);
\node[gray,font=\scriptsize] at (8.5,2.45) {200};
\node[gray,font=\scriptsize] at (8.5,3.85) {300};
\node[gray,font=\scriptsize] at (8.5,5.25) {400};
% river
\draw[thick,popblue!60] (4.5,6) .. controls (4.2,4) and (4.8,2) .. (4.5,0);
\node[popblue!60,font=\scriptsize] at (4.9,0.35) {R.};
% horizontal contact following 300 contour (west part)
\draw[very thick,poporange] (0,2.6) .. controls (2,3.4) and (3,2.0) .. (4.5,2.8);
% dipping contacts (east part) with V at river
\draw[very thick,popgreen] (4.5,2.8) -- (4.15,3.3) -- (4.5,3.8) -- (9,5.6);
\draw[very thick,popgreen] (4.5,0.9) -- (4.15,1.4) -- (4.5,1.9) -- (9,3.7);
% dip symbol
\draw[very thick,popdark] (7,4.9) -- (7.8,4.9);
\draw[very thick,popdark] (7.4,4.9) -- (7.4,4.55);
\node[popdark,font=\scriptsize] at (7.4,5.2) {20};
% units labels
\node[font=\scriptsize] at (1.2,5.2) {Unit C (hilltop)};
\node[font=\scriptsize] at (1.2,1.9) {Unit B};
\node[font=\scriptsize] at (1.2,0.5) {Unit A (valley)};
\node[font=\scriptsize] at (6.6,2.6) {Unit D};
\node[font=\scriptsize] at (7.6,4.3) {Unit E};
% north arrow
\draw[->,thick,popdark] (0.5,5.4) -- (0.5,5.95);
\node[font=\scriptsize,popdark] at (0.5,6.15) {N};
% numbered arrows showing description order
\draw[->,very thick,poppink] (8.6,5.7) -- (8.6,5.3);
\node[poppink,font=\scriptsize] at (8.85,5.9) {1};
\draw[->,very thick,poppink] (2.0,4.4) -- (1.6,4.0);
\node[poppink,font=\scriptsize] at (2.25,4.6) {2};
\draw[->,very thick,poppink] (6.2,5.6) -- (6.9,5.15);
\node[poppink,font=\scriptsize] at (5.9,5.75) {3};
\draw[->,very thick,poppink] (4.0,1.0) -- (4.35,1.35);
\node[poppink,font=\scriptsize] at (3.7,0.85) {4};
\end{tikzpicture}
\legende{Annotated map extract. The pink arrows recall the description order: \textbf{1} relief (contours, river), \textbf{2} rock units and their outcrops, \textbf{3} attitude of the strata (contact patterns, dip symbol), \textbf{4} structures and history (here: the change from horizontal to dipping behaviour across the map).}
\end{popfigure}

\begin{experience}[Guided description of the map extract above]
Apply the five-point plan to the figure:
\begin{enumerate}
\item \textbf{Relief.} The land rises from about \SI{200}{m} in the east to over \SI{400}{m} in the north-west; a river flows southwards through the centre of the map in a valley.
\item \textbf{Rock units.} Five units outcrop: A, B and C in the west, D and E in the east. Unit A occupies the valley floor, B the lower slopes, C the hilltops; D and E form bands trending roughly NE--SW.
\item \textbf{Attitude.} In the west, the contact between B and C follows the \SI{300}{m} contour: units A--C are \textbf{horizontal}. In the east, the contacts between D and E deflect into V's where they cross the river, and a symbol gives a dip of $20^\circ$; the strike lines fall in altitude towards the south-east, so units D--E \textbf{dip $20^\circ$ to the south-east}.
\item \textbf{Structures.} The contrast between horizontal western units and dipping eastern units suggests an \textbf{unconformity} or a structural boundary near the centre of the map (developed fully in Section 10).
\item \textbf{History (sketch).} Deposition of A, B, C horizontally; deposition (or tilting) of D--E; tilting of D--E to $20^\circ$; erosion to the present surface; excavation of the modern valley.
\end{enumerate}
\end{experience}

\begin{exoresolu}[Describing strata on a map extract]
On a $1:50\,000$ map, three units outcrop: a conglomerate, a sandstone and a shale. The conglomerate--sandstone contact follows the \SI{500}{m} contour around all the hills; the sandstone--shale contact follows the \SI{600}{m} contour. The highest point is \SI{680}{m}. Describe the strata and their attitude.
\tcblower
\textbf{Solution (short factual sentences, as required in the examination).}
\begin{enumerate}
\item The area is hilly, rising from below \SI{500}{m} to \SI{680}{m}.
\item Three sedimentary units outcrop: conglomerate, sandstone and shale.
\item Both contacts are parallel to contour lines (\SI{500}{m} and \SI{600}{m} contours respectively).
\item Therefore the beds are \textbf{horizontal} (dip $0^\circ$).
\item The conglomerate outcrops on the lower ground below \SI{500}{m}; the sandstone forms a band on the slopes between \SI{500}{m} and \SI{600}{m}; the shale caps the hills above \SI{600}{m}.
\item The sandstone is \SI{100}{m} thick; the shale is at least \SI{80}{m} thick (its top is eroded).
\item From superposition, the order from oldest to youngest is: conglomerate, sandstone, shale.
\end{enumerate}
\end{exoresolu}

\begin{exercice}[Outcrop patterns and dip]
On a map, a bed of limestone between two shales crosses a river valley. Its contacts make sharp V's that point \emph{upstream}, and the contours are widely spaced on the southern side of the limestone ridge but closely spaced on its northern side. A dip symbol near the ridge reads $15^\circ$.
\begin{enumerate}
\item State, with reasons, the attitude of the beds.
\item Give the dip direction, justifying from the slope asymmetry.
\item Explain why the limestone forms a ridge while the shales form lower ground.
\end{enumerate}
\end{exercice}

\begin{corrige}[Exercise 1]
\begin{enumerate}
\item The beds are \textbf{dipping} (not horizontal, because the contacts do not follow the contours; not vertical, because they deflect into V's across the valley). The symbol confirms a dip of $15^\circ$.
\item The gentle slope (widely spaced contours) on the southern side is the \textbf{dip slope}; the steep northern side is the \textbf{scarp slope}. The beds therefore dip $15^\circ$ towards the \textbf{south}. (The upstream V is consistent with a gentle downstream dip; the river flows northwards, so downstream is south, matching the dip direction. The slope asymmetry is the decisive evidence.)
\item The limestone is more \textbf{resistant to erosion} than the shales, so it stands out as a ridge (a cuesta), while the softer shales are worn down into lower ground — differential erosion.
\end{enumerate}
\end{corrige}

\begin{aretenir}[Key points of Section 9]
\begin{itemize}
\item Describe any map in the fixed order: \cle{relief} $\rightarrow$ \cle{rock units} $\rightarrow$ \cle{attitude (dip)} $\rightarrow$ \cle{structures} $\rightarrow$ \cle{history}.
\item \textbf{Horizontal beds:} contacts parallel to contours; outcrops mimic the relief; oldest beds in the valleys, youngest on the hilltops.
\item \textbf{Vertical beds:} straight contacts independent of relief; outcrop width $=$ true thickness; younging direction from legend or sedimentary structures only.
\item \textbf{Dipping beds:} V-rule in valleys, dip symbols (strike bar, dip tick, angle), and scarp/dip-slope asymmetry (gentle slope $=$ dip direction).
\item Dip without symbols: strike lines from contact--contour intersections; $\tan\theta = \dfrac{\text{vertical interval}}{\text{horizontal spacing}}$ (measured perpendicular to strike).
\item Read the \textbf{stratigraphic order} from the legend (youngest at top, unless indicated otherwise) and from superposition (youngest in dip direction for dipping beds, unless overturned).
\item Write \textbf{short factual sentences}, with locations and values: description is marked point by point.
\end{itemize}
\end{aretenir}

With the description of strata, dips and outcrop patterns mastered, you are ready for Section 10, where the same five-point plan will be applied to the more complex structures — folds, faults, unconformities and intrusions — and to the complete reconstruction of a region's geological history.

\coursec{Section 10: Description of Geological Structures on Maps II — Folds, Faults, Unconformities, Intrusions and Geological History}

In Section 9 we learned to read the simplest message of a geological map: how horizontal and dipping strata crop out, how the \cle{V-rule} reveals dip direction in valleys, and how outcrop width depends on dip and slope. Real maps, however, rarely show simple parallel bands. Beds bend, break, are cut by younger rocks, or are buried beneath younger cover. This section completes your map-reading toolkit: you will learn to recognise and describe \cle{folds}, \cle{faults}, \cle{unconformities} and \cle{intrusions} purely from map evidence, and then to assemble all the observations into a logical \cle{geological history} — the skill that earns the highest marks in the GCE Advanced Level practical paper.

\courssub{Recognising Folds on Maps}

\textbf{The map signature of a fold.}\par
Imagine a layered cake that has been bent into an arch or a trough and then sliced horizontally: the slice shows bands repeated symmetrically on either side of a central zone. A fold does exactly this to strata. On a geological map, a fold is recognised by three converging lines of evidence:

\begin{enumerate}
\item \textbf{Repetition of beds in mirror image:} the same sequence of units appears on both sides of a central zone, in reverse order (e.g.\ A--B--C--B--A across the fold axis).
\item \textbf{Mirror-image dips:} dip symbols on the two limbs point \emph{away from} the fold axis (anticline) or \emph{towards} it (syncline).
\item \textbf{Closure of outcrops (plunge):} if the fold \emph{plunges}, the outcrop pattern of each bed closes in a nose or U-shape pointing in the direction of plunge. A non-plunging (horizontal) fold gives parallel, straight bands with no closure.
\end{enumerate}

\begin{propriete}[Ages in the core of a fold]
In an \cle{anticline}, the \textbf{oldest} rocks occupy the core (axial zone) and successively younger beds occur outwards on both limbs. In a \cle{syncline}, the \textbf{youngest} rocks occupy the core and older beds occur outwards. This criterion is decisive: dips alone can be ambiguous where dip symbols are scarce, but ages never lie. (Not examinable as a proof; must be applied correctly.)
\end{propriete}

\begin{attention}[Repetition is not enough]
Repetition of beds across a line can be produced by a fold \emph{or} by a fault. Before concluding ``fold'', check the dip symbols: mirror-image dips confirm a fold; parallel dips with a repeated or omitted sequence indicate a fault.
\end{attention}

\begin{definition}[Fold morphology terms]
\begin{itemize}
\item \textbf{Axial plane trace:} the line on the map where the axial plane intersects the ground surface.
\item \textbf{Fold axis:} the line connecting points of maximum curvature on a folded surface; on a map it coincides with the axial plane trace for upright folds.
\item \textbf{Plunge:} the angle between the fold axis and the horizontal. A plunging fold shows closed outcrop patterns (noses); a non-plunging fold shows straight, parallel outcrop bands.
\item \textbf{Anticline / Syncline:} defined by relative ages in the core (oldest = anticline, youngest = syncline), not by shape (an anticline can form a valley, a syncline a hill — topography depends on erosion resistance).
\item \textbf{Symmetrical fold:} limbs dip at equal angles in opposite directions; axial plane is vertical.
\item \textbf{Asymmetrical fold:} limbs dip at different angles; axial plane is inclined.
\item \textbf{Overturned fold:} both limbs dip in the same direction, but one limb is overturned (dip > 90°). The younging direction (way-up criteria) is essential to identify this.
\item \textbf{Recumbent fold:} axial plane is nearly horizontal; one limb is completely inverted.
\end{itemize}
\end{definition}

\begin{popfigure}
\begin{tikzpicture}[scale=0.85]
% ---- Plunging anticline (left) ----
\begin{scope}
% Core (oldest) - gold
\fill[popgold!70] (0,0) .. controls (1.0,1.0) and (1.0,3.0) .. (0,4) -- (1.4,4) .. controls (2.0,3.0) and (2.0,1.0) .. (1.4,0) -- cycle;
% Middle - green
\fill[popgreen!60] (1.4,0) .. controls (2.0,1.0) and (2.0,3.0) .. (1.4,4) -- (2.8,4) .. controls (3.2,3.0) and (3.2,1.0) .. (2.8,0) -- cycle;
% Outer - blue
\fill[popblue!50] (2.8,0) .. controls (3.2,1.0) and (3.2,3.0) .. (2.8,4) -- (4.2,4) -- (4.2,0) -- cycle;
% Mirror left side
\fill[popblue!50] (0,0) -- (1.4,0) .. controls (2.0,1.0) and (2.0,3.0) .. (1.4,4) -- (0,4) -- cycle;
\fill[popgreen!60] (0,0) -- (1.4,0) .. controls (2.0,1.0) and (2.0,3.0) .. (1.4,4) -- (0,4) -- cycle;
\fill[popgold!70] (0,0) .. controls (1.0,1.0) and (1.0,3.0) .. (0,4) -- (1.4,4) .. controls (2.0,3.0) and (2.0,1.0) .. (1.4,0) -- cycle;
% Boundaries
\draw[thick,popdark] (0,0) .. controls (1.0,1.0) and (1.0,3.0) .. (0,4);
\draw[thick,popdark] (1.4,0) .. controls (2.0,1.0) and (2.0,3.0) .. (1.4,4);
\draw[thick,popdark] (2.8,0) .. controls (3.2,1.0) and (3.2,3.0) .. (2.8,4);
\draw[thick,popdark] (0,0) rectangle (4.2,4);
% Axial trace
\draw[very thick,poporange,dashed] (0.7,0) -- (0.7,4);
\node[popdark] at (0.7,2) {\small C};
\node[popdark] at (2.1,2) {\small B};
\node[popdark] at (3.5,2) {\small A};
\node[popdark] at (-0.3,2) {\small B};
% Dip symbols (away from axis)
\draw[popink,thick] (3.0,3.3) -- (3.4,3.3); \draw[popink,thick] (3.4,3.3) -- (3.4,3.0);
\draw[popink,thick] (1.0,0.7) -- (0.6,0.7); \draw[popink,thick] (0.6,0.7) -- (0.6,1.0);
\node[align=center] at (2.1,-0.6) {\small Plunging \cle{anticline}\\ \small oldest (C) in core\\ \small plunges North};
\end{scope}
% ---- Plunging syncline (right) ----
\begin{scope}[xshift=6.0cm]
% Core (youngest) - blue
\fill[popblue!50] (0,0) -- (1.2,0) .. controls (0.8,1.0) and (0.8,3.0) .. (1.2,4) -- (0,4) -- cycle;
% Middle - green
\fill[popgreen!60] (1.2,0) .. controls (0.8,1.0) and (0.8,3.0) .. (1.2,4) -- (2.6,4) .. controls (2.2,3.0) and (2.2,1.0) .. (2.6,0) -- cycle;
% Outer - gold
\fill[popgold!70] (2.6,0) .. controls (2.2,1.0) and (2.2,3.0) .. (2.6,4) -- (4.2,4) -- (4.2,0) -- cycle;
% Boundaries
\draw[thick,popdark] (1.2,0) .. controls (0.8,1.0) and (0.8,3.0) .. (1.2,4);
\draw[thick,popdark] (2.6,0) .. controls (2.2,1.0) and (2.2,3.0) .. (2.6,4);
\draw[thick,popdark] (0,0) rectangle (4.2,4);
% Axial trace
\draw[very thick,poporange,dashed] (3.4,0) -- (3.4,4);
\node[popdark] at (3.4,2) {\small A};
\node[popdark] at (1.9,2) {\small B};
\node[popdark] at (0.6,2) {\small C};
% Dip symbols (towards axis)
\draw[popink,thick] (1.8,3.3) -- (2.2,3.3); \draw[popink,thick] (1.8,3.3) -- (1.8,3.0);
\draw[popink,thick] (0.8,0.7) -- (0.4,0.7); \draw[popink,thick] (0.8,0.7) -- (0.8,1.0);
\node[align=center] at (2.1,-0.6) {\small Plunging \cle{syncline}\\ \small youngest (A) in core\\ \small plunges North};
\end{scope}
\end{tikzpicture}
\legende{Map patterns of plunging folds. Ages: A youngest, C oldest. Dips point away from the anticline axis and towards the syncline axis; outcrops close (nose) in the direction of plunge (North here).}
\end{popfigure}

\begin{exemplebox}[Identifying fold type and plunge from a map]
\textbf{Map evidence:} Units X (shale), Y (sandstone), Z (limestone) appear in order X--Y--Z--Y--X across a N--S zone. Dips are 30° away from the zone on both sides. Outcrops close to the North.
\textbf{Reasoning:} (1) Mirror repetition $\rightarrow$ fold. (2) Dips away from centre $\rightarrow$ anticline. (3) Oldest (X) in core $\rightarrow$ confirms anticline. (4) Closure to North $\rightarrow$ plunges North.
\textbf{Answer:} North-plunging anticline with X oldest in core.
\end{exemplebox}

\courssub{Recognising and Classifying Faults on Maps}

\textbf{The map signature of a fault.}\par
A fault is a fracture along which displacement has occurred. On a map it appears as a line (often drawn thicker, or ornamented with ticks on the downthrown side or balls on the upthrown side) along which geological contacts are abruptly \emph{cut}. The diagnostic evidence is:

\begin{itemize}
\item \textbf{Truncation and offset of contacts:} a bed boundary stops suddenly at the fault and reappears, displaced, on the other side.
\item \textbf{Omission of beds:} some units present on one side are missing on the other (typical of normal faults cutting dipping strata).
\item \textbf{Repetition of beds:} the same unit appears twice across the fault (typical of reverse/thrust faults).
\item \textbf{Fault-rock symbols:} some maps show fault breccia, gouge, or crushed rock along the fault line.
\item \textbf{Topographic expression:} fault scarps, linear valleys, or offset drainage may hint at faulting (though not definitive on a geological map alone).
\end{itemize}

\begin{definition}[Fault terminology]
\begin{itemize}
\item \textbf{Hanging wall:} the block above the fault plane.
\item \textbf{Footwall:} the block below the fault plane.
\item \textbf{Throw:} the vertical component of displacement (difference in elevation of a marker horizon across the fault).
\item \textbf{Heave:} the horizontal component of displacement measured perpendicular to the strike of the fault.
\item \textbf{Net slip:} the total displacement vector along the fault plane.
\item \textbf{Cut-off line:} the line of intersection between a bed and the fault plane.
\end{itemize}
\end{definition}

\textbf{Classifying the fault from map evidence.}\par
\begin{itemize}
\item \textbf{Dip-slip faults (normal or reverse):} contacts are offset \emph{across the strike}. Compare the two sides: the side where \emph{older} rocks are juxtaposed against younger is the \emph{upthrown} side. 
  \begin{itemize}
  \item \textbf{Normal fault:} hanging wall moves down relative to footwall $\Rightarrow$ \emph{omission} of beds in the downthrown block. The fault plane dips towards the downthrown side.
  \item \textbf{Reverse fault:} hanging wall moves up relative to footwall $\Rightarrow$ \emph{repetition} of beds in the upthrown block. The fault plane dips towards the upthrown side. A \textbf{thrust fault} is a low-angle reverse fault ($< 30^\circ$).
  \end{itemize}
\item \textbf{Strike-slip (wrench) faults:} contacts are shifted \emph{horizontally} along the fault; a vertical dyke or a fold axis is displaced sideways without omission or repetition. If the opposite block appears to have moved to the right, the fault is \emph{dextral}; to the left, \emph{sinistral}.
\item \textbf{Oblique-slip faults:} combine dip-slip and strike-slip components; show both omission/repetition and lateral offset.
\end{itemize}

\begin{methode}[Determining fault type and throw from a map]
\begin{enumerate}
\item Identify a distinctive marker bed (key bed) cut by the fault on both sides.
\item Note the age sequence on each side of the fault.
\item \textbf{Omission} of beds on one side $\rightarrow$ normal fault; that side is downthrown.
\item \textbf{Repetition} of beds on one side $\rightarrow$ reverse fault; that side is upthrown.
\item \textbf{Lateral offset} of a linear feature (fold axis, dyke) without omission/repetition $\rightarrow$ strike-slip fault; determine sense (dextral/sinistral).
\item \textbf{Throw calculation:} if structure contours or spot heights are given, throw = difference in elevation of the same marker horizon on the two sides. On a cross-section, throw = vertical separation of the cut-off lines.
\end{enumerate}
\end{methode}

\begin{exemplebox}[Calculating the throw]
The \cle{throw} is the vertical component of displacement. If a key bed (e.g., a limestone marker) crops out at \SI{400}{m} above sea level on the upthrown side of a fault and the same bed (projected) lies at \SI{250}{m} on the downthrown side, the throw is $400 - 250 = \SI{150}{m}$. On cross-sections, the throw is read directly as the vertical separation of a cut-off bed across the fault.
\end{exemplebox}

\begin{popfigure}
\begin{tikzpicture}[scale=0.9]
% Left block: folded sequence (anticline)
\fill[popgold!70] (0,0) rectangle (1.0,4);
\fill[popgreen!60] (1.0,0) rectangle (2.0,4);
\fill[popblue!50] (2.0,0) rectangle (3.0,4);
\fill[popgreen!60] (3.0,0) rectangle (4.0,4);
\fill[popgold!70] (4.0,0) rectangle (5.0,4);
% Right block: same sequence shifted right-laterally
\begin{scope}[xshift=5.6cm]
\fill[popgold!70] (0,0) rectangle (1.0,4);
\fill[popgreen!60] (1.0,0) rectangle (2.0,4);
\fill[popblue!50] (2.0,0) rectangle (3.0,4);
\fill[popgreen!60] (3.0,0) rectangle (4.0,4);
\fill[popgold!70] (4.0,0) rectangle (5.0,4);
\end{scope}
% Fault line
\draw[very thick,poporange] (5,0) -- (5.6,4);
\node[poporange,rotate=78] at (5.05,3.6) {\small fault};
% Fold axes truncated
\draw[very thick,popdark,dashed] (2.5,0) -- (2.5,4);
\draw[very thick,popdark,dashed] (8.1,0) -- (8.1,4);
\node[popdark] at (2.5,4.25) {\small axis};
\node[popdark] at (8.1,4.25) {\small axis};
% Offset arrow
\draw[<->,very thick,popink] (2.5,-0.35) -- (8.1,-0.35);
\node[popink] at (5.3,-0.7) {\small horizontal offset of fold axis};
% Unit labels
\node[popdark] at (0.5,2) {\small A};
\node[popdark] at (1.5,2) {\small B};
\node[popdark] at (2.5,2) {\small C};
\node[popdark] at (3.5,2) {\small B};
\node[popdark] at (4.5,2) {\small A};
\node[popdark] at (6.1,2) {\small A};
\node[popdark] at (7.1,2) {\small B};
\node[popdark] at (8.1,2) {\small C};
\node[popdark] at (9.1,2) {\small B};
\node[popdark] at (10.1,2) {\small A};
\end{tikzpicture}
\legende{A dextral strike-slip fault truncating and offsetting the axis of an anticline (unit C oldest, in the core). Note: no bed is omitted or repeated — the displacement is purely horizontal.}
\end{popfigure}

\begin{exemplebox}[Normal vs. reverse fault on a map]
\textbf{Map A:} Sequence Sandstone--Shale--Limestone dips 20° East. A N--S fault cuts them. West of fault: all three units present. East of fault: only Sandstone and Shale; Limestone is missing.
\textbf{Interpretation:} Omission of Limestone on east side $\rightarrow$ east side downthrown $\rightarrow$ \textbf{normal fault}, downthrow to East.
\textbf{Map B:} Same sequence, same dip. West of fault: Sandstone--Shale. East of fault: Sandstone--Shale--Limestone--Shale--Sandstone (repetition).
\textbf{Interpretation:} Repetition on east side $\rightarrow$ east side upthrown $\rightarrow$ \textbf{reverse fault}, upthrow to East.
\end{exemplebox}

\courssub{Recognising Unconformities on Maps}

An \cle{unconformity} records a gap in geological time: older rocks were deformed, uplifted and eroded before younger sediments were deposited on them. Map evidence includes:

\begin{itemize}
\item \textbf{Angular discordance:} the contacts of the older series trend in a different direction from those of the younger series; the younger beds \emph{truncate} (cut across) the older contacts.
\item \textbf{Overstep:} a single younger bed rests successively on several different older beds as it is traced along the unconformity line (the younger series overlaps the eroded edges of the older series).
\item \textbf{Basal conglomerate:} the lowest bed of the younger series is a conglomerate containing pebbles of the underlying rocks — often labelled on the map legend.
\item \textbf{Truncation of older structures:} fold axes, faults, and cleavage in the older series stop abruptly at the unconformity and do not continue into the younger cover.
\item \textbf{Relief on the unconformity surface:} the unconformity may cut across valleys and hills of the old landscape; the younger beds may show onlap onto palaeohighs.
\end{itemize}

\begin{definition}[Types of unconformity]
\begin{itemize}
\item \textbf{Angular unconformity:} younger beds rest on tilted/folded older beds with a clear angular discordance.
\item \textbf{Disconformity:} parallel bedding above and below, but a time gap shown by erosion surface, palaeosol, or fossil gap. Hard to see on a map unless a basal conglomerate is marked.
\item \textbf{Nonconformity:} sedimentary rocks rest on eroded igneous or metamorphic basement.
\item \textbf{Paraconformity:} bedding parallel, no visible erosion surface; gap known only from fossils. Not mappable from geometry alone.
\end{itemize}
\end{definition}

\begin{popfigure}
\begin{tikzpicture}[scale=0.9]
% Older folded series (tilted)
\fill[popgold!70] (0,0) -- (1.5,0) -- (2.5,4) -- (1.0,4) -- cycle;
\fill[popgreen!60] (1.5,0) -- (3.0,0) -- (4.0,4) -- (2.5,4) -- cycle;
\fill[popblue!50] (3.0,0) -- (4.5,0) -- (5.5,4) -- (4.0,4) -- cycle;
\fill[popgreen!60] (4.5,0) -- (6.0,0) -- (7.0,4) -- (5.5,4) -- cycle;
\fill[popgold!70] (6.0,0) -- (7.5,0) -- (8.5,4) -- (7.0,4) -- cycle;
% Unconformity line (horizontal)
\draw[very thick,poporange] (0,2.8) -- (8.5,2.8);
\node[poporange,above] at (4.25,2.8) {\small Unconformity};
% Younger horizontal cover
\fill[poppink!50] (0,2.8) rectangle (8.5,4);
% Labels
\node[popdark,rotate=45] at (1.2,1.5) {\small Oldest};
\node[popdark,rotate=45] at (3.7,1.5) {\small Youngest};
\node[popdark] at (4.25,3.4) {\small Young cover (horizontal)};
% Truncation arrows
\draw[popink,->,thick] (2.5,2.8) -- (2.5,2.2);
\draw[popink,->,thick] (4.0,2.8) -- (4.0,2.2);
\draw[popink,->,thick] (5.5,2.8) -- (5.5,2.2);
\node[popink,align=center] at (4.25,1.8) {\small Truncation of\\ \small older contacts};
\end{tikzpicture}
\legende{Angular unconformity: tilted, folded older series (dipping East) truncated by horizontal younger cover. Note the overstep — the cover rests on progressively younger older-beds from left to right.}
\end{popfigure}

\courssub{Recognising Intrusions on Maps}

Igneous intrusions are identified by their geometric relations with the country rocks:

\begin{itemize}
\item \textbf{Cross-cutting contacts:} the intrusion cuts across bedding contacts — proof that it is younger than the rocks it cuts.
\item \textbf{Baked margins (contact aureole):} a zone of metamorphosed rock (hornfels, spotted slate) may fringe the intrusion, shown by a special ornament or stipple on the map.
\item \textbf{Chilled margins:} fine-grained edges of the intrusion where it cooled rapidly against cold country rock.
\item \textbf{Xenoliths:} fragments of country rock enclosed in the intrusion — older than the intrusion.
\end{itemize}

\begin{definition}[Common intrusion types on maps]
\begin{itemize}
\item \textbf{Dyke:} narrow, sheet-like, discordant body cutting across bedding. Straight or slightly curved outcrop. Vertical or steeply dipping. Younger than all rocks it cuts.
\item \textbf{Sill:} sheet-like, \emph{concordant} body intruded parallel to bedding. Mimics a sedimentary layer. Identified by: (a) rock name in legend (igneous), (b) baked margins \emph{above and below}, (c) may transgress bedding gently (low-angle discordance) over long distances.
\item \textbf{Batholith:} large ($> 100\ \mathrm{km}^2$), irregular, discordant pluton. Oval to elliptical outcrop. Contacts cut country rock. Often shows contact aureole.
\item \textbf{Stock:} smaller version of a batholith ($< 100\ \mathrm{km}^2$).
\item \textbf{Laccolith:} concordant, lens-shaped intrusion that domes the overlying strata. Outcrop shows central igneous core with sedimentary strata dipping away radially.
\item \textbf{Lopolith:} concordant, saucer-shaped intrusion; strata dip inward towards the centre.
\end{itemize}
\end{definition}

\begin{attention}[A dyke is younger than everything it cuts]
Whatever the age of the surrounding rocks, a dyke (or any intrusion) is younger than the youngest bed it crosses — but older than any unconformity or fault that cuts \emph{it}. Always check both relations before placing an intrusion in the history.
\end{attention}

\begin{popfigure}
\begin{tikzpicture}[scale=0.9]
% Country rock: horizontal beds
\fill[popblue!30] (0,0) rectangle (8,1.5); \node[popdark] at (4,0.75) {\small Shale};
\fill[popgreen!40] (0,1.5) rectangle (8,3.0); \node[popdark] at (4,2.25) {\small Sandstone};
\fill[popgold!50] (0,3.0) rectangle (8,4.5); \node[popdark] at (4,3.75) {\small Limestone};
% Vertical dyke
\fill[poppurple!60] (3.5,0) rectangle (4.5,4.5);
\draw[thick,popdark] (3.5,0) -- (3.5,4.5);
\draw[thick,popdark] (4.5,0) -- (4.5,4.5);
\node[popdark,rotate=90] at (4.0,2.25) {\small Dolerite dyke};
% Baked margins
\draw[poppink,thick,dashed] (3.3,0) -- (3.3,4.5);
\draw[poppink,thick,dashed] (4.7,0) -- (4.7,4.5);
\node[poppink,rotate=90] at (3.3,2.25) {\small Baked margin};
\node[poppink,rotate=90] at (4.7,2.25) {\small Baked margin};
% Sill (concordant)
\fill[poppurple!60] (5.5,1.5) rectangle (7.5,2.0);
\draw[thick,popdark] (5.5,1.5) -- (7.5,1.5);
\draw[thick,popdark] (5.5,2.0) -- (7.5,2.0);
\node[popdark] at (6.5,1.75) {\small Sill};
% Baked margins on sill
\draw[poppink,thick,dashed] (5.5,1.3) -- (7.5,1.3);
\draw[poppink,thick,dashed] (5.5,2.2) -- (7.5,2.2);
\end{tikzpicture}
\legende{Intrusions on a map: a vertical dyke cuts across all beds (discordant) with baked margins on both sides; a sill lies parallel to bedding (concordant) with baked margins above and below.}
\end{popfigure}

\courssub{Establishing the Geological History}

The geological history is the ordered story of events that produced the map. It is built using the criteria of \cle{relative dating}. The history must be written as a \textbf{numbered list from the oldest event to the youngest event}, each justified by a criterion.

\begin{methode}[Step-by-step construction of a geological history]
\begin{enumerate}
\item \textbf{List all rock units} from the map legend with their ages (oldest to youngest) using superposition where strata are undisturbed.
\item \textbf{Identify all structures:} folds (type, plunge), faults (type, throw, sense), unconformities (type, what is truncated), intrusions (type, what they cut).
\item \textbf{Apply relative dating criteria} to order events:
  \begin{itemize}
  \item \textbf{Superposition:} in an undisturbed sedimentary sequence, the lowest bed is the oldest.
  \item \textbf{Cross-cutting relationships:} a fault, dyke, or unconformity is younger than every rock it cuts.
  \item \textbf{Unconformity:} the rocks below are older than those above; folding, faulting and erosion of the lower series all predate the upper series.
  \item \textbf{Inclusions:} fragments (pebbles in conglomerate, xenoliths in intrusion) are older than the rock that contains them.
  \item \textbf{Baked/chilled margins:} confirm intrusion is younger than country rock.
  \end{itemize}
\item \textbf{Write the history} as a numbered sequence: 1. Deposition of oldest unit... 2. Deposition of next unit... n. Folding event... n+1. Faulting... n+2. Erosion/unconformity... n+3. Deposition of cover... etc.
\item \textbf{Check consistency:} no event should contradict another (e.g., a fault cannot cut a dyke that cuts the fault).
\end{enumerate}
\end{methode}

\begin{popfigure}
\begin{tikzpicture}[scale=1.0, every node/.style={align=center}]
\node[draw,rounded corners,fill=popblueL,minimum width=4.5cm, align=center] (s) at (0,0){\small 1. Deposition of\\ \small sedimentary sequence\\ \small (superposition)};
\node[draw,rounded corners,fill=popgreenL,minimum width=4.5cm, align=center] (f) at (7,1.8){\small 2. Folding of\\ \small the sequence};
\node[draw,rounded corners,fill=popgreenL,minimum width=4.5cm, align=center] (ft) at (7,-1.8){\small 3. Faulting /\\ \small intrusion (cuts beds)};
\node[draw,rounded corners,fill=popgold!50,minimum width=4.5cm, align=center] (u) at (14,0){\small 4. Uplift, erosion,\\ \small unconformity surface};
\node[draw,rounded corners,fill=popblueL,minimum width=4.5cm, align=center] (c) at (21,0){\small 5. Deposition of\\ \small younger cover};
\draw[->,thick,popdark] (s) -- (f);
\draw[->,thick,popdark] (s) -- (ft);
\draw[->,thick,popdark] (f) -- (u);
\draw[->,thick,popdark] (ft) -- (u);
\draw[->,thick,popdark] (u) -- (c);
\end{tikzpicture}
\legende{Event-ordering logic: deposition first, then deformation (folding, faulting, intrusion), then erosion and younger cover. Anything truncated by the unconformity is older than the cover.}
\end{popfigure}

\begin{exoresolu}[Complete description and history of a compound map]
\textbf{Statement.} A map shows, from west to east, units A (conglomerate), B (sandstone) and C (limestone) repeated symmetrically about a N--S zone of unit A, with dips of $30^{\circ}$ away from this zone; the outcrops close northwards. This folded series is truncated along an E--W line by horizontal beds D (basal conglomerate) and E (shale). A vertical dolerite dyke cuts units A, B, C and D but stops at E. Describe the structures and give the geological history.
\tcblower
\textbf{Solution.} 
\emph{Structures:} 
\begin{enumerate}
\item The repetition A--B--C--B--A with outward dips and oldest unit A in the core indicates a \textbf{plunging anticline}, plunging north (closure of outcrops to the north).
\item The E--W line truncating the folded series, with horizontal beds above, is an \textbf{angular unconformity}; bed D is its basal conglomerate (contains pebbles of A, B, C).
\item The dolerite body is a \textbf{discordant intrusion (dyke)} — vertical, cross-cutting.
\end{enumerate}
\emph{Geological history (oldest to youngest):}
\begin{enumerate}
\item Deposition of units A, B and C (superposition: A oldest, C youngest).
\item Folding of A--C into a north-plunging anticline.
\item Uplift and erosion, producing the unconformity surface (truncates fold).
\item Deposition of the basal conglomerate D (on unconformity).
\item Intrusion of the dolerite dyke (cuts A, B, C, D $\rightarrow$ younger than D).
\item Deposition of shale E (dyke does not cut E $\rightarrow$ E younger than dyke).
\end{enumerate}
\end{exoresolu}

\begin{center}
\begin{tabularx}{\linewidth}{|c|X|l|}\hline
\rowcolor{popblueL} \textbf{Order} & \textbf{Event} & \textbf{Criterion used} \\ \hline
1 & Deposition of A, B, C & Superposition \\ \hline
2 & Folding (N-plunging anticline) & Truncated by unconformity \\ \hline
3 & Erosion (unconformity surface) & Angular discordance, overstep, basal conglomerate \\ \hline
4 & Deposition of D (basal conglomerate) & Rests on unconformity, contains pebbles of A--C \\ \hline
5 & Dolerite dyke intrusion & Cross-cuts A--D; baked margins \\ \hline
6 & Deposition of E (shale) & Not cut by dyke; horizontal, unconformable on D? (here conformable on D) \\ \hline
\end{tabularx}
\end{center}

\begin{attention}[Exam technique]
Present the geological history as a \textbf{numbered list from the oldest to the youngest event}, one event per line, each justified by a criterion. Never mix description and history: describe the structures first, then give the history. Vague statements such as ``the rocks were folded and faulted'' without order or evidence earn little credit. Always use the terminology: ``Deposition of...'', ``Folding of...'', ``Faulting (normal, downthrow to...)'', ``Intrusion of...'', ``Erosion and unconformity'', ``Deposition of...''.
\end{attention}

\begin{exercice}[Fault or fold?]
On a map, the sequence sandstone--limestone--sandstone is observed across a straight line, with all dips $25^{\circ}$ to the east. (a) Is this structure a fold or a fault? Justify. (b) What type of fault is it, and which side is downthrown? (c) Rewrite the answer assuming dips are $25^{\circ}$E on the west side and $25^{\circ}$W on the east side.
\end{exercice}

\begin{corrige}[Exercise 1]
(a) A \textbf{fault}: the dips are parallel (all eastward), so the repetition of sandstone is not due to folding (which requires mirror-image dips). (b) Repetition of a bed (sandstone appears twice) indicates a \textbf{reverse fault}; the western side, where older rocks (limestone is younger than sandstone? Wait — need age order) — assuming sandstone is older than limestone, the western side brings older sandstone against younger limestone, so western side is upthrown, eastern side is downthrown. (c) With mirror-image dips the structure is a \textbf{syncline} (dips converge towards the axis), with limestone — the youngest bed — in the core.
\end{corrige}

\begin{exercice}[Unconformity or intrusion?]
A map shows horizontal beds of shale (Unit X) overlying tilted beds of sandstone (Unit Y). The contact between X and Y is a straight line trending N--S. Unit X contains pebbles of Unit Y at its base. A vertical dyke cuts Unit Y but stops at the base of Unit X. (a) What is the contact between X and Y? (b) What is the relative age of the dyke? (c) Write the geological history.
\end{exercice}

\begin{corrige}[Exercise 2]
(a) \textbf{Angular unconformity}: angular discordance (horizontal vs tilted), basal conglomerate (pebbles of Y in X), truncation of Y contacts. (b) Dyke cuts Y but not X $\rightarrow$ dyke is younger than Y but older than X. (c) History: 1. Deposition of Y. 2. Tilting/folding of Y. 3. Erosion (unconformity). 4. Dyke intrusion (cuts Y). 5. Deposition of X (basal conglomerate with Y pebbles, then shale).
\end{corrige}

\begin{savaistu}[Map work in Cameroon]
Geological maps of the Mount Cameroon area show exactly these principles: recent lava flows (the youngest units) cross-cut and cover older pyroclastic deposits, while faults of the Cameroon Volcanic Line truncate older contacts. The Cretaceous sediments of the Rio del Rey Basin rest with angular unconformity on the Precambrian basement of the Congo Craton — a natural field laboratory for everything in this section.
\end{savaistu}

\begin{aretenir}[Key points]
\begin{itemize}
\item Fold: mirror-image repetition + opposed dips; anticline = oldest in core, syncline = youngest in core; closure shows plunge direction.
\item Fault: truncation/offset of contacts; omission = normal, repetition = reverse, sideways shift = strike-slip. Throw = vertical separation of a marker horizon.
\item Unconformity: angular discordance, overstep, basal conglomerate, truncated older structures. Types: angular, disconformity, nonconformity.
\item Intrusion: cross-cutting contacts and baked margins; younger than all it cuts. Dyke (discordant), sill (concordant), batholith (large discordant).
\item History: order events by superposition, cross-cutting, unconformity and inclusions; write a numbered list, oldest first, each justified.
\end{itemize}
\end{aretenir}

\coursec{Section 11: Interpretation of Geological Photographs}

In Sections 9 and 10 we learnt to read geological structures from maps — the \emph{plan view} of the Earth's surface. But the geologist in the field rarely sees the ground from directly above: he sees cliffs, hill slopes, road cuttings and distant mountains in \emph{perspective}. The GCE Advanced Level examination therefore frequently presents candidates with a \cle{geological photograph} — of Mount Cameroon, of a Mandara inselberg, of a road cutting near Bafoussam — and asks them to describe and interpret it. Photograph interpretation is simply map work translated into three dimensions: the same rocks, the same structures, the same logic, but read from visual clues rather than from printed symbols. This section gives you a rigorous, systematic method.

\courssub{Types of Geological Photographs}

\begin{definition}[Types of geological photographs]
A \cle{geological photograph} is any photographic image used to study rocks, landforms and structures. Four types are distinguished by the position and orientation of the camera or sensor:
\begin{itemize}
\item \cle{Ground photograph}: taken from the Earth's surface, horizontally or obliquely (e.g.\ a photograph of a cliff face or of Mount Cameroon from Buea). It shows structures in profile or perspective but covers a small area.
\item \cle{Oblique aerial photograph}: taken from an aircraft with the camera tilted, so that the horizon appears. It combines perspective (like a ground photograph) with a wide coverage; relief appears exaggerated towards the foreground and compressed towards the horizon.
\item \cle{Vertical aerial photograph}: taken with the camera pointing straight down. It is a near-map view used for photogeology; relief is revealed by shadows and, when two overlapping photographs are viewed together as a \cle{stereoscopic pair}, in apparent three dimensions.
\item \cle{Satellite image}: recorded by sensors aboard satellites; it covers very large areas (whole regions of Cameroon in one scene) and may include infrared bands that reveal vegetation and rock differences invisible to the eye.
\end{itemize}
\end{definition}

\begin{savaistu}[Photogeology in Cameroon]
Vertical aerial photographs and satellite images have been essential in mapping remote parts of Cameroon — the Adamaoua Plateau, the northern Benue plains and the dense forests of the South — where ground traverses on foot are slow and costly. Lineaments visible on satellite images have helped geologists trace the great fracture zones associated with the \cle{Cameroon Volcanic Line}.
\end{savaistu}

\courssub{A Systematic Approach to Photograph Interpretation}

The greatest danger in a photograph question is to jump straight to conclusions (``this is a fault!'') without first describing the evidence. Interpretation must proceed in a fixed order, from observation to deduction.

\begin{methode}[The 4-step photograph interpretation procedure]
\begin{enumerate}
\item \textbf{Locate and orient.} Identify the type of photograph (ground, oblique, vertical, satellite), the direction of view if given, the scale indicators (houses, trees, roads, people) and the position of the horizon. Estimate the approximate size of the area shown.
\item \textbf{Describe the landforms.} State objectively what is visible: hills, ridges, valleys, scarps, cones, flat plains; note their shapes, slopes and arrangement. Use neutral words (``a long straight ridge'', ``a steep slope facing left'').
\item \textbf{Identify rocks and structures.} Use photographic clues (tone, texture, vegetation, drainage, bedding traces) to recognise rock types and structures: dipping strata, fold closures, fault scarps, dykes, joints.
\item \textbf{Interpret the processes and history.} Only now deduce the geology: resistance of rocks to erosion, dip direction, folding, faulting, volcanism, and the sequence of events.
\end{enumerate}
\end{methode}

\begin{attention}[Exam tip: describe first, interpret after]
In GCE photograph questions, marks are awarded separately for \textbf{description} (what you \emph{see}) and for \textbf{interpretation} (what it \emph{means}). Never skip the description. Write ``a steep, straight, light-toned scarp crosses the photograph from top-left to bottom-right, with a stream offset along it'' \emph{before} writing ``this is a fault scarp''. A candidate who only writes ``fault'' loses the description marks — and if the identification is wrong, loses everything.
\end{attention}

\courssub{Reading Photographic Clues}

A photograph carries no legend and no symbols; the geologist must extract information from six families of visual clues.

\begin{center}
\begin{tabularx}{\linewidth}{|l|X|X|}
\hline
\rowcolor{popblueL} \textbf{Clue} & \textbf{What it is} & \textbf{Geological meaning} \\
\hline
Tone / colour & Lightness or darkness of the image & Light tones: bare rock, sand, limestone; dark tones: vegetation, wet ground, basalt \\
\hline
Texture & Smooth, rough, banded or mottled appearance & Banded texture = bedded strata; rough massive texture = crystalline rock; smooth = cultivated or alluvial land \\
\hline
Vegetation & Distribution and type of plant cover & Vegetation often follows rock type and soils; bare ridges may mark resistant sandstone or dykes \\
\hline
Drainage pattern & Arrangement of streams & Controlled by structure and rock type (see below) \\
\hline
Relief & Hills, ridges, scarps, depressions & Hard rocks stand out; soft rocks form valleys; scarps may mark faults or dipping strata \\
\hline
Lineaments & Long straight alignments on the image & Suspected faults, joints, dykes or lithological contacts \\
\hline
\end{tabularx}
\end{center}

\begin{definition}[Lineament]
A \cle{lineament} is any long, straight or gently curved feature visible on a photograph or image — an alignment of valleys, a tonal boundary, a straight ridge or a straight scarp.
\end{definition}

\begin{attention}[A lineament is only a SUSPECTED fault]
A straight line on a photograph is \emph{not} automatically a fault. Roads, fences, dykes, joint sets, lithological contacts and even plantation boundaries all produce lineaments. A lineament remains a \textbf{suspected} fault until confirmed by other evidence: offset of beds, a scarp with consistent downthrow, crushed rock, springs aligned along it, or displaced streams. In the examination, always write ``lineament, \emph{possibly} a fault'' and give the supporting evidence.
\end{attention}

\courssub{Drainage Patterns as Geological Indicators}

Because streams erode the weakest ground and avoid the hardest, the \emph{pattern} of a drainage network is one of the most powerful structural clues on any photograph or map.

\begin{definition}[Drainage pattern]
The \cle{drainage pattern} is the plan arrangement of a stream network. Four fundamental types are recognised:
\begin{itemize}
\item \cle{Dendritic}: branching like a tree, with tributaries joining at acute angles — develops on \textbf{homogeneous, horizontal or uniformly resistant rocks} (e.g.\ granite, flat-lying sediments).
\item \cle{Trellis}: long parallel main streams joined at right angles by short tributaries — develops on \textbf{alternating hard and soft dipping or folded strata}; the main streams occupy the soft-rock valleys.
\item \cle{Radial}: streams flowing outwards in all directions from a central high point — indicates a \textbf{volcanic cone or a dome} (e.g.\ Mount Cameroon).
\item \cle{Rectangular}: streams with right-angled bends following joint or fault directions — indicates \textbf{strongly jointed or faulted rocks}.
\end{itemize}
\end{definition}

\begin{propriete}[Drainage patterns and structure]
A \cle{trellis drainage} pattern indicates alternating hard and soft strata that are dipping or folded; the long subsequent streams follow the soft outcrops along strike. A \cle{radial drainage} pattern indicates a centre of relief such as a volcanic cone or an uplifted dome, with streams flowing down the outward slopes. (Recognition and interpretation of these patterns is directly examinable.)
\end{propriete}

\begin{popfigure}
\begin{tikzpicture}[scale=0.9, every node/.style={font=\small}]
% Dendritic
\begin{scope}[shift={(0,0)}]
\draw[popblue, thick] (0,0) -- (0,2.4);
\draw[popblue] (0,0.6) -- (-0.8,1.4); \draw[popblue] (0,0.6) -- (0.8,1.5);
\draw[popblue] (0,1.3) -- (-0.7,2.1); \draw[popblue] (0,1.3) -- (0.7,2.2);
\draw[popblue] (-0.8,1.4) -- (-1.3,2.0); \draw[popblue] (0.8,1.5) -- (1.3,2.1);
\node[align=center] at (0,-0.7) {Dendritic\\ \emph{homogeneous rock}};
\end{scope}
% Trellis
\begin{scope}[shift={(4.2,0)}]
\draw[popblue, thick] (-1.2,0) -- (-1.2,2.4);
\draw[popblue, thick] (0,0) -- (0,2.4);
\draw[popblue, thick] (1.2,0) -- (1.2,2.4);
\draw[popblue] (-1.2,0.7) -- (0,0.7); \draw[popblue] (0,1.5) -- (1.2,1.5);
\draw[popblue] (-1.2,1.8) -- (0,1.8); \draw[popblue] (0,0.4) -- (1.2,0.4);
\node[align=center] at (0,-0.7) {Trellis\\ \emph{dipping hard/soft strata}};
\end{scope}
% Radial
\begin{scope}[shift={(8.4,1.0)}]
\foreach \a in {0,45,90,135,180,225,270,315}{
\draw[popblue, thick] (0,0) -- (\a:1.15);}
\draw[popgold, fill=popgold!40] (0,0) circle (0.18);
\node[align=center] at (0,-1.9) {Radial\\ \emph{cone or dome}};
\end{scope}
% Rectangular
\begin{scope}[shift={(11.6,0)}]
\draw[popblue, thick] (-1.0,0) -- (-1.0,1.2) -- (0.2,1.2) -- (0.2,2.4);
\draw[popblue, thick] (0.2,1.2) -- (0.2,0.2) -- (1.2,0.2) -- (1.2,1.4);
\draw[popblue] (-1.0,0.5) -- (-1.6,0.5); \draw[popblue] (1.2,0.8) -- (1.7,0.8);
\node[align=center] at (0.1,-0.7) {Rectangular\\ \emph{jointed / faulted rock}};
\end{scope}
\end{tikzpicture}
\legende{The four fundamental drainage patterns and their geological significance.}
\end{popfigure}

\courssub{Identifying Structures on Photographs}

\textbf{Bedding and dip slopes.} Sedimentary strata appear as \emph{banding} — parallel stripes of differing tone or vegetation. Where a resistant bed dips gently, it forms a \cle{cuesta}: an asymmetric ridge with a gentle \cle{dip slope} following the top of the hard bed and a steep \cle{scarp slope} cutting across the bedding. Where strata dip steeply (more than about $45^\circ$), both slopes are steep and the ridge is called a \cle{hogback}.

\begin{popfigure}
\begin{tikzpicture}[scale=1.0, every node/.style={font=\small}]
% Photo frame
\draw[thick, popdark] (-0.4,-0.4) rectangle (10.6,4.4);
% Distant plain
\draw[popmuted] (0,1.0) -- (3.0,1.0);
% Cuesta profile: gentle dip slope up to crest, steep scarp down
\draw[very thick, popdark] (0.2,1.0) -- (5.2,3.4) -- (6.2,1.6) -- (10.2,1.6);
% Bedding traces on dip slope
\draw[poporange] (1.2,1.45) -- (5.35,3.15);
\draw[poporange] (2.2,1.9) -- (5.5,2.85);
\draw[poporange] (3.2,2.35) -- (5.6,2.55);
% Scarp face hachures
\draw[popgreen] (5.35,3.1) -- (5.7,2.6);
\draw[popgreen] (5.55,2.75) -- (5.9,2.25);
\draw[popgreen] (5.75,2.4) -- (6.05,1.95);
% Trees on dip slope
\draw[popgreen!60!black] (1.4,1.58) -- (1.4,1.76); \draw[popgreen!60!black, fill=popgreen!60!black] (1.4,1.85) circle (0.09);
\draw[popgreen!60!black] (2.6,2.15) -- (2.6,2.33); \draw[popgreen!60!black, fill=popgreen!60!black] (2.6,2.42) circle (0.09);
\draw[popgreen!60!black] (3.8,2.73) -- (3.8,2.91); \draw[popgreen!60!black, fill=popgreen!60!black] (3.8,3.0) circle (0.09);
% Annotations
\node[align=center, text=popblue] at (2.6,3.6) {gentle dip slope\\ (follows bedding)};
\draw[->, popblue] (2.6,3.25) -- (3.0,2.75);
\node[align=center, text=popgreen!50!black] at (8.0,3.4) {steep scarp slope\\ (cuts across beds)};
\draw[->, popgreen!50!black] (7.4,3.05) -- (5.9,2.6);
\node[align=center, text=poporange] at (1.6,0.4) {bedding traces (tone bands)};
\draw[->, poporange] (2.2,0.75) -- (2.9,1.95);
\node[align=center, text=popmuted] at (8.6,0.9) {soft-rock vale\\ (low flat ground)};
\draw[->, popmuted] (8.6,1.2) -- (8.4,1.55);
\node[align=center, text=popdark] at (5.2,4.1) {crest};
\draw[->, popdark] (5.2,3.85) -- (5.2,3.45);
\end{tikzpicture}
\legende{Annotated line-drawing of an oblique photograph of a cuesta: gentle dip slope, steep scarp slope, bedding traces and the soft-rock vale.}
\end{popfigure}

\textbf{Fold closures.} On vertical photographs, folded strata produce curved, concentric bands; the nose of a plunging anticline shows a \emph{V-shaped or U-shaped closure} of the outcrop pattern, with dips away from the core (anticline) or towards it (syncline).

\textbf{Fault scarps and lineaments.} A fault often appears as a straight, sharp break in the landscape: a \cle{fault scarp} separating high ground from low ground, truncated ridges, offset streams, or a line of springs and vegetation.

\begin{popfigure}
\begin{tikzpicture}[scale=1.0, every node/.style={font=\small}]
\draw[thick, popdark] (-0.4,-0.6) rectangle (10.6,4.2);
% Uplifted block (left) and downthrown block (right)
\draw[very thick, popdark] (0,2.8) -- (4.6,2.8) -- (5.0,1.2) -- (10.2,1.2);
% bedding on upthrown block
\draw[poporange] (0,2.3) -- (4.65,2.3); \draw[poporange] (0,1.85) -- (4.72,1.85);
% bedding on downthrown block (lower)
\draw[poporange] (5.05,0.9) -- (10.2,0.9); \draw[poporange] (5.1,0.5) -- (10.2,0.5);
% scarp hachures
\draw[popgreen] (4.66,2.5) -- (4.88,2.22);
\draw[popgreen] (4.74,2.1) -- (4.96,1.82);
\draw[popgreen] (4.82,1.7) -- (5.04,1.42);
% offset stream
\draw[popblue, thick] (2.0,3.6) -- (2.0,2.8);
\draw[popblue, thick] (2.0,2.8) -- (4.6,2.8);
\draw[popblue, thick, dashed] (4.6,2.8) -- (5.0,1.2);
\draw[popblue, thick] (5.0,1.2) -- (5.6,1.2) -- (5.6,0.2);
% springs along fault
\node[text=popblue] at (4.85,2.1) {\tiny $\bullet$};
\node[text=popblue] at (4.95,1.6) {\tiny $\bullet$};
% labels
\node[align=center, text=popgreen!50!black] at (7.6,3.5) {straight fault scarp};
\draw[->, popgreen!50!black] (7.2,3.25) -- (4.95,2.2);
\node[align=center, text=poporange] at (1.8,3.4) {beds on upthrown block};
\draw[->, poporange] (2.2,3.15) -- (2.4,2.35);
\node[align=center, text=poporange] at (8.6,0.2) {same beds lower: downthrow};
\draw[->, poporange] (8.2,0.45) -- (7.6,0.85);
\node[align=center, text=popblue] at (2.6,1.1) {stream offset\\ at scarp};
\draw[->, popblue] (3.0,1.35) -- (4.7,2.55);
\node[align=center, text=popblue] at (6.9,2.3) {springs along\\ fault line};
\draw[->, popblue] (6.5,2.1) -- (5.0,1.75);
\end{tikzpicture}
\legende{Line-drawing of a fault-scarp photograph: straight scarp, displaced bedding, offset stream and aligned springs.}
\end{popfigure}

\textbf{Dyke ridges.} A resistant igneous dyke crossing softer country rock forms a long, straight, narrow \emph{wall-like ridge}; a dyke softer than its host rock forms a straight trench. Both appear as lineaments.

\courssub{Identifying Landforms on Photographs}

\begin{itemize}
\item \cle{Cuesta} and \cle{hogback}: asymmetric ridges developed on dipping sedimentary strata (see the figure above).
\item \cle{Volcanic cone and crater}: a conical hill with radial drainage and a circular summit depression; dark-toned lava flows may radiate from the summit — the classic signature of \cle{Mount Cameroon}.
\item \cle{Inselberg}: an isolated steep-sided rocky hill rising abruptly from a flat erosion surface, typical of granitic and gneissic regions such as the \cle{Mandara Mountains} of the Far North.
\item \cle{River terraces}: flat, bench-like surfaces standing above the present floodplain, recording former river levels.
\item \cle{Landslides}: arcuate scars on steep slopes with hummocky, chaotic ground below — common on the steep volcanic flanks of the South-West Region during the rainy season.
\end{itemize}

\begin{savaistu}[Cameroon case studies]
\begin{itemize}
\item \textbf{Mount Cameroon}: on photographs, the volcano shows a near-perfect cone, radial drainage, dark recent lava flows cutting across older vegetated slopes, and a summit crater zone.
\item \textbf{Mandara Mountains}: photographs show bare, rounded granite inselbergs and boulder-strewn slopes rising from flat sandy plains — a textbook tropical erosion landscape.
\item \textbf{Benue trough}: around Garoua, photographs reveal flat-topped sedimentary scarps and cuestas developed on gently dipping Cretaceous sandstones, with trellis drainage in the valleys.
\end{itemize}
\end{savaistu}

\courssub{Linking Photograph, Map and Section}

The three tools of the geologist show the same reality in three ways: the \textbf{photograph} gives the perspective view, the \textbf{map} gives the plan view with symbols and contours, and the \textbf{cross-section} gives the vertical view. A cuesta appears on a photograph as an asymmetric ridge, on a map as a banded outcrop whose V-shaped boundary crossings of contour lines indicate dip, and on a section as a tilted bed between two horizontal limits. Examination questions increasingly combine all three: you may be asked to match a photograph to a map extract, or to sketch the section that explains the photographed landform. Always ask yourself: \emph{what structure seen on the photograph would produce this outcrop pattern on the map?}

\begin{exoresolu}[Interpreting an oblique photograph of the Benue scarp country]
An oblique aerial photograph shows a long, straight ridge. Its left slope is gentle and carries parallel light-and-dark bands; its right slope is steep and bare. Short streams descend the steep slope and join a long stream flowing along the flat low ground at its foot, parallel to the ridge. (a) Describe the landform. (b) Identify the structure and rock arrangement. (c) Name the drainage pattern and justify it.
\tcblower
\textbf{(a) Description.} The photograph shows an asymmetric ridge: a gentle, banded slope on one side, a steep bare slope on the other, and a longitudinal stream in the low ground parallel to the ridge crest.

\textbf{(b) Interpretation.} This is a \cle{cuesta}. The banded gentle slope is a \emph{dip slope} following the top of a resistant bed (e.g.\ sandstone); the steep slope is the \emph{scarp slope} cutting across the bedding. The strata therefore dip towards the left (towards the gentle slope). The low ground is a vale eroded in softer rock (e.g.\ shale).

\textbf{(c) Drainage.} The long stream along the vale, joined at right angles by short streams from the scarp, is characteristic of \cle{trellis drainage}, which develops where alternating hard and soft dipping strata crop out in parallel bands — exactly the situation deduced in (b).
\end{exoresolu}

\begin{exercice}[Reading a volcanic photograph]
A vertical aerial photograph of a volcanic region shows: a circular hill with a small central depression; streams radiating from the summit in all directions; a dark, rough-textured tongue extending from the summit down one flank, bordered by lighter vegetated ground. (a) Describe each feature. (b) Interpret the hill, the depression, the drainage and the dark tongue. (c) Which Cameroonian volcano could this photograph represent?
\end{exercice}

\begin{corrige}[Exercise 1]
\textbf{(a) Description:} circular conical hill; central circular depression; radial stream pattern; one dark, rough, tongue-shaped zone on a flank contrasting with lighter vegetated surroundings.
\textbf{(b) Interpretation:} the cone is a \cle{volcanic cone}; the depression is the \cle{crater}; the radial drainage reflects the outward slopes of the cone; the dark tongue is a relatively recent \cle{lava flow}, still bare of vegetation because of its youth and rough surface.
\textbf{(c)} This is typical of \cle{Mount Cameroon} (South-West Region), whose historical eruptions have produced dark recent lava flows crossing older vegetated flanks.
\end{corrige}

\begin{exercice}[From photograph to structure]
A ground photograph of a road cutting shows parallel light and dark bands inclined downwards from left to right at a moderate angle; the bands are truncated at the top of the cutting by a horizontal surface above which lie horizontal gravels. (a) Describe the photograph. (b) Identify the two structural features visible. (c) State the geological history recorded in the cutting.
\end{exercice}

\begin{corrige}[Exercise 2]
\textbf{(a) Description:} parallel inclined bands (bedding) dipping to the right, cut at the top by a horizontal surface overlain by horizontal deposits.
\textbf{(b) Interpretation:} the inclined bands are \cle{dipping sedimentary strata}; the horizontal surface separating them from the flat-lying gravels is an \cle{unconformity}.
\textbf{(c) History:} (1) deposition of the strata horizontally; (2) tilting (and possibly folding) of the strata; (3) uplift and erosion producing the horizontal erosion surface; (4) subsidence and deposition of the horizontal gravels on the eroded surface.
\end{corrige}

\begin{aretenir}[Key points of Section 11]
\begin{itemize}
\item Four types of photographs: ground, oblique aerial, vertical aerial, satellite images.
\item Interpret in four steps: \textbf{locate} $\rightarrow$ \textbf{describe landforms} $\rightarrow$ \textbf{identify rocks/structures} $\rightarrow$ \textbf{interpret processes}. Describe before interpreting.
\item Read the six clues: tone, texture, vegetation, drainage, relief, lineaments.
\item Drainage patterns diagnose structure: dendritic = homogeneous rock; trellis = dipping/folded hard--soft alternation; radial = cone or dome; rectangular = joints/faults.
\item A lineament is only a \emph{suspected} fault until confirmed by offsets, scarps, springs or crushed rock.
\item Cuestas, hogbacks, cones, inselbergs, terraces and landslides each have a distinctive photographic signature — Mount Cameroon, the Mandara inselbergs and the Benue scarps are your Cameroonian references.
\item Photograph, map and cross-section are three views of the same geology: learn to move fluently between them.
\end{itemize}
\end{aretenir}

\coursec{Section 12: Integration — Complete Map Interpretation and GCE Examination Technique}

\courssub{From Photographs Back to the Map: Why Integration Matters}

In Section 11 we learned to read the landscape itself: a photograph of Mount Cameroon's lava flows or of tilted sandstone beds in the Mamfe Basin tells us about dip, structure and history in three dimensions. The GCE examiner's favourite trick is to demand the \emph{reverse} journey: from the flat, symbolic language of a geological map back to the three-dimensional reality, and then to the fourth dimension --- \cle{time}, the geological history. No single skill is enough. A full map question tests, in one exercise, your knowledge of symbols, strike and dip, strike lines, the three-point problem, cross-section construction and the ordering of geological events. This final section trains you to combine all of these into one disciplined, timed, mark-winning workflow.

\begin{aretenir}[The one-sentence summary of the whole chapter]
A geological map is a coded message: your task is always the same --- \textbf{decode the outcrops} (symbols, contacts, topography), \textbf{reconstruct the geometry} (dip, strike lines, cross-section), then \textbf{narrate the history} (oldest event first, youngest last).
\end{aretenir}

\courssub{The Complete Map-Interpretation Workflow}

Before any timed practice, fix in your mind the logical order of operations. Skipping a step --- for example drawing a cross-section before establishing the dip --- is the root cause of most failed answers.

\begin{methode}[The seven-step workflow for any geological map]
\begin{enumerate}
\item \textbf{Read the marginal information}: title, scale, north arrow, legend, contour interval. Note the scale numerically (e.g. $1:25\,000$ means $1\ \mathrm{cm} = 250\ \mathrm{m}$).
\item \textbf{Identify the topography}: contour pattern, valleys, ridges, spot heights. Topography controls outcrop shape through the rule of V's.
\item \textbf{Identify the rock units}: list them from the legend (usually youngest at top) and locate each outcrop pattern on the map.
\item \textbf{Determine the geometry}: read dip symbols; draw \cle{strike lines} (structure contours) on at least two contacts; solve a three-point problem if dip symbols are absent; deduce dip direction and amount.
\item \textbf{Identify the structures}: folds (outcrop repetition and dip pattern), faults (displaced contacts, fault symbols), unconformities (truncated beds, basal conglomerate), intrusions (cross-cutting bodies, baked margins).
\item \textbf{Construct the cross-section} along the required line, with vertical exaggeration stated, all beds labelled, dip angles respected.
\item \textbf{Write the geological history}: a numbered list of events from oldest to youngest, each justified by a field relationship (superposition, cross-cutting, unconformity).
\end{enumerate}
\end{methode}

\begin{popfigure}
\begin{center}
\begin{tikzpicture}[
box/.style={draw=popblue, fill=popblueL, rounded corners=2pt, text width=3.1cm, align=center, minimum height=0.9cm, font=\small},
sbox/.style={draw=popgreen, fill=popgreenL, rounded corners=2pt, text width=3.1cm, align=center, minimum height=0.9cm, font=\small},
fleche/.style={-{Stealth[length=2.5mm]}, thick, popdark}]
\node[box, align=center] (a) at (0,0){\textbf{1. Margins}\\ scale, north, legend};
\node[box, align=center] (b) at (4.3,0){\textbf{2. Topography}\\ contours, valleys};
\node[box, align=center] (c) at (8.6,0){\textbf{3. Rock units}\\ outcrop patterns};
\node[sbox, align=center] (d) at (8.6,-2.1){\textbf{4. Geometry}\\ dip, strike lines, 3-point};
\node[sbox, align=center] (e) at (4.3,-2.1){\textbf{5. Structures}\\ folds, faults, intrusions};
\node[sbox, align=center] (f) at (0,-2.1){\textbf{6. Cross-section}\\ labelled, to scale};
\node[box, draw=poporange, fill=white, align=center] (g) at (4.3,-4.2){\textbf{7. Geological history}\\ oldest $\rightarrow$ youngest};
\draw[fleche] (a) -- (b);
\draw[fleche] (b) -- (c);
\draw[fleche] (c) -- (d);
\draw[fleche] (d) -- (e);
\draw[fleche] (e) -- (f);
\draw[fleche] (f.south) |- (g.west);
\end{tikzpicture}
\end{center}
\legende{The complete map-interpretation workflow. Steps 1--3 decode the map; steps 4--6 reconstruct the geometry; step 7 delivers the history.}
\end{popfigure}

\courssub{Timed Practice: A Complete GCE-Style Map Question}

The GCE Advanced Level typically allows about 30 minutes for a full map-work question. Train yourself with the following plan until it becomes automatic.

\begin{methode}[The 20-minute plan for a full GCE map question]
\begin{enumerate}
\item \textbf{Minutes 0--3}: read the whole question \emph{and} the map margins; underline the command words; note the scale and the mark allocation.
\item \textbf{Minutes 3--6}: description of relief and outcrop pattern (quick, factual sentences).
\item \textbf{Minutes 6--11}: constructions in pencil --- strike lines, three-point problem, dip calculation. Leave all construction lines visible.
\item \textbf{Minutes 11--16}: draw the cross-section; label every bed, the line of section, horizontal and vertical scales.
\item \textbf{Minutes 16--19}: write the geological history as a numbered, ordered list.
\item \textbf{Minute 20}: final check against the common-errors checklist (below).
\end{enumerate}
If the question carries more than 20 marks, scale each block proportionally --- roughly one minute per mark.
\end{methode}

\begin{exoresolu}[Worked integration: dip from a three-point problem, then history]
\textbf{Statement.} On a map at scale $1:10\,000$, three boreholes pierce the top of a sandstone bed: A ($200\ \mathrm{m}$ altitude), B ($150\ \mathrm{m}$) and C ($100\ \mathrm{m}$). A and B lie $500\ \mathrm{m}$ apart on a line bearing due east; C lies $400\ \mathrm{m}$ due south of B. (a) Find the strike and dip of the bed. (b) State, with reasons, where the bed outcrops relative to a $150\ \mathrm{m}$ contour.
\tcblower
\textbf{Solution.} (a) \emph{Strike:} the $150\ \mathrm{m}$ point on the bed occurs at B and also at a point D on line AC found by interpolation: altitude drops $100\ \mathrm{m}$ from A to C, so D is at $\tfrac{200-150}{200-100} = \tfrac{1}{2}$ of the way from A to C. Joining B to D gives a \cle{strike line} at $150\ \mathrm{m}$. A parallel strike line through C is at $100\ \mathrm{m}$. \emph{Dip:} the dip direction is perpendicular to the strike lines, towards decreasing altitude (towards C's side). On the map, the horizontal equivalent between the $150\ \mathrm{m}$ and $100\ \mathrm{m}$ strike lines is measured (say $250\ \mathrm{m}$ on the ground). Then
\[ \tan\theta = \frac{\text{vertical interval}}{\text{horizontal equivalent}} = \frac{50}{250} = 0.2 \quad\Rightarrow\quad \theta \approx 11.3^{\circ}. \]
The bed dips about $11^{\circ}$ towards the south(-east, as drawn). (b) Where the ground surface is exactly at $150\ \mathrm{m}$, the bed top coincides with the surface: the outcrop of the bed top follows the $150\ \mathrm{m}$ contour \emph{only} at the points where the $150\ \mathrm{m}$ strike line crosses that contour --- elsewhere the outcrop is found by joining crossings of strike lines with contours of equal value. \textbf{Reason:} a plane intersects the topography along lines joining points of equal elevation on both surfaces.
\end{exoresolu}

\courssub{Answering Technique: Decoding Command Words}

Examiners choose their verbs deliberately. Each command word implies a different type of answer, and answering a ``describe'' question with an explanation (or vice versa) wastes time and loses marks.

\begin{center}
\begin{tabularx}{\linewidth}{|l|X|X|}
\hline
\rowcolor{popblueL} \textbf{Command word} & \textbf{What the examiner wants} & \textbf{Typical marks signal} \\
\hline
\cle{Describe} & Facts only: what you see on the map (outcrop pattern, dip values, relationships). No interpretation. & 1 mark per correct observation \\
\hline
\cle{Explain} & Reasons and mechanisms: \emph{why} the outcrop is V-shaped, \emph{why} the beds are repeated. & 1 mark per linked reason \\
\hline
\cle{Deduce} & A conclusion drawn \emph{from evidence on the map}: state the conclusion \textbf{and} the evidence. & conclusion + justification \\
\hline
\cle{Construct} & A drawing: cross-section or strike lines, to scale, fully labelled, in pencil. & accuracy + labels + scale \\
\hline
\cle{Calculate} & Numerical working shown: gradient, dip angle, true thickness, real distance from the scale. & method + correct value + unit \\
\hline
\end{tabularx}
\end{center}

\begin{attention}[``Describe'' is not ``explain'']
In ``describe the outcrop pattern'', writing ``the beds are folded'' earns nothing if you never said \emph{what you see}: ``the sandstone outcrop is repeated symmetrically on both sides of the valley''. Give the observation first; interpretation belongs to ``deduce'' or ``explain'' questions.
\end{attention}

\courssub{Mark-Scheme Awareness: How Points Are Allocated}

Map-work mark schemes are remarkably predictable. Knowing the anatomy of the marks lets you check, before handing in, that every mark has been ``offered'' to the examiner.

\begin{itemize}
\item \textbf{Description (4--6 marks)}: one mark per distinct correct observation --- relief, drainage, number of rock units, outcrop widths, dip directions.
\item \textbf{Construction (4--6 marks)}: correct strike lines (1--2), correct dip direction (1), correct dip value within tolerance (1), correct use of scale (1).
\item \textbf{Cross-section (5--8 marks)}: correct line and scale (1), topography drawn (1), beds at correct angles (2), all units labelled (1--2), structures shown correctly (1--2).
\item \textbf{Geological history (4--6 marks)}: one mark per event \emph{in the correct chronological position}; an event correctly identified but wrongly ordered usually scores zero for that line.
\end{itemize}

\begin{attention}[The five most frequent mark-losing errors in map work]
\begin{enumerate}
\item \textbf{Wrong dip direction}: dipping beds \emph{towards} the higher strike line instead of the lower one. Always check: dip goes towards decreasing elevation.
\item \textbf{Forgotten or wrong scale}: distances measured in centimetres and quoted as metres; vertical exaggeration not stated on the section.
\item \textbf{Unlabelled cross-section}: no line of section, no scale, unnamed beds --- up to half the section marks vanish.
\item \textbf{Unordered history}: events listed randomly, or youngest first when the question demands oldest first.
\item \textbf{Erased construction lines}: the examiner cannot award method marks for work that is no longer visible.
\end{enumerate}
\end{attention}

\begin{attention}[Exam technique: pencil and visible constructions]
All constructions --- strike lines, interpolation points, projection lines from map to section --- must be done in \textbf{pencil}, \textbf{labelled} (e.g. ``strike line 150 m''), and \textbf{left visible}. Never erase construction lines: they carry method marks. Only the final cross-section lines may be inked or darkened if time permits.
\end{attention}

\courssub{Beyond the Syllabus: Digital Geological Maps and GIS}

\begin{savaistu}[Extra --- beyond syllabus: GIS in modern geology]
A \cle{GIS} (Geographic Information System) geological database stores each rock unit, fault and dip measurement as a digital layer linked to its exact coordinates, so maps can be queried, updated and combined with satellite images instantly. Recent examiners' reports mention digital maps favourably: knowing that agencies such as Cameroon's mining and geological services increasingly produce GIS-based maps shows you understand where the discipline is going --- but the paper-and-pencil skills of this chapter remain exactly what the GCE examines.
\end{savaistu}

\courssub{Revision Strategy and Self-Assessment}

Revise this chapter \emph{actively}: never re-read notes when you could redraw a strike line or re-solve a three-point problem. Work through past GCE map questions under timed conditions, then mark yourself with the grid below. Any row answered ``not yet'' defines your next revision session.

\begin{popfigure}
\begin{center}
\begin{tabular}{|p{6.2cm}|c|c|c|}
\hline
\rowcolor{popblueL} \textbf{Skill (chapter GEO03)} & \textbf{Secure} & \textbf{Shaky} & \textbf{Not yet} \\
\hline
Identify map types and read all marginal information & $\square$ & $\square$ & $\square$ \\
\hline
Use direction, bearings and the map scale correctly & $\square$ & $\square$ & $\square$ \\
\hline
Recognise all conventional symbols (dip, faults, folds) & $\square$ & $\square$ & $\square$ \\
\hline
State strike, dip, thickness and the rule of V's & $\square$ & $\square$ & $\square$ \\
\hline
Draw strike lines and solve the three-point problem & $\square$ & $\square$ & $\square$ \\
\hline
Construct a labelled cross-section (simple and folded) & $\square$ & $\square$ & $\square$ \\
\hline
Show faults, unconformities and intrusions on sections & $\square$ & $\square$ & $\square$ \\
\hline
Describe structures from outcrop patterns alone & $\square$ & $\square$ & $\square$ \\
\hline
Write an ordered, justified geological history & $\square$ & $\square$ & $\square$ \\
\hline
Interpret geological photographs and link them to maps & $\square$ & $\square$ & $\square$ \\
\hline
\end{tabular}
\end{center}
\legende{Self-assessment checklist for chapter GEO03. Re-test every ``shaky'' row with a timed past-paper question before the examination.}
\end{popfigure}

\begin{exercice}[Final integration exercise]
On the GCE-style map provided by your teacher (scale $1:25\,000$): (a) describe the relief and the outcrop pattern of each rock unit (4 marks); (b) construct strike lines on the upper contact of the limestone and calculate its dip (4 marks); (c) draw a labelled cross-section along line XY (6 marks); (d) deduce the type of fold present, justifying your answer (2 marks); (e) give the geological history of the area, oldest event first (4 marks). Time allowed: 30 minutes.
\end{exercice}

\begin{corrige}[Final integration exercise]
(a) Award one mark per correct observation: relief trend from contours, width and shape of each outcrop, repetition of units, relationship of outcrops to valleys. (b) Strike lines join points of equal elevation on the contact; dip direction is towards the lowest strike line; $\tan\theta = \dfrac{\text{contour interval}}{\text{horizontal spacing of strike lines}}$ --- method 2 marks, value within $\pm 2^{\circ}$ 2 marks. (c) Marks: topography along XY (1), beds drawn at the calculated dip (2), fold drawn symmetrically (1), all units and both scales labelled (2). (d) If dips face each other and the oldest rocks occupy the core, it is a \cle{syncline}; dips away from each other with oldest core indicate an \cle{anticline} --- conclusion 1 mark, evidence 1 mark. (e) One mark per event in correct order: deposition of oldest unit $\rightarrow$ deposition of younger units $\rightarrow$ folding $\rightarrow$ (faulting/intrusion if present) $\rightarrow$ erosion to present surface. Any event out of order loses its mark.
\end{corrige}

\begin{aretenir}[Chapter GEO03 --- the essentials in five lines]
\begin{enumerate}
\item Read the margins first: scale, north, legend, contours.
\item Strike lines are structure contours: dip is towards decreasing values, $\tan\theta = \dfrac{\text{vertical interval}}{\text{horizontal equivalent}}$.
\item Every cross-section: line, topography, correct dips, labels, two scales --- in pencil, constructions visible.
\item History is always oldest first, each event justified by superposition, cross-cutting or unconformity.
\item Match your answer to the command word, and spend roughly one minute per mark.
\end{enumerate}
You have now completed Map Work --- the most heavily weighted practical skill of the Upper Sixth Geology paper. Practise one full map per week, and this chapter will become your most reliable source of marks.
\end{aretenir}

\newpage

\coursec{Integration activity}

\coursec{Integration — Complete Map Interpretation and GCE Examination Technique}

\courssub{Aims of the activity}

This integration activity closes the chapter on \cle{map work} by mobilising, in one single authentic problem, every skill developed in Sections 1 to 11:

\begin{itemize}
\item read the \cle{scale} of a map and convert map distances into ground distances (Section 3);
\item recognise \cle{conventional symbols} and identify rock units, faults and intrusions (Section 4);
\item draw \cle{strike lines} and solve the \cle{three-point problem} to determine the dip of a bed (Section 6);
\item construct a \cle{geological cross-section} along a chosen line (Sections 7 and 8);
\item \cle{describe the structures} shown on the map and reconstruct the \cle{geological history} of the area (Sections 9 and 10);
\item interpret a \cle{geological photograph} and relate landforms to the underlying geology (Section 11);
\item apply all of the above to a real decision: siting a \cle{borehole} for village water supply, and defending the choice in a short technical report.
\end{itemize}

The activity is designed for teamwork (groups of 3 to 4 learners), followed by an oral defence of each team's decision, exactly as a geologist defends a siting report before a technical commission.

\courssub{Situation}

\textbf{The Bafut Water Supply Project.}\par

The village of \textbf{Nkwen--Bafut} (Mezam Division, North West Region) lies at the foot of a prominent escarpment. The Divisional Delegate of Water and Energy wishes to drill a new borehole to supply the village and to lay a pipeline from the borehole to a storage tank at the market square. Your school geological team has been contracted to advise on the best drilling site. You are given the simplified geological map below (Figure 1) and a photograph of the landscape taken from the plain, looking north towards the escarpment.

\begin{popfigure}
\begin{tikzpicture}[scale=1.05]
% grid
\draw[step=1cm, gray!35, very thin] (0,0) grid (10,7);
% north arrow
\draw[->, thick, popdark] (9.4,6.2) -- (9.4,6.9);
\node[above, font=\footnotesize] at (9.4,6.9) {N};
% scale bar
\draw[thick] (0.4,0.35) -- (2.4,0.35);
\draw[thick] (0.4,0.28) -- (0.4,0.42);
\draw[thick] (2.4,0.28) -- (2.4,0.42);
\node[below, font=\footnotesize] at (1.4,0.25) {0 \quad 1 km};
% boundaries dipping beds (strike N-S, dip E)
\draw[thick, popblue] (2.2,0.6) -- (2.2,6.6);
\draw[thick, popblue] (4.6,0.6) -- (4.6,6.6);
\draw[thick, popblue] (7.0,0.6) -- (7.0,6.6);
% units
\node[font=\footnotesize, popdark] at (1.1,5.6) {Alluvium};
\node[font=\footnotesize, popdark] at (3.4,5.6) {Shale};
\node[font=\footnotesize, popdark] at (5.8,5.6) {Sandstone};
\node[font=\footnotesize, popdark] at (5.8,5.1) {(aquifer)};
\node[font=\footnotesize, popdark] at (8.5,5.6) {Shale};
% fault
\draw[very thick, poporange, dashed] (0.6,1.2) -- (9.6,3.4);
\node[font=\footnotesize, poporange, rotate=13] at (5.2,2.05) {Fault F};
% dyke
\draw[very thick, poppink] (8.2,0.6) -- (8.8,6.6);
\node[font=\footnotesize, poppink, rotate=80] at (8.55,3.6) {Dyke};
% boreholes
\fill[popgreen] (5.0,4.6) circle (0.09); \node[right, font=\footnotesize] at (5.05,4.6) {BH1 (h = 120 m)};
\fill[popgreen] (6.2,3.0) circle (0.09); \node[right, font=\footnotesize] at (6.25,3.0) {BH2 (h = 60 m)};
\fill[popgreen] (5.6,1.6) circle (0.09); \node[right, font=\footnotesize] at (5.65,1.6) {BH3 (h = 90 m)};
% village and market
\fill[black] (1.2,2.2) rectangle (1.5,2.5); \node[below, font=\footnotesize] at (1.35,2.2) {Village};
\fill[black] (3.0,1.0) rectangle (3.3,1.3); \node[below, font=\footnotesize] at (3.15,1.0) {Market};
% proposed pipeline
\draw[thick, poppurple, dotted] (1.35,2.35) -- (3.15,1.15);
% escarpment hachures
\foreach \y in {1.0,1.8,2.6,3.4,4.2,5.0,5.8} {\draw[popgold, thick] (7.0,\y) -- (7.35,\y);}
\node[font=\footnotesize, popgold, rotate=90] at (7.35,3.6) {Escarpment};
\end{tikzpicture}
\legende{Simplified geological map of the Nkwen--Bafut area. Scale: 1 cm = 0.5 km (1:50\,000). Grid squares are 1 cm $\times$ 1 cm. The sandstone unit is the water-bearing bed (aquifer).}
\end{popfigure}

\textbf{Data supplied with the map.}
\begin{itemize}
\item The map scale is \textbf{1:50\,000} (1 cm on the map = 0.5 km on the ground).
\item The shale--sandstone and sandstone--shale boundaries run approximately \textbf{north--south}; the beds \textbf{dip eastwards}.
\item Three existing boreholes (BH1, BH2, BH3) all lie on the sandstone outcrop and each struck the \emph{top} of the water-bearing sandstone at a known altitude $h$ (given on the map). These three points define the attitude of the top of the aquifer.
\item A fault F (dashed line) cuts the area; field observations show the northern block has moved \textbf{down} relative to the southern block (normal fault).
\item A basaltic \textbf{dyke} cuts across the eastern shale; it is barren (no water).
\item The escarpment in the east corresponds to the resistant sandstone; the vales are carved in the soft shale.
\item The photograph (viewed in class) shows a steep, wooded scarp face to the east and a gentle dip slope descending westwards into a cultivated vale.
\end{itemize}

\courssub{Tasks}

\begin{enumerate}
\item \textbf{Scale and distances.} Measure on the map the length of the proposed pipeline route (village to market). Convert this map distance into a ground distance in kilometres, using the scale 1:50\,000.
\item \textbf{Three-point problem.} Using the three boreholes BH1 (120 m), BH2 (60 m) and BH3 (90 m), construct strike lines on the top of the aquifer sandstone and determine: (a) the direction of strike; (b) the direction of dip; (c) the angle of dip. Show all construction lines.
\item \textbf{Cross-section.} Construct a geological cross-section along the east--west line passing through BH2. Show the alluvium, the shale, the sandstone aquifer, the eastern shale, the fault F, the dyke and the escarpment. Use a vertical scale equal to the horizontal scale (no vertical exaggeration).
\item \textbf{Structural description.} Describe in correct geological language the dip of the strata, the fault (type, downthrow side, effect on outcrops) and the dyke (nature, age relative to the beds).
\item \textbf{Geological history.} Write the chronological succession of events that produced the present geology of the area, from the oldest to the youngest event.
\item \textbf{Photograph interpretation.} Explain how the scarp-and-vale landscape seen on the photograph confirms the geology shown on the map. Name the type of relief and the landform produced.
\item \textbf{Siting decision and report.} Two sites are proposed for the new borehole: \textbf{Site A}, on the sandstone outcrop just west of the escarpment, and \textbf{Site B}, in the alluvium of the vale, 1 km west of the sandstone outcrop. Using your cross-section and the dip calculated in Task 2, estimate the depth of the top of the aquifer beneath each site, then recommend one site in a one-page report to the Divisional Delegate of Water and Energy, justifying your choice (depth to water, cost, risk from the fault and the dyke).
\end{enumerate}

\courssub{Model answer}

\textbf{Task 1 — Scale and distances.}\par
The pipeline route measures about \textbf{4.4 cm} on the map. Since 1 cm represents 0.5 km:
\[
d = 4.4 \times 0.5 = 2.2\ \mathrm{km}.
\]
The pipeline will be approximately \textbf{2.2 km} long.

\textbf{Task 2 — Three-point problem.}\par
\begin{methode}[Solution]
\begin{enumerate}
\item Join BH1 (120 m, highest) to BH2 (60 m, lowest). Along this line, interpolate the position of the 90 m point: it lies halfway between BH1 and BH2, since 90 m is halfway between 120 m and 60 m.
\item Join this interpolated 90 m point to BH3 (90 m): this is the \textbf{90 m strike line}. It runs approximately north--south, confirming the map boundaries.
\item From BH2 (60 m), drop a perpendicular to the strike line: this is the \textbf{dip direction}, pointing \textbf{east}.
\item Measure the horizontal distance $d$ between the 90 m strike line and BH2 (about 1.2 cm on the map, i.e. 600 m on the ground). The height difference is $90 - 60 = 30$ m.
\end{enumerate}
\end{methode}
\[
\tan \theta = \frac{30}{600} = 0.05 \quad\Rightarrow\quad \theta \approx 2.86^{\circ} \approx 3^{\circ}\ \text{towards the east}.
\]
The aquifer dips about \textbf{3$^{\circ}$ east} (a gentle dip). Strike: N--S; dip: 3$^{\circ}$ E.

\textbf{Task 3 — Cross-section.}\par
From west to east the section shows: a thin veneer of alluvium in the vale; the western shale; the \textbf{sandstone aquifer} dipping gently east at about 3$^{\circ}$, its top passing through the altitudes given by BH1, BH2 and BH3; the eastern shale; the \textbf{fault F} throwing the northern block down (shown by the offset of the bed boundaries); the vertical \textbf{basalt dyke} cutting the eastern shale; and the \textbf{escarpment} at the sandstone--eastern shale boundary, where the resistant sandstone forms the scarp face.

\begin{popfigure}
\begin{tikzpicture}[scale=1.0]
% Horizontal scale: 1 cm = 0.5 km, vertical scale same (no exaggeration)
% We draw a simplified cross-section along east-west through BH2 (x=6.2 on map)
% Map x from 0 to 10 cm -> 0 to 5 km. BH2 at x=6.2 cm -> 3.1 km from west edge.
% We'll draw from x=0 to x=10 cm (0 to 5 km) at y=0 to 200 m approx.
% But for simplicity, schematic:
\draw[thick] (0,0) -- (10,0); % baseline
% Alluvium (west)
\draw[fill=popblueL!30] (0,0) -- (2.5,0) -- (2.5,0.3) -- (0,0.3) -- cycle;
\node[font=\footnotesize] at (1.25,0.15) {Alluvium};
% Western shale
\draw[fill=popgray!30] (2.5,0) -- (5.0,0) -- (5.0,1.0) -- (2.5,1.0) -- cycle;
\node[font=\footnotesize] at (3.75,0.5) {Shale};
% Sandstone aquifer dipping east at 3 deg
% Top of sandstone at BH2 (x=6.2 cm map = 3.1 km) altitude 60 m.
% We set vertical scale: 1 cm = 20 m? Let's do schematic.
% Actually, for a true scale cross-section, horizontal 1 cm = 0.5 km = 500 m, vertical 1 cm = 500 m.
% But here we just show schematic.
\draw[fill=popyellow!40] (5.0,0) -- (7.5,0) -- (7.5,0.8) -- (5.0,0.8) -- cycle;
\node[font=\footnotesize] at (6.25,0.4) {Sandstone (aquifer)};
% Eastern shale
\draw[fill=popgray!30] (7.5,0) -- (10,0) -- (10,1.2) -- (7.5,1.2) -- cycle;
\node[font=\footnotesize] at (8.75,0.6) {Shale};
% Fault F (dashed) - approximate position
\draw[dashed, thick, poporange] (4.0,0) -- (4.0,1.2);
\node[font=\footnotesize, poporange] at (4.0,1.3) {Fault F};
% Dyke (vertical)
\draw[thick, poppink] (9.0,0) -- (9.0,1.2);
\node[font=\footnotesize, poppink] at (9.0,1.3) {Dyke};
% Escarpment at sandstone/shale boundary (x=7.5)
\draw[thick, popgold] (7.5,0) -- (7.5,0.8);
% Dip arrow
\draw[->, thick, popblue] (6.25,0.8) -- (6.25,1.2);
\node[above, font=\footnotesize, popblue] at (6.25,1.2) {Dip 3$^{\circ}$ E};
% BH2 marker
\draw[popgreen, thick] (6.2,0) -- (6.2,0.8);
\node[below, font=\footnotesize, popgreen] at (6.2,0) {BH2};
% Scale bar
\draw[thick] (0.5,-0.5) -- (1.5,-0.5);
\node[below, font=\footnotesize] at (1.0,-0.55) {1 km (horiz.)};
\draw[thick] (0.5,-0.5) -- (0.5,-0.6);
\draw[thick] (1.5,-0.5) -- (1.5,-0.6);
\end{tikzpicture}
\legende{Schematic geological cross-section (west–east through BH2) with no vertical exaggeration. The gentle eastward dip of the sandstone aquifer is shown.}
\end{popfigure}

\textbf{Task 4 — Structural description.}\par
The strata form a \textbf{homoclinal (monoclinal) sequence} dipping uniformly at about 3$^{\circ}$ to the east; outcrop boundaries run N--S, parallel to the strike. Fault F is a \textbf{normal dip-slip fault} trending WNW--ESE; the northern block is downthrown, so the bed boundaries on the northern side are displaced eastwards relative to those on the southern side (right-lateral offset on the map). The dyke is a \textbf{discordant, near-vertical basic intrusion} cutting across the bedding; it is therefore \textbf{younger than all the beds it cuts} (cross-cutting relationship).

\textbf{Task 5 — Geological history.}\par
\begin{enumerate}
\item Deposition of the \textbf{lower shale} (oldest sediment).
\item Deposition of the \textbf{sandstone} (future aquifer) — a change to higher-energy conditions.
\item Deposition of the \textbf{upper shale}.
\item \textbf{Tilting} of the whole sequence towards the east (gentle tectonic dip).
\item \textbf{Faulting} along F (normal fault, tensional regime).
\item Intrusion of the \textbf{basalt dyke}.
\item \textbf{Erosion}: differential erosion carved the scarp in the resistant sandstone and the vales in the soft shale; recent \textbf{alluvium} was deposited in the vale (youngest event).
\end{enumerate}

\textbf{Task 6 — Photograph interpretation.}\par
The photograph shows a steep, wooded scarp face to the east and a gentle slope descending westwards into a cultivated vale. This is a \textbf{cuesta}: an asymmetric ridge produced by differential erosion of gently dipping alternating hard and soft strata. The \textbf{scarp slope} cuts across the bedding of the resistant sandstone; the \textbf{dip slope} follows the top of the sandstone at about 3$^{\circ}$. The vales coincide with the easily eroded shale. The photograph therefore confirms the dip direction (east), the gentle dip, and the lithological alternation read on the map.

\textbf{Task 7 — Siting decision (summary of the report).}\par
\emph{To the Divisional Delegate of Water and Energy, Mezam Division.}\par
At \textbf{Site A} (on the outcrop), the aquifer is at or near the surface: drilling is cheap, but the site lies close to the fault F and to the barren dyke, and the outcrop is a recharge zone where water may drain away downslope; yields are uncertain. At \textbf{Site B} (in the vale, 1 km west of the outcrop), the top of the aquifer lies at a depth given by the dip:
\[
\text{depth} = 1000 \times \tan 3^{\circ} \approx 52.4\ \mathrm{m},
\]
plus a few metres of alluvium, i.e. about \textbf{55--60 m} of drilling. The aquifer here is \textbf{confined beneath the shale}, which protects it from pollution and gives artesian pressure; the site is far from the fault and the dyke. We therefore recommend \textbf{Site B}: slightly deeper drilling, but a safer, better protected and more reliable supply for the village.

\begin{attention}[Common mistakes in this activity]
Confusing map distance with ground distance (always multiply by the scale); placing the 90 m interpolated point at the wrong fraction of BH1--BH2; measuring the dip along a line that is not perpendicular to the strike line; drawing the dyke as a bed (it is vertical and discordant); and recommending a site without quantifying the drilling depth from the calculated dip.
\end{attention}

\newpage

\coursec{Methods and worked examples}

\courssub{Method 1 — Converting and Using Map Scales}

\begin{methode}[Working with map scales]
\begin{enumerate}
\item \textbf{Identify the scale} given on the map (e.g. $1:25\,000$ or $1:50\,000$). A scale $1:n$ means that 1 unit measured on the map represents $n$ identical units on the ground.
\item \textbf{Map distance $\rightarrow$ ground distance:} multiply the measured map distance by $n$, then convert to metres or kilometres.
\[ d_{\text{ground}} = d_{\text{map}} \times n \]
\item \textbf{Ground distance $\rightarrow$ map distance:} convert the ground distance into the same unit as the map measurement, then divide by $n$.
\[ d_{\text{map}} = \frac{d_{\text{ground}}}{n} \]
\item \textbf{Areas:} areas scale with the \emph{square} of the linear scale: $A_{\text{ground}} = A_{\text{map}} \times n^{2}$.
\item \textbf{Reflex:} always check units before computing ($1\ \mathrm{km} = 1000\ \mathrm{m} = 100\,000\ \mathrm{cm}$).
\end{enumerate}
\end{methode}

\begin{exoresolu}[Scale conversion on a $1:25\,000$ map]
On a geological map at scale $1:25\,000$, the outcrop width of a sandstone bed measured perpendicular to the boundaries is $1.6\ \mathrm{cm}$. (a) Calculate the true ground width of the outcrop. (b) Two boreholes are $340\ \mathrm{m}$ apart on the ground; what is their separation on the map?
\tcblower
\textbf{(a) Ground width.} The scale $1:25\,000$ means $1\ \mathrm{cm}$ on the map represents $25\,000\ \mathrm{cm}$ on the ground.
\[ d_{\text{ground}} = 1.6\ \mathrm{cm} \times 25\,000 = 40\,000\ \mathrm{cm} = 400\ \mathrm{m}. \]
The true outcrop width is $\cle{400\ \mathrm{m}}$.

\textbf{(b) Map separation.} Convert $340\ \mathrm{m}$ to centimetres: $340\ \mathrm{m} = 34\,000\ \mathrm{cm}$. Then
\[ d_{\text{map}} = \frac{34\,000}{25\,000} = 1.36\ \mathrm{cm}. \]
The boreholes plot $\cle{1.36\ \mathrm{cm}}$ apart on the map.
\end{exoresolu}

\courssub{Method 2 — Direction, Bearings and Orientation}

\begin{methode}[Reading direction and orientation on a map]
\begin{enumerate}
\item \textbf{Locate the north arrow} (grid north). All bearings are measured \emph{clockwise from north}, from $000^{\circ}$ to $360^{\circ}$.
\item \textbf{Direction of a line} (e.g. a fault trace): place the protractor centre on the line with $0^{\circ}$ pointing north; read the clockwise angle. A line has two possible readings differing by $180^{\circ}$; quote the one requested (or both).
\item \textbf{Strike} of a bed is quoted as a bearing (e.g. N $040^{\circ}$ E, or simply $040^{\circ}$); the \textbf{dip direction} is always perpendicular to the strike, on the side towards which the beds descend.
\item \textbf{Reflex:} strike $\perp$ dip direction. If the strike is $040^{\circ}$, the dip direction is $040^{\circ} + 90^{\circ} = 130^{\circ}$ (or $310^{\circ}$ on the opposite side); the map evidence (younging, rule of V's, dip symbols) tells you which of the two to choose.
\item \textbf{Grid references:} read eastings first, then northings ("along the corridor, then up the stairs").
\end{enumerate}
\end{methode}

\begin{exoresolu}[Strike and dip direction of a limestone bed]
On a map, the boundary of a limestone bed runs in a straight line with bearing $060^{\circ}$ across flat ground. The bed dips towards the south-east. State (a) the strike of the bed, (b) the dip direction, (c) a complete dip--strike notation if the dip angle is $30^{\circ}$.
\tcblower
\textbf{(a)} The strike is the bearing of a horizontal line in the bedding plane; on flat ground a straight boundary is such a line: strike $= \cle{060^{\circ}}$ (equivalently $240^{\circ}$).

\textbf{(b)} The dip direction is perpendicular to the strike, on the south-east side:
\[ 060^{\circ} + 90^{\circ} = 150^{\circ}. \]
Dip direction $= \cle{150^{\circ}}$ (south-east), consistent with the statement.

\textbf{(c)} Full notation: $\cle{30^{\circ} \rightarrow 150^{\circ}}$ (dip amount $\rightarrow$ dip direction), with strike $060^{\circ}/240^{\circ}$. In field-book convention one writes: strike N $60^{\circ}$ E, dip $30^{\circ}$ SE.
\end{exoresolu}

\courssub{Method 3 — True Thickness and Dip from Outcrop Width}

\begin{methode}[Calculating dip and true thickness]
\begin{enumerate}
\item \textbf{Draw a cross-section} perpendicular to the strike, showing the ground surface and the two boundaries of the bed.
\item \textbf{On horizontal ground}, if $w$ is the outcrop width (measured perpendicular to strike) and $\theta$ the dip, then the \cle{true thickness} $t$ (measured perpendicular to the bedding planes) and the vertical thickness $h$ are:
\[ t = w\,\sin\theta, \qquad h = w\,\tan\theta. \]
\item \textbf{On sloping ground}, correct for the slope: if the ground slopes in the dip direction, the effective outcrop width changes; construct the section graphically rather than forcing a formula.
\item \textbf{Reflex:} for vertical beds ($\theta = 90^{\circ}$), $t = w$; horizontal beds have an outcrop width controlled only by topography.
\item Always state units and check that $t \leq w$ and $t \leq h$.
\end{enumerate}
\end{methode}

\begin{exoresolu}[Dip and thickness of a sandstone bed]
On flat ground, a sandstone bed outcrops over a width of $250\ \mathrm{m}$ measured perpendicular to strike. Its dip is $40^{\circ}$. Calculate (a) the true thickness of the bed, (b) its vertical thickness. Take $\sin 40^{\circ} = 0.643$, $\cos 40^{\circ} = 0.766$, $\tan 40^{\circ} = 0.839$.
\tcblower
\textbf{(a) True thickness.} The true thickness $t$ is measured perpendicular to bedding:
\[ t = w \sin\theta = 250 \times 0.643 = \cle{160.75\ \mathrm{m} \approx 161\ \mathrm{m}}. \]

\textbf{(b) Vertical thickness.} The vertical thickness $h$ satisfies $h = w \tan\theta$:
\[ h = 250 \times 0.839 = \cle{209.75\ \mathrm{m} \approx 210\ \mathrm{m}}. \]
\textbf{Check:} $t < w$ and $t < h$, as expected: the true (perpendicular) thickness is always the shortest distance between the two bedding planes. A common error is to use $\cos\theta$ instead of $\sin\theta$; remember that as the dip increases, the outcrop width for a given thickness increases, so $t = w\sin\theta$ behaves correctly ($\theta = 90^{\circ} \Rightarrow t = w$).
\end{exoresolu}

\courssub{Method 4 — Strike Lines (Structure Contours) and Apparent Dip}

\begin{methode}[Drawing and using strike lines]
\begin{enumerate}
\item \textbf{Definition reflex:} a \cle{strike line} (structure contour) joins points of equal altitude on a bedding plane; for a planar bed it is straight and parallel to the strike.
\item \textbf{To draw strike lines} from points of known height on the same plane: join the highest and lowest points, divide the line proportionally to find the position of the intermediate contour value, then draw parallels through points of equal height.
\item \textbf{Spacing:} equal contour intervals give equally spaced strike lines on a planar bed. The \cle{dip} is obtained from
\[ \tan\theta = \frac{\text{contour interval}}{\text{horizontal spacing of strike lines}}. \]
\item \textbf{Apparent dip:} in a cross-section that is not perpendicular to the strike, the observed dip $\theta'$ is smaller than the true dip $\theta$:
\[ \tan\theta' = \tan\theta \, \cos\beta \]
where $\beta$ is the angle between the section line and the true dip direction. (Equivalently, $\tan\theta' = \tan\theta\,\sin\gamma$ if $\gamma$ is the angle between the section line and the strike.)
\item \textbf{Reflex:} apparent dip is always $\leq$ true dip; it equals the true dip only when the section is perpendicular to the strike ($\beta = 0^{\circ}$).
\end{enumerate}
\end{methode}

\begin{exoresolu}[Dip from strike-line spacing and apparent dip]
Strike lines drawn on a bedding plane are spaced $200\ \mathrm{m}$ apart horizontally for a contour interval of $100\ \mathrm{m}$. (a) Calculate the true dip. (b) A cross-section is drawn at $30^{\circ}$ to the true dip direction; calculate the apparent dip shown on that section. Take $\tan 27^{\circ} \approx 0.510$; $\cos 30^{\circ} = 0.866$; $\tan 24^{\circ} \approx 0.445$.
\tcblower
\textbf{(a) True dip.}
\[ \tan\theta = \frac{\text{contour interval}}{\text{spacing}} = \frac{100}{200} = 0.5 \;\Rightarrow\; \theta \approx \cle{27^{\circ}}. \]

\textbf{(b) Apparent dip.} The section makes $\beta = 30^{\circ}$ with the true dip direction:
\[ \tan\theta' = \tan\theta \times \cos\beta = 0.5 \times 0.866 = 0.433 \;\Rightarrow\; \theta' \approx \cle{23^{\circ}}. \]
\textbf{Interpretation:} on the oblique section the beds appear to dip at only about $23^{\circ}$ instead of $27^{\circ}$. This is why, in the GCE examination, a cross-section must always be drawn \emph{perpendicular to the strike} unless the question states otherwise; otherwise the dips are distorted.
\end{exoresolu}

\courssub{Method 5 — The Three-Point Problem}

\begin{methode}[Solving the three-point problem]
\begin{enumerate}
\item \textbf{Plot} the three points $A$, $B$, $C$ (boreholes or outcrops) where the same bed is met at known heights; identify the highest ($H$), lowest ($L$) and intermediate ($M$).
\item \textbf{Join $H$ to $L$.} The height of $M$ lies between those of $H$ and $L$, so a point $M'$ of the same height as $M$ exists on the line $HL$. Locate $M'$ by proportional division:
\[ \frac{HM'}{HL} = \frac{h_{H} - h_{M}}{h_{H} - h_{L}} \quad \text{(map distances and heights used consistently).} \]
\item \textbf{Join $M$ to $M'$:} this is a strike line. Draw parallels to it for other contour values — the \cle{strike} is the bearing of these lines.
\item \textbf{Dip direction:} perpendicular to the strike lines, towards decreasing heights. \textbf{Dip amount:} $\tan\theta = \dfrac{\text{contour interval}}{\text{strike-line spacing}}$.
\item \textbf{Reflex:} check that the dip direction points from the highest towards the lowest point.
\end{enumerate}
\end{methode}

\begin{exoresolu}[Three boreholes on a coal seam]
Three boreholes $A$, $B$, $C$ drilled on horizontal ground meet the top of a coal seam at heights (above sea level) of $300\ \mathrm{m}$, $200\ \mathrm{m}$ and $100\ \mathrm{m}$ respectively. $A$, $B$, $C$ form an equilateral triangle of side $400\ \mathrm{m}$ with $AB$ running east--west ($A$ west of $B$) and $C$ south of $AB$. Determine the strike, dip direction and dip of the seam. Take $\tan 27^{\circ} \approx 0.5$.
\tcblower
\textbf{Step 1 — Rank the heights:} $A = 300\ \mathrm{m}$ (highest), $C = 100\ \mathrm{m}$ (lowest), $B = 200\ \mathrm{m}$ (intermediate).

\textbf{Step 2 — Locate the point $B'$ on $AC$ at height $200\ \mathrm{m}$:} the height drops from $300\ \mathrm{m}$ at $A$ to $100\ \mathrm{m}$ at $C$, a drop of $200\ \mathrm{m}$ over the full length $AC = 400\ \mathrm{m}$. A drop of $100\ \mathrm{m}$ (to reach $200\ \mathrm{m}$) therefore occurs at half the distance: $B'$ is the \cle{midpoint of $AC$}.

\textbf{Step 3 — Strike:} join $B$ to $B'$: both points lie at $200\ \mathrm{m}$, so $BB'$ is the $200\ \mathrm{m}$ strike line. Set coordinates $A = (0,0)$, $B = (400,0)$, $C = (200, -200\sqrt{3})$ (with $y$ positive northward). Then $B' = (100, -100\sqrt{3})$, and the line $BB'$ makes an angle of $30^{\circ}$ with the east--west line, descending towards the south-west. Its bearing is therefore $060^{\circ}/240^{\circ}$: strike $= \cle{060^{\circ}}$ (equivalently $240^{\circ}$).

\textbf{Step 4 — Dip direction:} perpendicular to the strike, on the side of decreasing height (towards $C$, the lowest point): dip direction $= 060^{\circ} + 90^{\circ} = \cle{150^{\circ}}$ (towards the south-east).

\textbf{Step 5 — Dip amount:} the $200\ \mathrm{m}$ strike line passes through $B$ and $B'$; the $100\ \mathrm{m}$ strike line is the parallel through $C$. The spacing is the perpendicular distance from $C$ to the line $BB'$. With the coordinates above, this perpendicular distance is exactly $200\ \mathrm{m}$. Hence
\[ \tan\theta = \frac{\text{contour interval}}{\text{spacing}} = \frac{100}{200} = 0.5 \;\Rightarrow\; \theta \approx \cle{27^{\circ}}. \]
\textbf{Answer:} strike $060^{\circ}/240^{\circ}$, dip $\approx 27^{\circ}$ towards $150^{\circ}$ (SE). \textbf{Check:} the dip direction points from high ($A$, $300\ \mathrm{m}$) towards low ($C$, $100\ \mathrm{m}$) — correct.
\end{exoresolu}

\courssub{Method 6 — Constructing a Geological Cross-Section}

\begin{methode}[Drawing a cross-section from a geological map]
\begin{enumerate}
\item \textbf{Choose the section line} $XY$ (usually given; if free, draw it perpendicular to the strike) and mark it on the map.
\item \textbf{Transfer the topography:} place a strip of paper along $XY$, tick every contour crossing with its altitude, then plot these heights on the profile frame using the vertical scale; join with a smooth curve.
\item \textbf{Transfer the geology:} on the same strip, tick every geological boundary crossing $XY$; project these points onto the topographic profile.
\item \textbf{Draw the boundaries underground:} from each boundary point, draw the contact at the \cle{true dip} (read from dip symbols or strike lines) in the correct direction. Boundaries of parallel strata must remain parallel.
\item \textbf{Complete the structures:} extend beds beneath unconformities (truncated below, with the younger series above), offset beds at faults (show the fault plane with its dip), and keep intrusions cross-cutting.
\item \textbf{Finish:} colour or ornament each unit, add a key, title, scales (state any vertical exaggeration), and label the structures (fault, fold axis, unconformity).
\end{enumerate}
\end{methode}

\begin{exoresolu}[Section across a dipping sequence cut by a fault]
On a map at $1:25\,000$, section line $XY$ runs north--south, perpendicular to an east--west strike. From $X$ (north) to $Y$ (south) it crosses: shale (boundary at $2.0\ \mathrm{cm}$ from $X$), sandstone (boundary at $4.0\ \mathrm{cm}$), then a fault at $6.0\ \mathrm{cm}$, then limestone to $Y$. All beds dip $30^{\circ}$ south; the fault is vertical. The ground is horizontal at $150\ \mathrm{m}$. Describe the section you would draw and state the ground positions of the boundaries.
\tcblower
\textbf{Step 1 — Ground positions (scale $1:25\,000$, so $1\ \mathrm{cm} = 250\ \mathrm{m}$):}
shale/sandstone boundary at $2.0\ \mathrm{cm} \rightarrow 500\ \mathrm{m}$ from $X$; sandstone/fault boundary at $4.0\ \mathrm{cm} \rightarrow 1000\ \mathrm{m}$; fault at $6.0\ \mathrm{cm} \rightarrow 1500\ \mathrm{m}$.

\textbf{Step 2 — Topography:} the ground is flat at $150\ \mathrm{m}$, so draw a straight horizontal profile line.

\textbf{Step 3 — Boundaries:} at $500\ \mathrm{m}$ and $1000\ \mathrm{m}$ draw the two contacts dipping $30^{\circ}$ towards the south (descending towards $Y$). Keep them parallel: shale overlies sandstone, and the sequence youngs southwards.

\textbf{Step 4 — Fault:} at $1500\ \mathrm{m}$ draw a \cle{vertical line} through the full depth of the section. South of the fault, draw the limestone. Because the map shows limestone immediately south of the fault, against sandstone to the north, the southern block has moved \emph{up} relative to the northern block (a deeper/younger-in-sequence unit is brought against the sandstone at the surface) — equivalently, the northern block moved down. Show the offset by drawing the contacts in the southern block at a higher level than their equivalents in the northern block.

\textbf{Step 5 — Finish:} ornament the units (dots for sandstone, brick pattern for limestone, dashes for shale), label "vertical fault", add the title "Cross-section along $XY$", the horizontal scale $1:25\,000$, the vertical scale as chosen (state any exaggeration), and a key.

\textbf{Examiner's reflex:} marks are awarded for correct dip direction, parallel boundaries, correct sense of fault offset, and a labelled key — not for artistic colouring.
\end{exoresolu}

\courssub{Method 7 — Describing Structures and Writing the Geological History}

\begin{methode}[Describing map structures and deducing geological history]
\begin{enumerate}
\item \textbf{Inventory first:} list the rock units from the key (oldest $\rightarrow$ youngest using superposition, cross-cutting and included-fragment criteria).
\item \textbf{Outcrop pattern:} V-shaped boundary deflections across valleys indicate the dip direction (the "rule of V's"); parallel bands = dipping strata; closed loops = folds, hills or basins.
\item \textbf{Folds:} identify the axial trace; oldest rocks in the core $=$ \cle{anticline}, youngest in the core $=$ \cle{syncline}; dips away from/towards the axis confirm. Plunging folds give nose-shaped (hooked) outcrop patterns closing in the plunge direction for anticlines.
\item \textbf{Faults:} look for truncated and offset boundaries and repeated or omitted strata; state the type (normal/reverse/strike-slip) from the relative movement and the dip of the fault plane.
\item \textbf{Unconformities:} the basal bed of the younger series truncates the older beds and structures.
\item \textbf{History:} narrate the events in order: deposition (oldest first) $\rightarrow$ tilting/folding $\rightarrow$ erosion (unconformity) $\rightarrow$ further deposition $\rightarrow$ faulting/intrusion (the youngest events cut everything else) $\rightarrow$ final erosion shaping the present relief.
\end{enumerate}
\end{methode}

\begin{exoresolu}[Geological history of a map area]
A map shows, from oldest to youngest in the key: mica schist (basement); then conglomerate, sandstone and limestone (dipping $25^{\circ}$ E, in parallel bands); a basalt dyke cutting the schist and all three sedimentary units; and finally a west--east vertical fault truncating all units, with the sandstone--limestone boundary shifted $1\ \mathrm{km}$ to the west on the northern side. Write the geological history of the area.
\tcblower
\textbf{Order of events (oldest first):}
\begin{enumerate}
\item \textbf{Formation of the mica schist} — metamorphism of pre-existing rocks, forming the basement.
\item \textbf{Erosion of the basement}, producing the surface on which the sediments rest (an unconformity between the schist and the conglomerate).
\item \textbf{Deposition of conglomerate, then sandstone, then limestone} in horizontal layers (principle of original horizontality and superposition), recording a transgressive (deepening) environment.
\item \textbf{Tilting} of the sedimentary sequence to its present $25^{\circ}$ easterly dip (a tectonic event), since the beds were originally horizontal.
\item \textbf{Intrusion of the basalt dyke:} it cuts the schist and all three sedimentary units, so by the principle of \cle{cross-cutting relationships} it is younger than all of them; its fine grain records rapid cooling of magma injected along a fracture.
\item \textbf{Faulting} along the west--east vertical fault: the westward shift of the boundaries on the northern block indicates \cle{strike-slip movement} — the northern block moved west relative to the southern block. The fault cuts and offsets every unit, including the dyke, so it is the youngest tectonic event.
\item \textbf{Final erosion} sculpting the present landscape and exposing the units at the surface.
\end{enumerate}
\textbf{Key principles used:} superposition (sediments young upward), original horizontality (tilting came after deposition), cross-cutting relationships (the dyke is younger than everything it cuts; the fault is younger than everything it offsets). In the GCE examination, present the history as a numbered chronological list and justify each step with one criterion — that is where the marks are.
\end{exoresolu}

\courssub{Method 8 — Interpreting Geological Photographs}

\begin{methode}[Analysing a geological photograph]
\begin{enumerate}
\item \textbf{Observe before interpreting:} describe only what is visible (layers, fractures, textures, landforms, vegetation contrasts) before naming structures.
\item \textbf{Bedding vs.\ jointing:} bedding planes are continuous, parallel and bound different lithologies; joints are fractures without displacement, often in regular sets.
\item \textbf{Dip estimation:} use the slope of bedding traces relative to the horizontal (skyline, water surface) and any scale object (person, house, tree) for size.
\item \textbf{Structures:} look for fold hinges (curved layers), fault scarps (offset layers, linear vegetation lines), dykes (dark or light cross-cutting walls), unconformities (angular discordance between two sets of layers).
\item \textbf{Landform--rock links:} resistant rocks (sandstone, lava) form ridges and cliffs; soft rocks (shale) form valleys and gentle slopes — e.g. the cliffed lava flows of Mount Cameroon versus the lowlands on weathered basement.
\item \textbf{Conclude} with the geological significance: environment of deposition, tectonic episode, or volcanic activity illustrated.
\end{enumerate}
\end{methode}

\begin{exoresolu}[Photograph of a sea cliff with tilted layers cut by a dark wall]
A photograph shows a sea cliff in which pale, layered rocks dip about $20^{\circ}$ to the right and are cut across by a near-vertical dark wall of fine-grained rock about $3\ \mathrm{m}$ wide (scaled against a person standing nearby). The dark rock cuts the layers at a high angle and shows no offset of the layers on either side. Interpret the photograph.
\tcblower
\textbf{Observation:} pale, parallel, continuous layers $=$ \cle{bedded sedimentary rocks} (e.g. sandstone); the beds dip $\approx 20^{\circ}$ to the right of the view. A dark, vertical, tabular body $\approx 3\ \mathrm{m}$ thick cuts across the bedding at a high angle; the beds on either side are \emph{not} displaced.

\textbf{Interpretation:}
\begin{itemize}
\item The layered rocks were deposited horizontally and later \textbf{tilted} by tectonic movements to their present $20^{\circ}$ dip.
\item The dark body is a \cle{dolerite or basalt dyke}: a near-vertical sheet of fine-grained (quickly cooled) igneous rock intruded along a fracture. Its fine grain indicates rapid cooling of magma near the surface.
\item The absence of offset across the dyke shows it is an \textbf{intrusion, not a fault}: the magma simply filled a tensional fracture.
\item \textbf{Relative ages (cross-cutting):} sediments $\rightarrow$ tilting $\rightarrow$ dyke intrusion (the dyke is the youngest feature, since it cuts the tilted beds).
\end{itemize}
\textbf{Cameroon link:} identical basaltic dykes cutting sedimentary or basement rocks occur along the \cle{Cameroon Volcanic Line}, e.g. around the Bambouto and Manengouba highlands. \textbf{Exam tip:} always state the scale evidence ("the person gives a thickness of about $3\ \mathrm{m}$") — examiners reward the use of scale and the distinction between observation and interpretation.
\end{exoresolu}

\newpage

\coursec{Exercises}

\courssub{I check my knowledge}

\begin{exercice}[MCQ: geological maps and their elements]
For each question, choose the single correct answer and justify your choice in one sentence.
\begin{enumerate}
\item A geological map on which formations are grouped by age rather than by rock type is called:
\begin{enumerate}
\item a lithological map;
\item a chronological (stratigraphic) map;
\item a structural map;
\item a topographic map.
\end{enumerate}
\item On a geological map, a circle with a central dot accompanied by a number (e.g. 35) at an outcrop indicates:
\begin{enumerate}
\item the thickness of the bed in metres;
\item the strike and dip of the strata;
\item the trend of a fold axis;
\item a borehole with its reference number.
\end{enumerate}
\item A map at the scale 1:50\,000 is, compared with a map at 1:10\,000:
\begin{enumerate}
\item a larger-scale map showing more detail;
\item a smaller-scale map showing less detail;
\item a map of identical scale;
\item a map whose scale cannot be compared.
\end{enumerate}
\item The angle between the horizontal and an inclined bed, measured in a vertical plane perpendicular to the strike, is:
\begin{enumerate}
\item the apparent dip;
\item the true dip;
\item the strike;
\item the plunge.
\end{enumerate}
\end{enumerate}
\end{exercice}

\begin{exercice}[True or false — justify each answer]
State whether each assertion is true or false, and justify your answer in one or two sentences.
\begin{enumerate}
\item The outcrop of a vertical bed crosses contour lines without being deflected by them.
\item On flat ground, the width of outcrop of a bed of given true thickness increases as the dip decreases.
\item In the three-point problem, the three points used to determine the attitude of a bed must lie on the same straight line.
\item Repetition of beds on a map can be produced either by folding or by faulting.
\item On a geological cross-section drawn without vertical exaggeration, the horizontal and vertical scales must be equal.
\item A dyke always appears on a map as a band whose outcrop is deflected by valleys in the same way as a horizontal bed.
\end{enumerate}
\end{exercice}

\begin{exercice}[Definitions]
Define precisely each of the following terms, using correct geological vocabulary:
\begin{enumerate}
\item geological map;
\item strike of a bed;
\item true dip and apparent dip;
\item outcrop;
\item unconformity;
\item scale of a map;
\item grid reference (six-figure);
\item geological cross-section.
\end{enumerate}
\end{exercice}

\begin{exercice}[Direct calculations: scale, distance and bearing]
On a topographic map at the scale 1:25\,000, the villages of Bafut and Wum are $8.4\ \mathrm{cm}$ apart on the map.
\begin{enumerate}
\item Calculate the real horizontal distance between the two villages, in metres and in kilometres.
\item Express the scale of the map as a verbal statement ("1 cm represents \ldots").
\item The line joining Bafut to Wum makes an angle of $52^{\circ}$ clockwise from grid north. State the bearing of Wum from Bafut, and the back-bearing of Bafut from Wum.
\item On the same map, a lava flow of Mount Cameroon covers an area of $3.6\ \mathrm{cm^{2}}$. Calculate the real area covered, in $\mathrm{km^{2}}$.
\end{enumerate}
\end{exercice}

\courssub{I practise}

\begin{exercice}[Reading conventional symbols]
On a given geological map extract (provided by your teacher or taken from a past GCE paper), ten conventional symbols are labelled (a) to (j): dip-and-strike symbol, vertical-bed symbol, horizontal-bed symbol, anticlinal axis, synclinal axis, normal fault, borehole, spring, dyke, and fossil locality.
\begin{enumerate}
\item Identify and name each of the ten symbols.
\item For each symbol, state precisely what information it gives about the geology of the area.
\item Redraw the dip-and-strike symbol for a bed dipping $40^{\circ}$ towards the south-east, and label its parts (strike line, dip arrow, dip angle).
\end{enumerate}
\end{exercice}

\begin{exercice}[Thickness of a dipping bed]
A sandstone bed outcrops over a width of $200\ \mathrm{m}$ measured on flat, horizontal ground, perpendicular to the strike. The bed dips at $30^{\circ}$.
\begin{enumerate}
\item Draw a labelled sketch cross-section showing the outcrop width, the dip, the true thickness and the vertical thickness.
\item Calculate the \cle{true thickness} of the bed.
\item Calculate the \cle{vertical thickness} of the bed (the thickness that would be traversed by a vertical borehole).
\item Explain why the vertical thickness is always greater than the true thickness for a dipping bed.
\end{enumerate}
\end{exercice}

\begin{exercice}[Solving a three-point problem]
Three boreholes A, B and C are drilled on flat ground at an elevation of $500\ \mathrm{m}$. They intersect the top of the same sandstone bed at the following absolute elevations: A: $420\ \mathrm{m}$; B: $360\ \mathrm{m}$; C: $300\ \mathrm{m}$. On the map, B lies $4\ \mathrm{cm}$ due east of A, and C lies $5\ \mathrm{cm}$ due south of A.
\begin{enumerate}
\item Plot the three boreholes on a suitable base map at the scale 1:10\,000.
\item Solve the three-point problem graphically by constructing strike lines (structure contours) for the top of the sandstone bed.
\item Deduce the strike, the dip direction and the amount of dip of the bed.
\item State two practical applications of this construction in exploration geology.
\end{enumerate}
\end{exercice}

\begin{exercice}[Folding versus faulting: repetition of beds]
A map shows the following sequence of outcrops from west to east: sandstone -- shale -- limestone -- shale -- sandstone.
\begin{enumerate}
\item Show, with the aid of two labelled map sketches and two short cross-sections, how such a repetition of beds can be produced (i) by folding and (ii) by faulting.
\item State the map evidence (dip symbols, ages of beds, disposition of outcrops) that allows you to distinguish between the two interpretations.
\item For each case, state whether the repeated beds appear in the same order or in mirror image, and explain why.
\end{enumerate}
\end{exercice}

\courssub{I prepare for the GCE}

\begin{exercice}[GCE-style compound map question (20 marks)]
You are given a geological map extract at the scale 1:50\,000 showing sedimentary formations from Cretaceous sandstone to Recent alluvium, folded into a plunging anticline, cut by a normal fault trending north--south, and intruded by a basaltic dyke.
\begin{enumerate}
\item State the scale of the map in two different ways, and give the six-figure grid reference of the borehole marked on the extract. \textbf{(4 marks)}
\item Describe fully the geological structures shown on the map, using the dip symbols and outcrop patterns as evidence. \textbf{(6 marks)}
\item Construct a geological cross-section along the line X--Y drawn across the map, showing the fold, the fault and the dyke. \textbf{(6 marks)}
\item Write the geological history of the area as a numbered sequence of events, from the oldest to the youngest, justifying the relative age of each event. \textbf{(4 marks)}
\end{enumerate}
\end{exercice}

\begin{exercice}[Cross-section across a plunging anticline cut by a fault (15 marks)]
A map extract shows a plunging anticline with a north-easterly plunge, whose core is occupied by granite and whose flanks are formed of schist, quartzite and limestone in outward succession. The fold is displaced by a normal fault striking east--west, and dip symbols are provided on both flanks.
\begin{enumerate}
\item Using the dip symbols and the outcrop pattern, determine the direction and amount of dip on each flank of the fold, and the direction of plunge of the axis. \textbf{(4 marks)}
\item Construct a geological cross-section along the line X--Y across the fold and the fault, respecting the scale of the map and without vertical exaggeration. \textbf{(7 marks)}
\item From the displacement of the outcrops across the fault, determine the downthrown side and explain your reasoning. \textbf{(2 marks)}
\item State two criteria visible on the map that confirm that the fold plunges, and not that it is horizontal. \textbf{(2 marks)}
\end{enumerate}
\end{exercice}

\begin{exercice}[Interpretation of geological photographs (15 marks)]
Two photographs are provided for study. Photograph 1 is an \emph{oblique} aerial photograph of a cuesta landscape in the Benue trough region, showing an asymmetric ridge with one gentle slope and one steep slope, and a stream network. Photograph 2 is a \emph{vertical} aerial photograph showing a radial drainage pattern around a circular feature on the flank of Mount Cameroon.
\begin{enumerate}
\item On photograph 1, identify the dip slope and the scarp slope of the cuesta, and deduce the direction of dip of the underlying strata. \textbf{(4 marks)}
\item Name and describe the drainage pattern developed on photograph 1, and explain its relationship with the alternating hard and soft beds. \textbf{(3 marks)}
\item On photograph 2, name the drainage pattern and interpret the circular landform in terms of the underlying geology. \textbf{(4 marks)}
\item State two advantages and two limitations of oblique photographs compared with vertical aerial photographs for geological interpretation. \textbf{(4 marks)}
\end{enumerate}
\end{exercice}

\newpage

\coursec{Solutions to the exercises}

\begin{corrige}[Exercise 1 — MCQ: geological maps and their elements]
\begin{enumerate}
\item \textbf{Answer: (b) a chronological (stratigraphic) map.} On a \cle{chronological map} the units shown are time-stratigraphic units (formations grouped by age, e.g. Cretaceous, Tertiary), whereas a \cle{lithological map} distinguishes units by rock type regardless of age.
\item \textbf{Answer: (b) the strike and dip of the strata.} The \cle{dip-and-strike symbol} (a T-shaped or arrow symbol with a number) records the direction of the strike line and the amount and direction of dip measured at that outcrop; the number (35) is the dip angle in degrees.
\item \textbf{Answer: (b) a smaller-scale map showing less detail.} The representative fraction $1/50\,000$ is smaller than $1/10\,000$, so 1 cm represents a greater ground distance ($500\ \mathrm{m}$ instead of $100\ \mathrm{m}$) and less detail can be shown.
\item \textbf{Answer: (b) the true dip.} The \cle{true dip} is the maximum angle of inclination of a bed, measured from the horizontal in the vertical plane perpendicular to the strike; any other vertical section gives a smaller, \cle{apparent dip}.
\end{enumerate}
\end{corrige}

\begin{corrige}[Exercise 2 — True or false]
\begin{enumerate}
\item \textbf{True.} A \cle{vertical bed} has a horizontal trace that is a straight line independent of relief: its outcrop is not deflected by contour lines (the ``rule of V's'' gives no V).
\item \textbf{True.} On flat ground the outcrop width is $w = t/\sin\delta$ (where $t$ is the \cle{true thickness} and $\delta$ the dip); as $\delta$ decreases, $\sin\delta$ decreases and $w$ increases, until a horizontal bed outcrops over the whole area.
\item \textbf{False.} The three points must \emph{not} be collinear: three collinear points define only a line, not a plane, so the attitude of the bed could not be determined. They must form a triangle.
\item \textbf{True.} \cle{Folding} repeats beds in mirror-image order on either side of an axis; \cle{faulting} can juxtapose and repeat beds in the same order across the fault plane. Both produce repetition on the map.
\item \textbf{True.} Absence of \cle{vertical exaggeration} means by definition that the vertical scale equals the horizontal scale; only then are dips represented at their true angular values.
\item \textbf{False.} A \cle{dyke} is a steeply dipping or vertical intrusion; its outcrop crosses valleys in a straight or nearly straight line, whereas a horizontal bed follows the contour lines and is strongly deflected by valleys.
\end{enumerate}
\end{corrige}

\begin{corrige}[Exercise 3 — Definitions]
\begin{enumerate}
\item \textbf{\cle{Geological map}:} a map, drawn on a topographic base at a given scale, showing the distribution, boundaries, nature, age and attitude of rock formations at or near the Earth's surface, using colours, ornament and conventional symbols.
\item \textbf{\cle{Strike of a bed}:} the direction (bearing) of the horizontal line formed by the intersection of the bedding plane with a horizontal plane; it is perpendicular to the direction of true dip.
\item \textbf{\cle{True dip} and \cle{apparent dip}:} the true dip is the maximum angle between the bedding plane and the horizontal, measured in the vertical plane perpendicular to the strike; the apparent dip is the (smaller) angle of inclination measured in any other vertical section, with $\tan\delta' = \tan\delta\,\sin\theta$, where $\theta$ is the angle between the section and the true dip direction.
\item \textbf{\cle{Outcrop}:} the area over which a rock formation is exposed at the Earth's surface (or would be exposed if soil and superficial deposits were removed); also the line on a map representing the intersection of a geological boundary with the topography.
\item \textbf{\cle{Unconformity}:} a surface of erosion or non-deposition separating two series of strata whose deposition was interrupted by tectonic deformation and/or erosion, so that the younger series rests on older, often differently oriented, rocks; it marks a gap (hiatus) in the geological record.
\item \textbf{\cle{Scale of a map}:} the constant ratio between a distance measured on the map and the corresponding horizontal distance on the ground, expressed as a representative fraction (e.g. $1:25\,000$), verbally, or graphically.
\item \textbf{\cle{Grid reference (six-figure)}:} the coordinates of a point read on the kilometre grid of a topographic map: first the easting (two figures plus the estimated tenth of a kilometre, read along the bottom), then the northing (three figures, read up the side), locating the point within a $100\ \mathrm{m}$ square.
\item \textbf{\cle{Geological cross-section}:} a diagram representing the arrangement of rock formations along a vertical plane cutting the Earth's surface, constructed from the outcrop pattern, dip data and topography of a geological map, at stated horizontal and vertical scales.
\end{enumerate}
\end{corrige}

\begin{corrige}[Exercise 4 — Direct calculations: scale, distance and bearing]
\begin{enumerate}
\item At 1:25\,000, 1 cm on the map represents $25\,000\ \mathrm{cm} = 250\ \mathrm{m}$ on the ground. Hence
\[ d = 8.4\ \mathrm{cm} \times 250\ \mathrm{m\,cm^{-1}} = 2100\ \mathrm{m} = 2.1\ \mathrm{km}. \]
\item Verbal statement: ``1 cm on the map represents 250 m on the ground'' (equivalently ``4 cm represent 1 km'').
\item The bearing of Wum from Bafut is the clockwise angle from grid north: $\mathrm{N}\,052^{\circ}$, i.e. a bearing of $052^{\circ}$. The back-bearing differs by $180^{\circ}$:
\[ 052^{\circ} + 180^{\circ} = 232^{\circ}, \]
so the bearing of Bafut from Wum is $232^{\circ}$.
\item Areas scale as the square of the scale factor. Since 1 cm represents $0.25\ \mathrm{km}$, $1\ \mathrm{cm^{2}}$ represents $(0.25)^{2} = 0.0625\ \mathrm{km^{2}}$. Therefore
\[ A = 3.6\ \mathrm{cm^{2}} \times 0.0625\ \mathrm{km^{2}\,cm^{-2}} = 0.225\ \mathrm{km^{2}}. \]
\end{enumerate}
\end{corrige}

\begin{corrige}[Exercise 5 — Reading conventional symbols]
\begin{enumerate}
\item Identification of the ten symbols:
\begin{center}
\begin{tabularx}{\linewidth}{|l|X|}
\hline
\rowcolor{popblueL} Label & Symbol \\
\hline
(a) & Dip-and-strike symbol (T with arrow and dip figure) \\
\hline
(b) & Vertical-bed symbol (crossed T / strike line with ticks on both sides) \\
\hline
(c) & Horizontal-bed symbol (circle with central dot, or ``+'' convention) \\
\hline
(d) & Anticlinal axis (line with arrows pointing away from the axis) \\
\hline
(e) & Synclinal axis (line with arrows pointing towards the axis) \\
\hline
(f) & Normal fault (thick line, tick or ball on the downthrown side) \\
\hline
(g) & Borehole (small circle, often numbered, with collar elevation) \\
\hline
(h) & Spring (small blue dot with a short tail) \\
\hline
(i) & Dyke (narrow band or thick line of igneous ornament crossing other units) \\
\hline
(j) & Fossil locality (crossed-hammer or star symbol, or labelled point) \\
\hline
\end{tabularx}
\end{center}
\item Information given by each symbol:
\begin{itemize}
\item \textbf{Dip-and-strike:} direction of strike, direction and amount of true dip at that outcrop.
\item \textbf{Vertical bed:} strata dipping at $90^{\circ}$; outcrop width equals true thickness.
\item \textbf{Horizontal bed:} strata with $0^{\circ}$ dip; boundaries follow contour lines.
\item \textbf{Anticlinal axis:} trace of the axial plane of an upfold; beds dip away from the axis, oldest rocks in the core.
\item \textbf{Synclinal axis:} trace of the axial plane of a downfold; beds dip towards the axis, youngest rocks in the core.
\item \textbf{Normal fault:} a fracture with dip-slip displacement; the tick marks the downthrown block.
\item \textbf{Borehole:} point of subsurface data; gives depth/elevation of horizons encountered.
\item \textbf{Spring:} emergence of groundwater, often along a fault or an impermeable boundary.
\item \textbf{Dyke:} a discordant tabular intrusion, younger than the rocks it cuts.
\item \textbf{Fossil locality:} site where fossils were collected, used for dating and correlating the beds.
\end{itemize}
\item Dip-and-strike symbol for a bed dipping $40^{\circ}$ to the south-east:
\begin{popfigure}
\begin{center}
\begin{tikzpicture}[scale=1.0]
\draw[popink, thick] (-2,0) -- (2,0);
\draw[poporange, thick, -latex] (0,0) -- (0.85,-0.85);
\node[poporange, anchor=west] at (0.95,-0.85) {$40^{\circ}$};
\node[popink, anchor=south] at (-1.4,0.05) {strike line (NW--SE)};
\node[poporange, anchor=south west] at (0.15,-0.35) {dip arrow};
\draw[popmuted, dashed, -latex] (0,1.0) -- (0,0.15);
\node[popmuted, anchor=south] at (0,1.0) {N};
\end{tikzpicture}
\end{center}
\legende{Dip-and-strike symbol: strike NW--SE, dip $40^{\circ}$ towards the SE.}
\end{popfigure}
The long line is the strike line (perpendicular to the dip direction), the arrow points in the dip direction (SE), and the figure 40 is the dip angle in degrees.
\end{enumerate}
\end{corrige}

\begin{corrige}[Exercise 6 — Thickness of a dipping bed]
\begin{enumerate}
\item Labelled sketch cross-section:
\begin{popfigure}
\begin{center}
\begin{tikzpicture}[scale=0.9]
\draw[popink] (0,0) -- (8,0);
\draw[popblue, thick] (1,0) -- (7,-3.46);
\draw[popblue, thick] (1.8,0) -- (7.8,-3.46);
\draw[popmuted, dashed] (1,0) -- (1,-2.6);
\draw[popmuted, dashed] (1.8,0) -- (1.8,-3.2);
\draw[popgreen, thick, latex-latex] (2.6,-0.92) -- (3.4,-0.92);
\node[popgreen, anchor=south] at (3.0,-0.9) {$t$};
\draw[popgreen, thick, latex-latex] (1.55,-0.1) -- (1.55,-1.5);
\node[popgreen, anchor=west] at (1.65,-0.8) {$t_v$};
\draw[popdark, latex-latex] (1,0.25) -- (1.8,0.25);
\node[popdark, anchor=south] at (1.4,0.3) {$w$};
\draw[popmuted] (3.2,0) arc (0:-30:2.2);
\node[popmuted] at (3.55,-0.35) {$\delta$};
\node[popink, anchor=west] at (8.1,0) {surface};
\end{tikzpicture}
\end{center}
\legende{Cross-section perpendicular to strike: outcrop width $w$, dip $\delta$, true thickness $t$ (perpendicular to bedding) and vertical thickness $t_v$.}
\end{popfigure}
\item The outcrop width, dip and true thickness are related by
\[ t = w\,\sin\delta = 200 \times \sin 30^{\circ} = 200 \times 0.5 = 100\ \mathrm{m}. \]
\item The vertical thickness is
\[ t_v = w\,\tan\delta = 200 \times \tan 30^{\circ} = 200 \times 0.5774 \approx 115.5\ \mathrm{m}. \]
\item The true thickness is measured \emph{perpendicular} to the bedding planes, which is the shortest distance between them. A vertical borehole crosses the bed obliquely (the bed is inclined), so the path through the bed is longer: geometrically $t = t_v \cos\delta$, and since $\cos\delta < 1$ for any non-vertical bed, $t_v > t$ always.
\end{enumerate}
\end{corrige}

\begin{corrige}[Exercise 7 — Solving a three-point problem]
\begin{enumerate}
\item Base map: at 1:10\,000, 1 cm represents $100\ \mathrm{m}$. Plot A at the origin, B $4\ \mathrm{cm}$ due east of A, C $5\ \mathrm{cm}$ due south of A. Annotate each point with the elevation of the bed top: A $= 420\ \mathrm{m}$, B $= 360\ \mathrm{m}$, C $= 300\ \mathrm{m}$.
\item Construction of strike lines (structure contours):
\begin{itemize}
\item The elevation difference between A and B is $420 - 360 = 60\ \mathrm{m}$; between B and C it is $360 - 300 = 60\ \mathrm{m}$; between A and C it is $120\ \mathrm{m}$.
\item Along AC, the $360\ \mathrm{m}$ level lies where the elevation has dropped $60\ \mathrm{m}$ out of $120\ \mathrm{m}$, i.e. halfway: point D, $2.5\ \mathrm{cm}$ south of A.
\item Join B to D: BD is the $360\ \mathrm{m}$ strike line (structure contour).
\item Through A and C draw parallels to BD: these are the $420\ \mathrm{m}$ and $300\ \mathrm{m}$ strike lines respectively. The strike lines are equally spaced, confirming a plane bed.
\end{itemize}
\begin{popfigure}
\begin{center}
\begin{tikzpicture}[scale=0.85]
\draw[popmuted, -latex] (0,0) -- (0,1.4) node[anchor=south]{N};
\fill[popink] (0,0) circle (0.07) node[anchor=south west]{A (420 m)};
\fill[popink] (4,0) circle (0.07) node[anchor=south west]{B (360 m)};
\fill[popink] (0,-5) circle (0.07) node[anchor=north west]{C (300 m)};
\fill[poporange] (0,-2.5) circle (0.07) node[anchor=west]{D (360 m)};
\draw[popblue, thick] (-1.3,-1.63) -- (5.3,-2.7);
\draw[popblue, thick] (-1.3,0.87) -- (5.3,-1.2);
\draw[popblue, thick] (-1.3,-4.13) -- (5.3,-6.2);
\node[popblue, anchor=west] at (5.3,-1.2) {420 m};
\node[popblue, anchor=west] at (5.3,-2.7) {360 m};
\node[popblue, anchor=west] at (5.3,-6.2) {300 m};
\draw[poporange, thick, -latex] (1.6,-1.05) -- (2.15,-2.4);
\node[poporange, anchor=west] at (2.2,-1.8) {dip dir.};
\draw[popmuted, dashed] (0,0) -- (4,0) -- (0,-5) -- cycle;
\end{tikzpicture}
\end{center}
\legende{Three-point problem: D on AC at the 360 m level; BD is the 360 m strike line; parallels give the 420 m and 300 m strike lines.}
\end{popfigure}
\item Results:
\begin{itemize}
\item \textbf{Strike:} direction of BD, approximately N$\,120^{\circ}$E / N$\,300^{\circ}$E (i.e. ESE--WNW, about $30^{\circ}$ south of east), read with a protractor on the construction.
\item \textbf{Dip direction:} perpendicular to the strike lines, towards decreasing elevations, i.e. towards the south-south-west (about N$\,210^{\circ}$).
\item \textbf{Amount of dip:} the spacing between the 420 m and 360 m strike lines measured on the map is about $2.3\ \mathrm{cm} \approx 230\ \mathrm{m}$ on the ground for a vertical interval of $60\ \mathrm{m}$:
\[ \tan\delta = \frac{60}{230} \approx 0.26 \quad\Rightarrow\quad \delta \approx 15^{\circ}. \]
\end{itemize}
\item Applications:
\begin{itemize}
\item predicting the depth at which a borehole will intersect a mineral seam, aquifer or oil-bearing bed before drilling;
\item constructing structure-contour maps of a reservoir or coal seam to plan mining works, tunnels and water-supply boreholes.
\end{itemize}
\end{enumerate}
\end{corrige}

\begin{corrige}[Exercise 8 — Folding versus faulting: repetition of beds]
\begin{enumerate}
\item \textbf{(i) Folding:} an upright syncline whose axis trends north--south produces, from west to east, the mirror-image sequence sandstone -- shale -- limestone -- shale -- sandstone, with the youngest bed (limestone) in the core and dips converging towards the axis.
\begin{popfigure}
\begin{center}
\begin{tikzpicture}[scale=0.8]
\draw[popink, thick] (0,0) .. controls (2,-1.6) and (3,-1.6) .. (5,0);
\draw[popink, thick] (0,0.6) .. controls (2,-1.0) and (3,-1.0) .. (5,0.6);
\draw[popink, thick] (0,1.2) .. controls (2,-0.4) and (3,-0.4) .. (5,1.2);
\draw[popink, thick] (0,1.8) -- (5,1.8);
\node at (2.5,2.1) {sandstone};
\node at (2.5,0.75) {shale};
\node at (2.5,-0.55) {limestone};
\draw[popmuted, dashed] (2.5,1.8) -- (2.5,-1.8);
\node[popmuted, anchor=north] at (2.5,-1.85) {axis};
\draw[poporange, -latex] (1.2,1.0) -- (1.7,0.55);
\draw[poporange, -latex] (3.8,1.0) -- (3.3,0.55);
\node[poporange, anchor=south] at (1.1,1.0) {dip};
\node[poporange, anchor=south] at (4.0,1.0) {dip};
\end{tikzpicture}
\end{center}
\legende{Cross-section of a syncline: mirror-image repetition, dips converge towards the axis.}
\end{popfigure}
\textbf{(ii) Faulting:} a faulted homoclinal sequence (e.g. a reverse fault, or a normal fault repeating a dipping series) can juxtapose the same succession twice in the \emph{same} order: sandstone -- shale -- limestone, then again shale -- sandstone across the fault trace, depending on the geometry of displacement and erosion.
\begin{popfigure}
\begin{center}
\begin{tikzpicture}[scale=0.8]
\draw[popink, thick] (0,0) -- (2.2,-1.4);
\draw[popink, thick] (0,0.7) -- (2.2,-0.7);
\draw[popink, thick] (0,1.4) -- (2.2,0);
\draw[popink, thick] (2.8,0.4) -- (5,-1.0);
\draw[popink, thick] (2.8,1.1) -- (5,-0.3);
\draw[popink, thick] (2.8,1.8) -- (5,0.4);
\draw[popdark, very thick] (2.2,1.9) -- (2.8,-1.6);
\node[popdark, anchor=south] at (2.5,1.9) {fault};
\node at (1.0,1.05) {sh};
\node at (1.0,0.35) {ls};
\node at (1.0,-0.5) {ss};
\node at (4.0,1.45) {sh};
\node at (4.0,0.7) {ls};
\node at (4.0,-0.1) {ss};
\end{tikzpicture}
\end{center}
\legende{Cross-section of a faulted dipping series: the same succession reappears across the fault (ss = sandstone, sh = shale, ls = limestone).}
\end{popfigure}
\item Distinguishing evidence on the map:
\begin{itemize}
\item \textbf{Dip symbols:} in a fold, dips are opposed on the two limbs and converge towards (syncline) or diverge from (anticline) the axis; in faulting, dips are generally in the same direction on both sides.
\item \textbf{Ages of beds:} a fold core contains the youngest (syncline) or oldest (anticline) bed, with ages symmetric about the axis; a fault cuts and displaces boundaries abruptly without age symmetry.
\item \textbf{Disposition of outcrops:} fold closures are curved and continuous around the nose of the fold; a fault trace is a straight or gently curved line that truncates and offsets boundaries, often marked by a tick on the downthrown side.
\end{itemize}
\item Order of repetition:
\begin{itemize}
\item \textbf{Folding:} the beds reappear in \emph{mirror image} (reverse order), because the limbs are the two flanks of the same folded succession read in opposite directions.
\item \textbf{Faulting:} the beds reappear in the \emph{same order}, because the fault displaces blocks of the same unfurled succession, which is then re-exposed by erosion in its original stratigraphic order.
\end{itemize}
\end{enumerate}
\end{corrige}

\begin{corrige}[Exercise 9 — GCE-style compound map question (20 marks)]
\textbf{Indicative mark scheme and model answers.}
\begin{enumerate}
\item \textbf{(4 marks)} Scale: stated as a representative fraction, $1:50\,000$ (1 mark), and verbally, ``1 cm on the map represents $500\ \mathrm{m}$ on the ground'' or ``2 cm represent 1 km'' (1 mark). Six-figure grid reference of the borehole: correct easting read first (3 figures) then northing (3 figures), e.g. 345 678 — 1 mark for the correct easting, 1 mark for the correct northing (accept the value read on the actual extract; the examiner expects the candidate to read the grid lines, estimate tenths, and quote eastings before northings).
\item \textbf{(6 marks)} Description of structures — award 1 mark per valid, evidence-based statement, e.g.:
\begin{itemize}
\item the sedimentary series (Cretaceous sandstone to Tertiary/Recent units) is folded; dip symbols on the two flanks point in opposite directions, away from a central axis (evidence of an anticline);
\item the oldest formation (sandstone) occupies the core, with progressively younger beds outwards, confirming an anticline;
\item the outcrops close around a nose and the dips decrease along the axis in one direction, showing that the anticline \emph{plunges};
\item a north--south fault (tick on the downthrown side) offsets the boundaries: a normal fault;
\item a narrow, straight basaltic dyke cuts across all the sedimentary boundaries without being deflected by relief, indicating a steeply dipping discordant intrusion;
\item Recent alluvium occupies the valley floors, unconformably overlying the older formations.
\end{itemize}
Examiner expectation: every structural statement must be supported by map evidence (dip symbols, ages, outcrop pattern), not merely asserted.
\item \textbf{(6 marks)} Cross-section along X--Y:
\begin{itemize}
\item correct topographic profile transferred from contour intersections (1 mark);
\item equal horizontal and vertical scales, stated on the section (1 mark);
\item boundaries projected from the map and beds drawn with the dips given by the symbols, symmetric about the fold axis (2 marks);
\item fault drawn as a steep plane with correct downthrown side and offset of boundaries (1 mark);
\item dyke drawn as a near-vertical band cutting all units; formations labelled/ornamented with a key (1 mark).
\end{itemize}
\item \textbf{(4 marks)} Geological history (oldest to youngest), 1 mark per correctly ordered and justified event:
\begin{enumerate}
\item deposition of the sedimentary series (sandstone, then younger beds up to the youngest sedimentary formation) — superposition;
\item compressional folding producing the plunging anticline — the fold affects all the sedimentary beds, so it post-dates them;
\item normal faulting displacing the folded beds — the fault cuts the fold, so it is younger;
\item intrusion of the basaltic dyke cutting sediments, fold and fault, followed by erosion and deposition of Recent alluvium in the valleys — cross-cutting relationships and superposition.
\end{enumerate}
\end{enumerate}
\end{corrige}

\begin{corrige}[Exercise 10 — Cross-section across a plunging anticline cut by a fault (15 marks)]
\textbf{Indicative mark scheme and model answers.}
\begin{enumerate}
\item \textbf{(4 marks)} Reading the dip symbols: on the north-western flank the beds dip towards the NW, on the south-eastern flank towards the SE (dips directed \emph{away} from the axis — characteristic of an anticline) — 1 mark per flank; the amounts are read from the figures beside the symbols (e.g. $30^{\circ}$ NW and $40^{\circ}$ SE, values read on the extract — 1 mark). The outcrops of the granite core close towards the south-west and the fold opens towards the north-east, so the axis \emph{plunges towards the north-east} (1 mark).
\item \textbf{(7 marks)} Construction of the cross-section:
\begin{itemize}
\item transfer the relief along X--Y from the contour lines to a profile at the map scale, vertical scale equal to horizontal scale (no exaggeration) (2 marks);
\item mark the intersection of each boundary (granite/schist, schist/quartzite, quartzite/limestone) and of the fault with X--Y (1 mark);
\item draw the beds on each flank with the measured dips, converging downwards beneath the axis so that granite occupies the core of an upfold (2 marks);
\item draw the east--west fault as a steep plane offsetting the boundaries, with the correct downthrown side (1 mark);
\item label formations, add title, scales and key (1 mark).
\end{itemize}
\item \textbf{(2 marks)} Downthrown side: on the downthrown block, \emph{younger} beds are preserved against older ones, and a given boundary is displaced towards the stratigraphically lower (older) side. Comparing the outcrop positions of the same boundary on either side of the fault, the side where the boundaries are shifted in the direction of dip of the beds (or where younger formations abut the fault) is the downthrown side — identified here from the offset of the schist/quartzite boundary (1 mark for the correct side, 1 mark for the reasoning).
\item \textbf{(2 marks)} Criteria proving plunge (any two, 1 mark each):
\begin{itemize}
\item the outcrops form a \emph{nose} (V-shaped or U-shaped closure) around the end of the fold instead of parallel bands;
\item the granite core \emph{widens} in the direction opposite to the plunge (and narrows/closes in the plunge direction);
\item dip symbols show a component of dip along the axis (oblique dips on the flanks near the nose), and the strike lines converge towards the closure.
\end{itemize}
\end{enumerate}
\end{corrige}

\begin{corrige}[Exercise 11 — Interpretation of geological photographs (15 marks)]
\textbf{Indicative mark scheme and model answers.}
\begin{enumerate}
\item \textbf{(4 marks)} On photograph 1, the \emph{dip slope} is the long, gentle slope of the cuesta, which follows the inclination of the resistant bed (1 mark); the \emph{scarp slope} is the short, steep slope cutting across the bedding (1 mark). Since the dip slope coincides with the bedding surface of the hard bed, the strata dip in the direction of the gentle slope (1 mark), at an angle comparable to the gradient of that slope (1 mark).
\item \textbf{(3 marks)} The drainage is a \emph{trellis} (rectangular) pattern (1 mark): long subsequent streams occupy the strike valleys eroded in the soft beds (e.g. shale), short consequent streams run down the dip slope and scarp, and right-angled junctions dominate (1 mark). The pattern is controlled by the alternation of hard beds (forming the cuesta ridges) and soft beds (forming the valleys), so streams are guided along the strike of the weak strata (1 mark).
\item \textbf{(4 marks)} Photograph 2 shows a \emph{radial} drainage pattern (1 mark): streams flow outwards in all directions from a central high point (1 mark). The circular feature is interpreted as a volcanic cone (or dome) on the flank of Mount Cameroon: the resistant volcanic rocks build a cone, and the consequent streams radiate down its slopes (1 mark); the circular plan and radial drainage together indicate a localised centre of uplift/eruption, i.e. an igneous body such as a volcanic neck, cone or small intrusive dome (1 mark).
\item \textbf{(4 marks)} Oblique photographs — any two advantages and two limitations, 1 mark each:
\begin{itemize}
\item \textbf{Advantages:} they show relief, slopes and landforms in a natural perspective close to ground view, making scarps, cuestas and dips easy to recognise; they cover a large area to the horizon in a single view and are cheap to obtain.
\item \textbf{Limitations:} scale varies from foreground to background, so distances and areas cannot be measured accurately; they cannot be used for stereoscopic measurement or precise mapping (no reliable heights), and distant features are foreshortened or hidden by relief.
\end{itemize}
\end{enumerate}
\end{corrige}

\newpage

\coursec{Summary sheet}

\begin{popfigure}
\begin{tikzpicture}[
  grow cyclic,
  mindmap,
  concept color=popblueL,
  every node/.style={concept, circular drop shadow},
  root concept/.append style={concept color=popblue, font=\bfseries\small, text=white, minimum size=2.6cm},
  level 1 concept/.append style={level distance=3.6cm, sibling angle=51.4, font=\scriptsize, minimum size=1.9cm},
  level 2 concept/.append style={level distance=2.2cm, font=\tiny, minimum size=1.3cm}
]
\node[root concept]{Geological Map Work}
  child[concept color=poporange]{ node {Maps: types \& features}
    child { node {topographic vs geological} }
    child { node {scale, legend, north} }
  }
  child[concept color=popgreen]{ node {Orientation \& scale}
    child { node {bearings, grid north} }
    child { node {RF, graphic scale} }
  }
  child[concept color=popgold]{ node {Symbols \& strata geometry}
    child { node {strike, dip symbols} }
    child { node {true vs apparent dip} }
  }
  child[concept color=popink]{ node {Strike lines \& 3-point problem}
    child { node {join equal heights} }
    child { node {$\tan\delta = h/d$} }
  }
  child[concept color=poppurple]{ node {Cross-sections}
    child { node {V.E.\ choice} }
    child { node {folds, faults, intrusions} }
  }
  child[concept color=popblue]{ node {Structures \& history}
    child { node {outcrop patterns} }
    child { node {relative dating} }
  }
  child[concept color=popgreen]{ node {Photographs \& exam skill}
    child { node {landscape reading} }
    child { node {full interpretation} }
  };
\end{tikzpicture}
\legende{Mind map of the chapter \emph{Map Work}: from reading a geological map to constructing cross-sections and writing a complete geological history.}
\end{popfigure}

\begin{bilan}
\textbf{Key definitions}
\begin{itemize}
\item \cle{Geological map}: a map showing the distribution, nature and age of rock units at (or just below) the surface, together with their structural attitude (strike and dip).
\item \cle{Strike}: the direction of the horizontal line contained in an inclined plane; it is expressed as a bearing and has no unique sense (e.g.\ N$040^{\circ}$E is the same line as S$040^{\circ}$W).
\item \cle{Dip}: the maximum angle of inclination of a plane below the horizontal, measured in a vertical plane \emph{perpendicular to the strike}; it is always quoted with a direction and an angle (e.g.\ $30^{\circ}$ towards the SE).
\item \cle{Outcrop}: the area where a rock unit appears at the surface; its width depends on the dip of the beds and on the topography.
\item \cle{Strike line} (structure contour): a line joining points of equal height on the same geological surface (bed boundary, vein, fault plane).
\item \cle{Unconformity}: a buried erosion surface separating two series of strata, recording a gap in deposition.
\item \cle{Vertical exaggeration} (V.E.): the ratio of the vertical scale to the horizontal scale on a cross-section; it must always be stated on the section.
\end{itemize}

\textbf{Formulas and relations to know}
\begin{itemize}
\item Scale: $\text{scale} = \dfrac{\text{map distance}}{\text{ground distance}}$ (both in the same unit). At $1/25\,000$, $1\ \mathrm{cm}$ on the map represents $25\,000\ \mathrm{cm} = 250\ \mathrm{m}$ on the ground.
\item True dip from strike lines: $\tan\delta = \dfrac{\text{vertical interval between two strike lines}}{\text{horizontal distance between them measured perpendicular to strike}}$, i.e.\ $\delta = \arctan\left(\dfrac{h}{d}\right)$.
\item Apparent dip in any other direction: $\tan\alpha = \tan\delta \times \sin\theta$, where $\theta$ is the angle between the section line and the true dip direction. Note that $\alpha \leq \delta$ always, and $\alpha = \delta$ only when the section is perpendicular to strike ($\theta = 90^{\circ}$).
\item True thickness of a dipping bed on flat ground: $t = w\sin\delta$, where $w$ is the outcrop width measured perpendicular to strike. For a vertical bed $t = w$; for a horizontal bed the thickness is simply the difference in height between its two boundaries read from the contours.
\end{itemize}

\textbf{Methods}
\begin{itemize}
\item \emph{Three-point problem}: (1) join the highest and the lowest of the three points; (2) by proportional interpolation along that line, locate the point having the intermediate height; (3) join this point to the intermediate point: this is a strike line; (4) the dip direction is perpendicular to the strike line, down the height gradient; (5) draw a second strike line through one of the remaining points and compute $\tan\delta = h/d$.
\item \emph{Cross-section}: (1) draw the section line on the map; (2) transfer every geological boundary intersection and every contour crossing onto the edge of a strip of paper; (3) plot the topographic profile at the chosen vertical scale; (4) project the observed dips (converted to apparent dips if the section is oblique to strike); (5) draw the boundaries at depth, respecting folds, faults, unconformities and intrusions; (6) add ornament, title, orientation, horizontal and vertical scales, V.E.\ and key.
\item \emph{Geological history}: apply the principles of superposition, cross-cutting relationships, included fragments and unconformities to order all events from oldest to youngest, justifying each step.
\end{itemize}

\textbf{Common mistakes}
\begin{itemize}
\item Confusing dip direction with strike: the two are perpendicular, never parallel.
\item Measuring the dip along a line that is not perpendicular to strike: this gives an \emph{apparent} dip, always smaller than the true dip.
\item Forgetting that V-shaped outcrop deflections in valleys depend on the dip of the bed relative to the valley gradient (rule of V's): a horizontal bed's V points upstream; a vertical bed shows no V.
\item Mixing up relative ages across faults and intrusions: the cross-cutting body (fault, dyke, batholith) is always \emph{younger} than the rocks it cuts.
\item Starting an interpretation without first reading the legend, the scale and the north arrow.
\item Forgetting to state the vertical exaggeration, or drawing dipping boundaries as vertical lines on an exaggerated section without recomputing the apparent dip.
\end{itemize}

\textbf{GCE tips}
\begin{itemize}
\item Always state the strike as a bearing (e.g.\ N$040^{\circ}$E) and the dip with both direction and angle (e.g.\ $30^{\circ}$ SE).
\item Label every cross-section fully: title, orientation of the section line, horizontal and vertical scales, V.E., and a key to the ornaments used.
\item For photographs: first locate and orient the view, then describe the landforms, identify the rocks and structures, and only then interpret the processes responsible.
\item Present the geological history as a numbered chronological sequence (oldest $\rightarrow$ youngest), with one justification per event.
\item Show all working in dip and thickness calculations: the method marks are often worth more than the final number.
\end{itemize}
\end{bilan}

\findecours

\end{document}
